\documentclass[11pt,reqno]{amsart}
\usepackage[localtheorems]{raeez-math-template}

\usepackage{array}

% ---------- Theorem environments ----------
\theoremstyle{plain}
\newtheorem{theorem}{Theorem}[section]
\newtheorem{proposition}[theorem]{Proposition}
\newtheorem{lemma}[theorem]{Lemma}
\newtheorem{corollary}[theorem]{Corollary}
\newtheorem{conjecture}[theorem]{Conjecture}

\theoremstyle{definition}
\newtheorem{definition}[theorem]{Definition}
\newtheorem{construction}[theorem]{Construction}
\newtheorem{example}[theorem]{Example}

\theoremstyle{remark}
\newtheorem{remark}[theorem]{Remark}
\newtheorem{question}[theorem]{Question}

% ---------- Claim-status tags ----------
\providecommand{\ClaimStatusProvedHere}{\textnormal{[proved here]}}
\providecommand{\ClaimStatusProvedElsewhere}{\textnormal{[proved elsewhere]}}
\providecommand{\ClaimStatusConjectured}{\textnormal{[conjectured]}}
\providecommand{\ClaimStatusHeuristic}{\textnormal{[physical heuristic]}}
\providecommand{\ClaimStatusConditional}{\textnormal{[conditional]}}
\providecommand{\ClaimStatusOpen}{\textnormal{[open]}}

% ---------- Operators ----------
\DeclareMathOperator{\Ext}{Ext}
\DeclareMathOperator{\Hom}{Hom}
\DeclareMathOperator{\End}{End}
\DeclareMathOperator{\Tor}{Tor}
\DeclareMathOperator{\tr}{tr}
\DeclareMathOperator{\Coh}{Coh}
\DeclareMathOperator{\Aut}{Aut}
\DeclareMathOperator{\id}{id}
\DeclareMathOperator{\Sp}{Sp}
\DeclareMathOperator{\Conf}{Conf}
\DeclareMathOperator{\HH}{HH}
\DeclareMathOperator{\HC}{HC}
\DeclareMathOperator{\Mod}{Mod}
\DeclareMathOperator{\Spec}{Spec}
\DeclareMathOperator{\Fact}{Fact}
\DeclareMathOperator{\av}{av}

% ---------- Notation ----------
\newcommand{\CY}{\mathrm{CY}}
\newcommand{\AdS}{\mathrm{AdS}}
\newcommand{\Cbb}{\mathbb{C}}
\newcommand{\Rbb}{\mathbb{R}}
\newcommand{\Zbb}{\mathbb{Z}}
\newcommand{\Pbb}{\mathbb{P}}
\newcommand{\Qbb}{\mathbb{Q}}
\newcommand{\Hbb}{\mathbb{H}}
\newcommand{\cM}{\mathcal{M}}
\newcommand{\cU}{\mathcal{U}}
\newcommand{\cH}{\mathcal{H}}
\newcommand{\cP}{\mathcal{P}}
\newcommand{\cO}{\mathcal{O}}
\newcommand{\cA}{\mathcal{A}}
\newcommand{\cE}{\mathcal{E}}
\newcommand{\cC}{\mathcal{C}}
\newcommand{\cT}{\mathcal{T}}
\newcommand{\sfG}{\mathsf{G}}
\newcommand{\sfL}{\mathsf{L}}
\newcommand{\sfC}{\mathsf{C}}
\newcommand{\sfM}{\mathsf{M}}
\newcommand{\sfB}{\mathsf{B}}
\newcommand{\bfH}{\mathbf{H}}
\newcommand{\at}{\mathrm{at}}
\newcommand{\Bar}{\mathrm{B}}
\newcommand{\Cobar}{\Omega}
\newcommand{\hbar}{{\hslash}}

\title[Atiyah--Connes obstruction]{The Atiyah--Connes obstruction to chiral bar--cobar specialisation}

\author{Raeez Lorgat}
\address{}
\date{\today}

\begin{document}

\begin{abstract}
Let $X$ be a smooth projective Calabi--Yau threefold with holomorphic
volume form $\Omega_X$ and fixed cyclic dg enhancement
$\mathrm{Perf}_{\mathrm{dg}}(X)$ of $D^b(\Coh X)$. The
Shklyarov--Brav--Dyckerhoff Calabi--Yau trace transports Connes'
chain operator $B$ to a cochain $S^1$-operator $\Delta_2$ on
Hochschild cochains; the ternary Stasheff operation $m_3$ acts by
Deligne's conjecture on the same complex. The commutator
$\gamma_X=[m_3,\Delta_2]$ is $\delta_{\mathrm{Hoch}}$-closed of total
Hochschild degree $4$; its cohomology class projects under the
Hochschild--Kostant--Rosenberg isomorphism to
$\beta_X\in H^2(X,\wedge^2T_X)$, and the Calabi--Yau contraction
produces the \emph{Atiyah--Connes class}
\[
  \alpha_X
  :=
  \iota_{\Omega_X}\bigl(\pi^{\mathrm{HKR}}_{2,2}[\gamma_X]\bigr)
  \in H^2(X,\Omega^1_X).
\]
The construction is independent of the Dolbeault representative of the
Atiyah class and invariant under cyclic $A_\infty$-quasi-equivalence
preserving the cyclic-HKR datum.

On a smooth curve $C$, $H^2(C,\Omega^1_C)=0$, so ordinary coherent
restriction cannot detect $\alpha_X$: any bridge to the curve-side
theory must pass through a \emph{stage-two specialisation}
$\Sp^{\mathrm{ch}}_{\Sigma,C}=\operatorname{res}_C\circ\int_\Sigma$
(factorisation homology followed by restriction). Let
$A_C=\Sp^{\mathrm{ch}}_{\Sigma,C}\Phi^{\mathrm{FA}}_3
(\mathrm{Perf}_{\mathrm{dg}}(X))$ and let
$\Defcyc^{\mathrm{mod}}(A_C)$ denote the completed modular cyclic
deformation complex with genus-zero arity-three projection
$\pi_{3,0}$.  The paper's main theorem is that stage-two
specialisation carries $\alpha_X$ onto the first bar--cobar
obstruction:
\[
  \operatorname{sp}_{\Sigma,C}(\alpha_X)
  =
  o^{(0)}_3(A_C)
  \in H^2\!\left(
    \pi_{3,0}\Defcyc^{\mathrm{mod}}(A_C)
  \right).
\]
Consequently $\alpha_X=0$ kills the first genus-zero chiral
bar--cobar obstruction on $C$.

Further consequences are conditional. Full bar--cobar inversion
requires an independent Koszul/PBW condition, completed Positselski
hypotheses, or a higher-obstruction closure theorem. The comparison
between the zero locus of $\alpha_X$ and Costello--Li factorisation
canonicality, the Vol.~III principal $K3\times E$
$\mathbf H_{\Delta_5}$ source-recognition target, and the
derived-centre Mukai conductor trace
$\hbar^2K^{\kappa_{\mathrm{ch}}}_{\mathrm{Muk}}=-1$ is recorded as a
comparison diagram of derived zero loci under explicit hypotheses and
remains conjectural in full generality. Three geometries exercise the
theorem: the generic quintic has $101$-dimensional target and non-zero
Yukawa detector; $K3\times E$ reduces its obstruction to a single
cyclic scalar $\lambda_{\mathrm{Cyc}}(S,E,\mathfrak c)$; and Clemens'
non-K\"ahler small-resolution threefolds require a heterotic-with-torsion
Courant/CDO replacement of the Dolbeault construction.
\end{abstract}

\maketitle

\setcounter{tocdepth}{2}

% ==================================================
\section{Introduction}
% ==================================================

\emph{Unifying principle.} A smooth projective Calabi--Yau threefold
$X$ carries a single cyclic Hochschild obstruction class
\[
  \alpha_X=\iota_{\Omega_X}\pi^{\mathrm{HKR}}_{2,2}[m_3,\Delta_2]
  \in H^2(X,\Omega^1_X)=H^{1,2}(X)
\]
built from the ternary Stasheff operation $m_3$ on Hochschild cochains
and the cochain Connes rotation $\Delta_2$, the transport of Connes'
chain operator $B=(1-\lambda)sN$ under the Shklyarov--Brav--Dyckerhoff
cyclic Calabi--Yau trace. The construction is carried out in
\S\ref{sec:cocycle}; representative-independence, Hodge type, and
invariance under cyclic $A_\infty$-quasi-equivalence are proved in
Propositions~\ref{prop:cocycle-rep},~\ref{prop:atiyah-cup-derivation}
and Corollary~\ref{cor:hodge}. The Hodge-theoretic notation
$H^{2,1}(X)$ used elsewhere in the programme is the Calabi--Yau-dual
convention under $\Omega_X$, not the sheaf-cohomology target.

The curve-side target is not $H^2(C,\Omega^1_C)$, which vanishes for
every smooth curve, but the completed modular cyclic deformation
complex $\Defcyc^{\mathrm{mod}}(A_C)$ of the curve-side chiral algebra
\[
  A_C=\Sp^{\mathrm{ch}}_{\Sigma,C}
  \Phi^{\mathrm{FA}}_3(\mathrm{Perf}_{\mathrm{dg}}(X)),
\]
where $\Phi^{\mathrm{FA}}_3$ is Vol.~III's stage-one holomorphic
factorisation-algebra functor and $\Sp^{\mathrm{ch}}_{\Sigma,C}$ is
stage-two specialisation (factorisation homology over a two-cycle
$\Sigma\subset X$ followed by restriction to the reference curve
$C$). The Ayala--Francis derivative
\[
  d\Sp^{\mathrm{ch}}_{\Sigma,C}:
  \Def^{E_3,\mathrm{cyc}}(\Phi^{\mathrm{FA}}_3(X))
  \to
  \Defcyc^{\mathrm{mod}}(A_C)
\]
is a filtered $L_\infty$-morphism of formal moduli problems
(Theorem~\ref{thm:sp-differentiable}); its arity-three genus-zero
component identifies $\alpha_X$ with the first bar--cobar obstruction
(Theorem~\ref{thm:sp-alpha-first-barcobar-obstruction}):
\[
  \operatorname{sp}_{\Sigma,C}(\alpha_X)
  :=
  \pi_{3,0}d\Sp^{\mathrm{ch}}_{\Sigma,C}(\alpha_X)
  =
  o^{(0)}_3(A_C)
  \in H^2\!\left(\pi_{3,0}\Defcyc^{\mathrm{mod}}(A_C)\right).
\]
Vanishing of $\alpha_X$ kills the first genus-zero chiral bar--cobar
obstruction of $A_C$ on the reference curve.

The comparison with Theorem~B of Volume~I (strict or completed
Positselski bar--cobar inversion), with the holomorphic
factorisation-algebra canonicality on the Costello--Li locus, with the
Vol.~III principal $K3\times E$ $\bfH_{\Delta_5}$
source-recognition target, and with the Bruinier--Mukai conductor trace identity
$\hbar^2K^{\kappa_{\mathrm{ch}}}_{\mathrm{Muk}}=-1$ is stated as a
diagram of derived zero loci
(Theorem~\ref{thm:atiyah-connes-bridge}) under explicit
comparison hypotheses $(B1)$--$(B4)$. The five-vertex pentagon is the
conjectural simultaneous closure of these bridges
(Conjecture~\ref{conj:atiyah-connes-pentagon}).

Three geometries exercise the theorem. The generic quintic
$Q_5\subset\Pbb^4$ has $h^{1,2}(Q_5)=h^{2,1}(Q_5)=101$ and non-zero
Yukawa coupling, with the Atiyah--Connes residue
$\Psi_f(\alpha_{Q_5})\in R_f^{10}$ an open Griffiths computation. The
product $K3\times E$ has
$\at_{K3\times E}=\pi_{K3}^*\at_{K3}\neq 0$ and
$H^2(K3\times E,\Omega^1)$ of dimension $20+1$; product naturality of
the cyclic-HKR datum kills the $20$-dimensional mixed summand and
reduces vanishing of $\alpha_{K3\times E}$ to a single scalar
$\lambda_{\mathrm{Cyc}}(S,E,\mathfrak c)$ on the one-dimensional line
$H^2(S,\cO_S)\otimes H^0(E,\Omega^1_E)$. Clemens' non-K\"ahler
small-resolution threefolds violate the $\partial\bar\partial$-lemma,
so the K\"ahler Dolbeault representative fails and the construction
must be rebuilt in an $H$-twisted Courant-algebroid CDO setting with
closed three-form Strominger flux $H$.

% ==================================================
\section{Setup and notation}
\label{sec:setup}
% ==================================================

\subsection{Calabi--Yau categories and Serre functor}
\label{subsec:cy-serre}
Fix a smooth projective Calabi--Yau threefold $X$ over $\Cbb$: a smooth projective variety of dimension $3$ with trivial canonical bundle $\omega_X \cong \cO_X$ and the Hodge hypothesis $h^{1,0}(X) = h^{2,0}(X) = 0$ (so that $h^{3,0}(X) = 1$ by Serre duality). Simple connectedness of $X$ is a sufficient condition; $K3 \times E$ fails ($\pi_1 = \Zbb^2$, $h^{1,0} = 1$) and enters the argument after a $\Zbb/n$ Borcea--Voisin orbifold twist absorbing the $b_1(E) = 2$ discrepancy (§\ref{sec:clemens}). Bondal--Kapranov 1989 constructs the bounded derived category $D^b(X) = D^b(\Coh X)$ and Kuznetsov 2004 and Keller 2003 characterise it as a Calabi--Yau-$d$ category in the sense that the Serre functor
\[
  \mathbb{S}_{D^b(X)}(-) \;=\; (- \otimes \omega_X)[d] \;\simeq\; [d]
\]
is naturally isomorphic to the $d$-fold shift, with $d = \dim X = 3$. The Hochschild cohomology $\HH^\bullet(X) = \HH^\bullet(D^b(X))$ decomposes by Hochschild--Kostant--Rosenberg (Caldararu 2005) as
\[
  \HH^n(X) \;\cong\; \bigoplus_{p+q = n} H^p(X, \wedge^q T_X),
\]
and the two-term filtration $\HH^{\leq 2}$ encodes first-order noncommutative deformations (Gerstenhaber 1964, Kontsevich 2003).

\medskip
\noindent\textbf{Dg-enhancement and $\mathrm{Pr}^{\mathrm{L}}_{k,\mathrm{st}}$-framing.} The Serre equivalence $\mathbb{S} \simeq [3]$ is read in $\mathrm{Pr}^{\mathrm{L}}_{k,\mathrm{st}}$, the $(\infty,1)$-category of $k$-linear presentable stable $\infty$-categories with colimit-preserving functors (Lurie $\mathrm{HA}$~\S1.1.1, \S4.8.1, \S5.5.3; To\"en--Vezzosi 2007 \S2; Cohn 2016 \S3). The enhancement $\cD^{b}_{\mathrm{dg}}(X)$ lies in this $(\infty,1)$-category; its uniqueness up to quasi-equivalence is Lunts--Orlov 2010 Theorem~2.8. The reconstruction theorem of Bondal--Orlov 2001 (Theorem~3.1) forces $D^b(Y) \simeq D^b(X)$ to imply $Y \simeq X$ whenever $\omega_X$ is ample or anti-ample; for $X$ Calabi--Yau the triviality $\omega_X \simeq \cO_X$ sits at the degenerate boundary where the reconstruction fails and the Calabi--Yau-$d$ autoequivalence $\mathbb{S} \simeq [d]$ becomes the intrinsic invariant of $\cD^{b}_{\mathrm{dg}}(X)$ as an object of $\mathrm{Pr}^{\mathrm{L}}_{k,\mathrm{st}}$. Compact objects coincide with perfect complexes, $\mathrm{Perf}(X) \subset \cD^{b}_{\mathrm{dg}}(X)$ (Bondal--Van den Bergh 2003 Theorem~3.1.1), and the Hochschild cohomology $\HH^\bullet(X)$ is a derived invariant of the underlying object of $\mathrm{Pr}^{\mathrm{L}}_{k,\mathrm{st}}$ (Keller 2003 Theorem~A).

\subsection{The Atiyah class: local versus global, computed in $D^b(\Coh X)$}
\label{subsec:atiyah-local-global}
Given a locally free sheaf $\cE$ on $X$, its Atiyah class $\at_\cE \in \Ext^1_X(\cE, \cE \otimes \Omega^1_X)$ is the obstruction to splitting the jet exact sequence
\[
  0 \;\longrightarrow\; \cE \otimes \Omega^1_X \;\longrightarrow\; J^1 \cE \;\longrightarrow\; \cE \;\longrightarrow\; 0.
\]
The Ext group is computed in the bounded derived category of coherent sheaves $D^b(\Coh X)$ (Bondal--Kapranov 1989, Keller 2003, Lunts--Orlov 2010); for $\cE$ perfect the class admits the equivalent description as the obstruction to lifting $\cE \in D^b(\Coh X)$ to a perfect complex over the first-order thickening $X \times \Spec k[\epsilon]/\epsilon^2$, via Illusie 1971--72 Chapter IV \S2.3. Markarian 2009 Theorem 1.1 and Ramadoss 2008 Theorem 3.3.1 identify $\at_\cE$ as the image of $\id_\cE$ under the Chern-character tangent map in the $(\infty,1)$-stable category $D^b(\Coh X)$; the $(\infty,1)$-categorical formulation of Toen 2014 Theorem 1.4 realises $\at$ as a natural transformation $\id_{D^b(X)} \Rightarrow (-) \otimes \Omega^1_X[1]$ of $\infty$-functors.

The local-global dichotomy is load-bearing. Fix an affine cover $\cU_X = \{U_\alpha\}$ of $X$ trivialising $\cE$. On each chart $U_\alpha$ the jet sequence splits; a choice of smooth connection $\nabla_\alpha$ on $\cE|_{U_\alpha}$ witnesses a local splitting. The Čech $1$-cocycle of connection differences
\[
  \at_\cE^{\check{\mathrm C}} \;:=\; \{\nabla_\alpha - \nabla_\beta\}_{\alpha\beta} \;\in\; C^1(\cU_X, \mathcal{H}om(\cE, \cE \otimes \Omega^1_X))
\]
represents the global class $\at_\cE \in H^1(X, \mathcal{H}om(\cE, \cE \otimes \Omega^1_X)) \cong \Ext^1_X(\cE, \cE \otimes \Omega^1_X)$ via the Čech-to-Ext spectral sequence; passing to $D^b(\Coh X)$ the isomorphism is the Čech-Ext comparison of Hartshorne \emph{Residues and Duality} III \S7. The local datum is a connection on a chart; the global invariant in $D^b(\Coh X)$ is the cohomology class of its patching discrepancy.

Four statements of this paper live strictly at global scope. The Hodge decomposition of Corollary~\ref{cor:hodge} is a global Dolbeault statement. The Yukawa non-vanishing $\int_{Q_5} \at_{Q_5}^{\cup 3} \neq 0$ of Proposition~\ref{prop:quintic-nonvanish} is a global integral pairing. The Künneth additivity $\at_{K3 \times E} = \at_{K3} \boxplus \at_E$ of Proposition~\ref{prop:k3e-vanish} is a decomposition of a global class under a global product structure. The vanishing $[m_3, B^{(2)}]_X = 0$ lives in $H^2(X, \Omega^1_X)$. None admits a chart-local avatar: on any chart $U_\alpha$ the restriction $\at_\cE|_{U_\alpha}$ vanishes trivially, witnessed by $\nabla_\alpha$. Conversely, the chain-level witnesses of §\ref{sec:cocycle} and Theorem~\ref{thm:clemens-extension} --- Massey primitives $\alpha$ with $d\alpha = f \circ g$, Kotake--Kato chart-local heat kernels, chart-local smooth connections $\nabla_\alpha$ and their $\bar\partial$-derivatives --- are by construction chart-local; the content of each vanishing theorem is the obstruction to patching these chart-local primitives into one global primitive. Local vanishing is automatic. Global vanishing is the theorem.

Kapranov 1999 (Prop.~4.4) establishes that for $\cE = T_X$ the tangent bundle, the Atiyah class $\at_X := \at_{T_X} \in \Ext^1_{D^b(\Coh X)}(T_X, T_X \otimes \Omega^1_X)$ generates, by iterated self-composition in $D^b(\Coh X)$, a graded Lie algebra structure on the shift $T_X[-1]$ in the $(\infty,1)$-stable category $D^b(\Coh X)$, with brackets
\[
  [\at_X, \at_X, \ldots, \at_X] \;\in\; \Ext^{n-1}_{D^b(\Coh X)}(T_X, T_X^{\otimes (n-1)} \otimes \Omega^1_X),
\]
whose Jacobiator is controlled by the next term in the Bianchi identity of the underlying Kähler connection. For $X$ Calabi--Yau of dimension $3$, $\at_X$ lifts to
\[
  \at_X \;\in\; H^1(X, \End(T_X) \otimes \Omega^1_X) \;\subset\; H^1(X, T_X \otimes \Omega^1_X \otimes \Omega^1_X),
\]
with the enclosing $\End$ factor tracking the $E_1$-associative structure on $T_X[-1]$ within $D^b(\Coh X)$. The Čech representative is $\at_X^{\check{\mathrm C}} \in C^1(\cU_X, \End(T_X) \otimes \Omega^1_X)$; the Dolbeault representative used in §\ref{sec:cocycle} is $\bar\partial \nabla \in \Omega^{0,1}(X, \End(T_X) \otimes \Omega^1_X)$ for a global Chern connection $\nabla$ of a Hermitian metric. The two representatives agree under the Dolbeault--Čech comparison in the analytic topos: the local connection and the global $\bar\partial$-failure are two faces of one obstruction class in the $(\infty,1)$-stable category $D^b(\Coh X)$.

\subsection{Connes' cyclic operator and its $(\infty,1)$-operadic home}
\label{sub:B2-operad}
For a unital associative algebra $A$, the cyclic bar complex $\mathrm{CC}^{\mathrm{per}}_\bullet(A)$ carries Connes' cyclic operator $B : \HH_n(A) \to \HH_{n+1}(A)$ (Connes 1985), constructed as $B = (\id - \lambda)s N$, where $s$ is the extra degeneracy, $\lambda$ the cyclic rotation, and $N = 1 + \lambda + \lambda^2 + \cdots + \lambda^n$. Three operadic homes are distinct and load-bearing.

The first is Connes' cyclic category $\Lambda$ (Connes 1983): $\mathrm{CC}^{\mathrm{per}}_\bullet(A)$ is a $\Lambda$-object, and $B$ is the image under the nerve $N_\Lambda$ of the generator of $\pi_1(B\Lambda) = \pi_1(BS^1) = \Zbb$.

The second is the cyclic operad $\mathrm{Cyc}$ (Getzler--Kapranov 1995, Loday 1998 \S6.1, Markl--Shnider--Stasheff 2002 \S5.4), a coloured operad enriched over pointed simplicial sets whose algebras are associative algebras equipped with an invariant symmetric pairing; the Connes complex is the natural home for its cohomology. $\mathrm{Cyc}$ is \emph{not} the same as $E_1$: $\mathrm{Cyc}$-algebras are associative-with-trace, not just associative.

The third is the framed little-$2$-discs operad $fE_2 \simeq E_2 \rtimes SO(2)$ of Getzler 1994 and Getzler--Jones 1994 (\S1 Theorem 1, \S2 Theorem 2.1): Hochschild cochains $C^\bullet(A, A)$ of an associative algebra carry a canonical $fE_2$-algebra structure whose $E_2$-part is the Gerstenhaber bracket--cup-product pair and whose $SO(2)$-part is generated by Connes' $B$. $fE_2$ lives in $\mathrm{Op}^\infty_2$, strictly above $E_1$; the $E_1$-structure on $C^\bullet(A, A)$ is the image of the inclusion $E_1 \hookrightarrow fE_2$. The Connes operator $B$ is an $fE_2$-operation --- specifically the $SO(2)$-action --- not an $E_1$-operadic operation.

In factorisation-homology form: Lurie (\emph{Higher Algebra}, Theorem 5.5.3.8) and Ayala--Francis 2015 Theorem 5.5 establish that for an $E_1$-algebra $A$ in a stable presentable $(\infty,1)$-category, cyclic homology \emph{equals} factorisation homology over the circle, $\HC_\bullet(A) = \int_{S^1} A$. The Connes operator $B$ is the image of the fundamental class $[S^1] \in H_1(S^1; \Zbb)$ under the canonical $S^1$-action on $\int_{S^1} A$ obtained by rotation; this $S^1$-action is the $SO(2)$-part of the $fE_2$-action on Hochschild cochains, restricted along the loop $S^1 \hookrightarrow SO(2)$.

At bidegree $2$, $B^{(2)}: \HH_2(A) \to \HH_3(A)$ \emph{is} the obstruction to lifting a Hochschild $2$-cycle to a cyclic $3$-cycle and \emph{is} the degree-two component of the $S^1$-rotation. Dually, on Hochschild cochains $\HH^\bullet(A) = \pi_\bullet \End_A(A)$, under HKR (Caldararu 2005 Proposition 4.4),
\[
  B^{(2)} \;=\; (-) \lrcorner \Omega^1_X \qquad\text{on polyvector fields},
\]
and at the $(\infty,1)$-level,
\[
  B^{(2)} \;=\; [1]_{S^1} \qquad\text{the degree shift induced by the $S^1$-action on } T_{D^b(X)} \text{ (Toen--Vezzosi 2015 Theorem 2.6.1)}.
\] The pairing with $m_3$ below is computed in the $fE_2$-operadic structure on $C^\bullet(\End^\bullet(\cE), \End^\bullet(\cE))$, with output living in the negative cyclic bicomplex $\HC^-_\bullet(A)$.

\subsection{The ternary $A_\infty$-bracket and its $(\infty,1)$-operad}
\label{sub:m3-operad}
Beilinson 1978 describes $\End^\bullet(\cE)$, for $\cE$ a complex of locally free sheaves, as an algebra over the topological $A_\infty$-operad $\cA_\infty = \{K_n\}_{n \geq 2}$, where $K_n$ is the Stasheff associahedron of dimension $n - 2$ (Stasheff 1963). The operadic suspension $C_\bullet(\cA_\infty)$ is a cofibrant resolution of the associative operad $\mathrm{Ass}$ in the model category $\mathrm{Op}(\mathrm{Ch}_k)$ of dg-operads over the ground field $k$, and Boardman--Vogt 1973 \S4 identifies $\cA_\infty$ with the little-$1$-cubes operad $E_1$. Passing to the $(\infty,1)$-category $\mathrm{Op}^\infty_1$ of $\infty$-operads via Lurie's operadic nerve (Lurie, \emph{Higher Algebra}, Definition 2.1.1.10), the two operads represent the same object: the unique $E_1 = \cA_\infty \in \mathrm{Op}^\infty_1$ whose algebras in the $(\infty,1)$-category $\mathrm{Mod}_k^\infty$ of chain complexes are homotopy-coherent associative algebras (Lurie, \emph{Higher Algebra}, Theorem 4.1.8.4). The operadic depth is $n = 1$: every $m_n$ of this paper lives in $E_1 = \cA_\infty$, not in $E_k$ for $k \geq 2$.

The binary product $m_2$ is composition; the higher brackets
\[
  m_n \;:\; \End^\bullet(\cE)^{\otimes n} \;\longrightarrow\; \End^\bullet(\cE)[2 - n]
\]
are the Stasheff operations indexed by the cellular chain $[K_n] \in C_{n-2}(\cA_\infty(n))$. The ternary bracket $m_3$ is the pentagonal cell $[K_3] \in C_1(\cA_\infty(3))$, classifying the Massey triple-product obstruction, non-zero precisely when the associator
\[
  (f \circ g) \circ h - f \circ (g \circ h)
\]
fails to vanish on the nose. Kontsevich 2003 Theorem 4.6 supplies the formality $L_\infty$-quasi-isomorphism $T^\bullet_{\mathrm{poly}}(X) \simeq D^\bullet_{\mathrm{poly}}(X)$ --- an $E_1$-operadic statement in $\mathrm{Op}^\infty_1$ at the level of $L_\infty = \mathrm{Lie}_\infty$-operadic algebras; the corresponding associative-operadic statement is the $\cA_\infty$-formality of Kontsevich--Soibelman 2000 Theorem 2.4, identifying $C^\bullet(\cO_X, \cO_X)$ as an $\cA_\infty$-algebra in $\mathrm{Mod}_k^\infty$. Under this formality, the $m_3$ of this paper is the image of $[K_3]$ on the $\End^\bullet(\cE)$-side; dually, it is the order-three Taylor coefficient of the Maurer--Cartan cochain solving the $(\infty,1)$-operadic lift of $\End^\bullet(\cE)$ to an $E_1$-algebra in $\mathrm{Mod}_k^\infty$. Caldararu--Willerton 2010 Theorem 1.6 and Calaque--Van den Bergh 2010 \emph{equate} the Yoneda product on $\Ext^\bullet(T_X, T_X)$ with the Kapranov $L_\infty$-bracket on the tangent complex $T_{\cM_{\mathrm{cx}}(X)}$. Operadic depth of $m_3$ is $1$; appeals to $E_k$ for $k \geq 2$ refer not to $m_3$ itself but to the $fE_2$-Gerstenhaber--BV structure of \S\ref{sub:B2-operad} in which the commutator $[m_3, B^{(2)}]$ is computed.

\subsection{The common ambient for $m_3$ and $B^{(2)}$}
\label{sub:common-ambient}
Both operators act on the Hochschild cochain complex $\mathrm{CC}^\bullet(\End \cE)$; this is the scope on which the commutator $[m_3, B^{(2)}]$ is defined. Getzler--Jones 1994 (\S1 Theorem~1, \S2 Theorem~2.1) place both inside the same $fE_2$-operadic structure: $m_3$ is the pentagonal cell $[K_3]$ inside the $E_1 \hookrightarrow fE_2$ associative part, and $B^{(2)}$ is the bidegree-two piece of the $SO(2)$-action. At cochain level, $m_3 : \mathrm{CC}^\bullet \to \mathrm{CC}^{\bullet - 1}$ is the degree-$(-1)$ endomorphism pre-composing with the triple Stasheff bracket on three consecutive bar slots (shift $-1 = 2 - 3$ is the Stasheff degree of $[K_3]$), and $B^{(2)} : \mathrm{CC}^\bullet \to \mathrm{CC}^{\bullet + 1}$ is the degree-$(+1)$ cyclic shift (shift $+1$ is the cochain degree of $[S^1]$ under the $S^1$-action). The graded commutator $[m_3, B^{(2)}] = m_3 \circ B^{(2)} + B^{(2)} \circ m_3$ has total degree zero, commutes with the Hochschild differential by the $fE_2$-Cartan relations (Getzler 1994 Prop.~1.2, Getzler--Jones 1994 \S2), and descends to a class in Hochschild cohomology. HKR (Caldararu 2005 Prop.~4.4) translates the ambient into polyvector fields: $m_3$ becomes wedge-multiplication by $\at_X$ (Kapranov 1999 Prop.~4.4), $B^{(2)}$ becomes contraction against the CY volume form $\Omega_X \in H^0(X, \Omega^3_X)$, and the commutator lands in $H^2(X, \Omega^1_X)$ as announced in \S\ref{sec:cocycle}.

\subsection{Chain-level and $(\infty,1)$-categorical lanes}
Every construction here admits two models: a chain-level model --- explicit complexes, named differentials, witnessed chain homotopies --- and an $(\infty,1)$-categorical model --- Lurie's $\mathrm{HA}$ formalism for derived categories, Francis--Gaitsgory for factorisation. Chain-level statements compute; $(\infty,1)$-categorical statements express universal properties. Each proposition carries the tag of the lane in which its proof is given. The obstruction class is a cyclic-HKR construction in both lanes once the datum of Definition~\ref{def:cyclic-hkr-datum} is fixed; equivalence with the four comparison targets requires the hypotheses stated in Theorem~\ref{thm:main}.

\subsection{The reference curve and the two-stage factorisation as derived functors}
\label{subsec:two-stage-derived}
Fix a smooth algebraic curve $C$ over $\Cbb$ and an immersed real $2$-cycle $\Sigma_2 \hookrightarrow X$ whose tubular neighbourhood retracts onto $C$. The two-stage Calabi--Yau-to-chiral factorisation of Volume~III reads
\[
  \Phi_3 \;=\; \Sp^{\mathrm{ch}}_{\Sigma_2, C} \circ \Phi^{\mathrm{FA}}_3,
\]
with stage one producing the canonical holomorphic $E_3$-factorisation algebra (Costello--Gwilliam--Li 2021 locality, Kontsevich--Tamarkin 2000 $E_d$-formality; the Weiss-topology prefactorisation-algebra data and the compactness hypothesis under which this is a factorisation algebra in the Costello--Gwilliam sense are inscribed in \S\ref{sub:fa-precision} below) and stage two giving its factorisation homology over $\Sigma_2$ followed by restriction to $C$. The reference curve $C$ is the geometric locus on which Theorem~B of Volume~I lives: Theorem~A's bar--cobar equivalence \emph{equals} the $E_1$-chiral specialisation of $\Phi_3(X)$ to $C$.

\subsection{Factorisation-algebra precision: Weiss topology, structure maps, locality}
\label{sub:fa-precision}
Stage one $\Phi^{\mathrm{FA}}_3$ produces a prefactorisation algebra on the underlying real $6$-manifold $X^{\mathrm{top}}$, locally constant in the open-set argument and holomorphic in the $\bar\partial$-dependence; the data under which it is a factorisation algebra in the sense of Costello--Gwilliam (CG FA Vol~1 Definition~6.1.4, Vol~2 Definition~3.1.1) are as follows. Let $\mathrm{Disj}(X)$ denote the coloured operad whose colours are open subsets $U \subset X^{\mathrm{top}}$ and whose $(U_1, \ldots, U_n; V)$-operations exist precisely when the $U_i$ are pairwise disjoint and $\bigsqcup_i U_i \subset V$. A Weiss cover $\mathfrak{U}^{\mathrm{Wss}}$ of $V \subset X^{\mathrm{top}}$ is a cover by opens of the form $U_1 \sqcup \cdots \sqcup U_n$ with each $U_i \subset V$ open and the $U_i$ pairwise disjoint; the Weiss topology is the Grothendieck topology on $\mathrm{Open}(X^{\mathrm{top}})$ generated by Weiss covers. A prefactorisation algebra valued in cochain complexes is a symmetric monoidal functor $\cF : \mathrm{Disj}(X) \to \Mod_k$; it is a factorisation algebra when it satisfies Weiss-descent: for every Weiss cover of $V$,
\[
  \mathrm{colim}_{\mathfrak{U}^{\mathrm{Wss}}} \cF(U_{\alpha_1} \sqcup \cdots \sqcup U_{\alpha_n}) \;\xrightarrow{\sim}\; \cF(V)
\]
is a quasi-isomorphism. For disjoint $U_1, \ldots, U_n \subset V$ the factorisation structure map $\mu^{\cF}_{U_1, \ldots, U_n; V} : \cF(U_1) \otimes \cdots \otimes \cF(U_n) \to \cF(V)$ records the inclusion of disjoint neighbourhoods, with symmetry, associativity, and unit axioms those of a symmetric monoidal functor on the Weiss operad.

On a contractible real chart $U \cong \Rbb^6$, locally constant
factorisation algebras on $U$ are equivalent to $E_6$-algebras in
$\Mod_k$ (Lurie HA Theorem~5.4.5.9; Francis--Gaitsgory identify the
real case with $E_n$ for $n$ the real dimension). The $E_3$ notation in
this paper refers instead to the holomorphic complex-three-dimensional
refinement: $\cF$ is enriched by a compatible $\bar\partial$-action and
its structure maps vary holomorphically in configuration space. The
passage from the underlying real $E_6$ factorisation algebra to this
holomorphic $E_3$-type object is a comparison theorem, not a formal
identity.

Canonicality --- the statement that $\Phi^{\mathrm{FA}}_3 : \mathrm{CY}_3 \to \mathrm{Fact}^{E_3, \mathrm{hol}}_X$ lifts uniquely up to contractible choice --- is therefore conditional on Kontsevich--Tamarkin formality in the holomorphic refinement together with the Costello--Gwilliam locality axiom (CG FA Vol~2 Theorem~8.6.9, the holomorphic refinement of the locality theorem of Costello 2011 Chapter~5 paired with the cylindrical-end compact-exhaustion of Costello 2011 Theorem~13.4.1). The locality theorem requires a compact exhaustion of $X^{\mathrm{top}}$ with uniformly bounded heat-kernel propagators: for a smooth projective $\mathrm{CY}_3$ this is automatic (the projective embedding supplies a compact Kähler metric); on non-compact complex-threefold targets such as $K3 \times \Cbb$ it requires cylindrical-end compatibility on the affine factor together with the twisted holomorphic extension of CG FA Vol~2 Theorem~8.6.9 on the K3 factor. When the compactness hypothesis fails without cylindrical-end control, the stage-one output is a prefactorisation algebra not yet proven to satisfy Weiss-descent, the identification with the holomorphic $E_3$ structure holds only in the homotopy category of prefactorisation algebras, and canonicality of Theorem~\ref{thm:main}(iv) reduces to the conditional status of these hypotheses. Every occurrence of $\Phi^{\mathrm{FA}}_3$ below is read with the hypotheses and structure maps of this subsection in place.

\medskip
\noindent\textbf{Both stages as $\mathrm{Pr}^{\mathrm{L}}_{k,\mathrm{st}}$-derived functors.} Stage one is $\Phi^{\mathrm{FA}}_3 : \cD^{b}_{\mathrm{dg}}(X) \to \mathrm{Fact}^{E_3}_X$, a morphism in $\mathrm{Pr}^{\mathrm{L}}_{k,\mathrm{st}}$: $\cD^{b}_{\mathrm{dg}}(X)$ is compactly generated by $\mathrm{Perf}(X)$ (Bondal--Van den Bergh 2003 Theorem~3.1.1), Costello--Gwilliam--Li 2021 locality pins the value on each compact generator, Kontsevich--Tamarkin 2000 $E_d$-formality sends perfect complexes to compact $E_3$-factorisation algebras, and Lurie $\mathrm{HA}$ Corollary~5.5.2.9 extends to the unique colimit-preserving functor from the Ind-completion; $\Phi^{\mathrm{FA}}_3$ is the left adjoint in a colimit-preserving adjunction (Lurie $\mathrm{HA}$ Theorem~5.5.3.18). Stage two $\Sp^{\mathrm{ch}}_{\Sigma_2, C} = \mathrm{res}_C \circ \int_{\Sigma_2} : \mathrm{Fact}^{E_3}_X \to \mathrm{Fact}^{E_1}_{X/\Sigma_2} \to \mathrm{ChirAlg}^{\mathrm{ch}}_C$ is the composite of factorisation homology (Ayala--Francis 2015 Definition~3.1) and restriction to $C$; both legs preserve colimits. The composite $\Phi_3 : \cD^{b}_{\mathrm{dg}}(X) \to \mathrm{ChirAlg}^{\mathrm{ch}}_C$ is a $1$-morphism in $\mathrm{Pr}^{\mathrm{L}}_{k,\mathrm{st}}$; a single CY-$d$ category admits a family of $E_1$-chiral shadows indexed by $(\Sigma_{d-1}, C)$.

% ==================================================
\section{Comparison targets from the three volumes}
\label{sec:climaxes}
% ==================================================

This section records imported theorem surfaces and conditional targets.
It does not prove the Atiyah--Connes obstruction; the construction of
the obstruction begins in Section~\ref{sec:cocycle}, and the comparison
with these targets is exactly Theorem~\ref{thm:main}.

\subsection{Volume I -- Theorem B, chiral Positselski}
Let $A$ be an augmented chiral algebra on the curve $C$ with augmentation ideal $\bar A$. The ordered chiral bar complex is $\Bar^{\mathrm{ch}}_C(A) = T^c_C(s^{-1}\bar A)$, a cofree ordered chiral coalgebra with Maurer--Cartan differential
\[
  D \cdot \Theta_A + \tfrac{1}{2}[\Theta_A, \Theta_A] = 0
\]
governed by the universal obstruction $\Theta_A$ in the ordered convolution dGLA.

\begin{theorem}[Chiral Positselski, Volume~I, Theorem~B; \ClaimStatusProvedElsewhere]
\label{thm:climax-I}
For $A$ an admissible chiral algebra on $C$ whose underlying chiral module lies in the Koszul stratum $\cU^{\mathrm{adm}}(X) \subset \cM_{\mathrm{cx}}(X)$, the chiral bar and cobar constructions form an adjoint equivalence
\[
  \Cobar^{\mathrm{ch}}_C \Bar^{\mathrm{ch}}_C(A) \;\xrightarrow{\sim}\; A \qquad\text{in}\qquad D^{\mathrm{co}}_{\mathrm{ch}}(C).
\]
The inverse equivalence is $\Bar^{\mathrm{ch}}_C \Cobar^{\mathrm{ch}}_C$ on the Koszul-comatched side.
\end{theorem}

The proof (Volume~I, §7) adapts Positselski's 2011 §6.5 contracting-homotopy construction: one writes an explicit chain homotopy $h$ on the relative tensor product $\Bar \otimes \Cobar$ satisfying $[d, h] = \id - p$, with $p$ the projection onto the Koszul-admissible subcomplex. The homotopy witness $h$ exists precisely when the completed modular cyclic obstruction classes vanish through all weights; on $\cU^{\mathrm{adm}}$ they do.

\subsection{Volume II -- the Dijkgraaf--Verlinde--Verlinde identity}
The Volume~II comparison target identifies the protected BPS-index
sector of the $K3\times T^2$ compactification with the reciprocal of
the Igusa cusp form $\Phi_{10}$.

\begin{theorem}[DVV protected index, Volume~II comparison target; \ClaimStatusProvedElsewhere]
\label{thm:climax-II}
For $Z \in \Hbb_2$ in the Siegel upper half-space of genus $2$,
\[
  Z^{\AdS_3 \times K3}_{\mathrm{3dQG, BPS}}(Z) \;=\; \frac{1}{\Phi_{10}(Z)},
\]
with $\Phi_{10}$ the unique cusp form of weight $10$ for the paramodular group $\Gamma_{10}$ (Gritsenko 1999). The identity holds for the twisted-BPS-index-refined partition function (Dijkgraaf--Verlinde--Verlinde 1997); the equivalent interpretation as $1/4$-BPS dyon-counting on type-IIB on $K3 \times T^2$ is Sen 2008.
\end{theorem}

The left-hand side \emph{equals} the generating function of $1/4$-BPS multiplicities in the corresponding CHL orbifold (Witten $1/4$-BPS lectures, Sen 2008), and the equality with $1/\Phi_{10}$ holds to all orders in the low-energy expansion.

\begin{theorem}[Physical identification of the DVV symbols; \ClaimStatusProvedElsewhere]
\label{thm:dvv-physics}
Each symbol in $Z^{\AdS_3 \times K3}_{\mathrm{3dQG, BPS}}(Z) = 1/\Phi_{10}(Z)$ is a physical quantity, and the identity equates these quantities.
\begin{itemize}
  \item $Z = \bigl(\begin{smallmatrix} \tau & z \\ z & \sigma \end{smallmatrix}\bigr) \in \Hbb_2$ is a period matrix in the genus-$2$ Siegel upper half-space $\Hbb_2 = \{Z \in \mathrm{Mat}_2(\Cbb) : Z^T = Z,\; \mathrm{Im}(Z) > 0\}$. The diagonal entries $Z_{11} = \tau$ and $Z_{22} = \sigma$ are the moduli of the left-moving (thermal) circle $S^1_\tau$ and the right-moving (angular-momentum) circle $S^1_\sigma$ in the $\AdS_3$ boundary torus $T^2 = S^1_\tau \times S^1_\sigma$; the off-diagonal entry $Z_{12} = z$ is the R-symmetry chemical potential coupling the two circles, which under electromagnetic duality maps to the mixing parameter between $L_0$ and $J_0$ on the $\AdS_3$ boundary CFT.
  \item $Z^{\AdS_3 \times K3}_{\mathrm{3dQG, BPS}}(Z) = \tr_{\cH^{\mathrm{BPS}}}\bigl[(-1)^F q^{L_0 - c/24} \bar q^{\bar L_0 - c/24} y^{J_0}\bigr]$ is the Witten index over $1/4$-BPS bulk states of type-IIB string theory on $K3 \times T^2$, with $(q, \bar q, y) = (e^{2\pi i \tau}, e^{2\pi i \sigma}, e^{2\pi i z})$.
  \item $\Phi_{10}(Z) = \prod_{(k, l, m) > 0} (1 - q^k y^l \bar q^m)^{c(4km - l^2)}$ is the Igusa cusp form realised as the Borcherds product over the positive-root lattice $\mathrm{II}_{2,1} \oplus \langle -2 \rangle$, with exponents $c(n)$ the Fourier coefficients of the $K3$ elliptic genus.
\end{itemize}
The equation $Z^{\AdS_3 \times K3}_{\mathrm{3dQG, BPS}} = 1/\Phi_{10}$ equates the bulk BPS Hilbert-space trace with the boundary automorphic denominator on the full doubly-elliptic period matrix $Z \in \Hbb_2$.
\end{theorem}

\begin{proof}
Under $Z = \bigl(\begin{smallmatrix} \tau & z \\ z & \sigma \end{smallmatrix}\bigr)$ the entries identify with the thermal modulus on the left factor, the angular-momentum modulus on the right factor, and the R-charge fugacity on the off-diagonal. Maldacena--Moore--Strominger 1999~\S3 identifies these three circle bundles with the three charges of a D1-D5-P black hole; the Witten-index character $\tr(-1)^F q^{L_0} \bar q^{\bar L_0} y^{J_0}$ over the protected Hilbert space is Sen 2008 equation (3.15). The Borcherds product factorisation of $1/\Phi_{10}$ is Gritsenko--Nikulin 1998 Theorem~2.1; the identification of the exponent function $c(n)$ with the $K3$ elliptic-genus Fourier coefficient is DVV 1997 equation (1.5). The bulk-boundary equality is the DVV identity.
\end{proof}

\begin{remark}[BPS-index scope versus off-shell gravitational path integral]
\label{rem:bps-off-shell}
Theorem~\ref{thm:climax-II} holds as an identity of BPS Witten indices, not as an identity of off-shell partition functions. States contributing to the trace are $1/4$-BPS, annihilated by half the $\cN = 4$ supercharges, and the insertion $(-1)^F$ cancels non-BPS multiplets in short/long pairs (Bhattacharyya--Sen 2008, \S2.3); what survives is the rigid cohomological invariant, the Witten index of the supersymmetric Hilbert space, not the full path-integral measure. The full off-shell 3D gravity path integral $Z^{\mathrm{full}}_{\mathrm{3dQG}}$ receives contributions from non-BPS black holes, from $\AdS_3$-fluctuation modes beyond the mini-superspace reduction, and from sums over $\mathrm{SL}_2(\Zbb)$ images of bulk saddles (Maloney--Witten 2007); these contributions are not captured by $1/\Phi_{10}$. The equality is exact in the supergravity large-charge limit (Sen 2008 \S1); at finite charge, logarithmic corrections (Banerjee--Gupta--Sen 2011) modify the leading-order Cardy formula and are captured by $1/\Phi_{10}$ through the Fourier-coefficient asymptotics of the unique weight-$10$ Siegel cusp form.
\end{remark}

\begin{theorem}[DVV as the K3-boundary protected-index stratum; \ClaimStatusConditional]
\label{thm:dvv-as-boundary}
Conditional on the stage-two specialisation and character comparison of Theorem~\ref{thm:main}\textup{(iv)--(v)} (principal K3 locus and trace normalisation), the DVV identity is the $K3$-boundary protected-index stratum of the comparison. Evaluating that conditional bridge and taking the Witten index produces $\chi_{\mathrm{BPS}}\bigl(\Phi_3(K3 \times T^2)\bigr) = 1/\Phi_{10}$.
\end{theorem}

\begin{proof}
Stage one $\Phi^{\mathrm{FA}}_3(K3 \times T^2)$ is the $E_3$-factorisation algebra of the type-IIB $K3 \times T^2$ compactification; stage-two specialisation $\Sp^{\mathrm{ch}}_{K3, T^2}$ integrates over the $K3$ fibre and restricts to the elliptic-curve factor. Under the Vol.~III principal-locus source-recognition gate, the stage-two target is compared with $\bfH_{\Delta_5}$ and the genus-one character is compared with the R-charge-refined Witten index on the left side of the DVV identity. Without that gate only the scalar automorphic target $1/\Phi_{10}$ remains, whose Borcherds denominator is Borcherds 1998 Theorem~10.1.
\end{proof}


\subsection{Volume III -- $\bfH_{\Delta_5}$ source-recognition gate}
The Volume~III input is a conditional principal-locus
$K3\times E$ specialisation, not a bare CY3 functor applied to the
surface $K3$. Let $Y=S\times E$ with $S$ a projective K3 surface on
the Mukai-principal locus and $E$ an elliptic curve. In this
comparison, $\bfH_{\Delta_5}$ appears only as the Vol.~III
source-recognition target for the stage-two chiral specialisation of
the CY3 object
$\mathrm{Perf}(Y)$ on the elliptic factor.

\begin{theorem}[$\bfH_{\Delta_5}$ source-recognition target, Volume~III; \ClaimStatusConditional]
\label{thm:climax-III}
Assume the Vol.~III principal-locus source-recognition hypotheses:
Kontsevich--Tamarkin $E_3$-formality, Costello--Gwilliam--Li locality,
Ayala--Francis factorisation-homology Fubini, the BL-2/4/6 comparison
witnesses, and the BL-7 Borcherds-weight normalisation. Then the
comparison target is the equivalence of $E_1$-chiral bialgebras on $E$
\[
  \Sp^{\mathrm{ch}}_{S,E}\!\left(\Phi^{\mathrm{FA}}_3
  (\mathrm{Perf}_{\mathrm{dg}}(S\times E))\right)
  \;\simeq\; \bfH_{\Delta_5}.
\]
The scalar targets split as
\[
  \kappa_{\mathrm{BKM}}(\Delta_5)=5,\qquad
  \Phi_{10}=\Delta_5^2,\qquad
  K_{\mathrm{Muk}}=2c_+(\widetilde H(S,\Zbb))=8.
\]
The value $8$ is the Mukai conductor of the chiral bialgebra, not the
BKM weight and not a proof that $\alpha_{S\times E}=0$.
\end{theorem}

The source-recognition package consists of the Vol.~III
$E_3$-formality/locality input under the principal-locus hypotheses
and the Ayala--Francis factorisation-homology Fubini comparison
(Ayala--Francis 2015, Theorem~3.16).

\begin{theorem}[Physical interpretation of $\bfH_{\Delta_5}$ as $K3 \times E$ compactification; \ClaimStatusConditional]
\label{thm:hdelta5-physics}
Assume the source-recognition gate of Theorem~\ref{thm:climax-III}. Then the K3 chiral bialgebra $\bfH_{\Delta_5}$ is the comparison target for the worldsheet chiral-algebra sector of type-IIA string theory compactified on $K3 \times E$, with the following physical scalar identifications:
\begin{itemize}
  \item The central charge $K = 8$ equals twice the positive-signature rank of the Mukai lattice $\Lambda_{K3} = H^{\mathrm{even}}(K3, \Zbb) = 2U \oplus 2(-E_8) = \mathrm{II}_{4, 20}$ (Mukai 1987), whose positive-definite sub-lattice $\Lambda^+_{K3}$ has signature $(4, 0)$: four self-dual bosons $(h^{0,0} + h^{1,1}_{\mathrm{alg}} + h^{2,0} + h^{2,2}) = (1 + 1 + 1 + 1) = 4$ at the Narain-locus $\mathrm{II}_{4, 4} \subset \mathrm{II}_{4, 20}$. The factor $2$ in $K = 2 c_+(L) = 8$ is the Bruinier normalisation (Bruinier 2002 Proposition~5.1) pairing holomorphic and antiholomorphic sectors in the Heegner--Chern reciprocity.
  \item The trace parameter $\hbar^2 = -1/8$ is the tree-level Wald entropy normalisation on the attractor geometry $\AdS_2 \times S^2 \times K3$ of a $1/4$-BPS type-IIB black hole (Sen 2008 equation~(4.17)), equal to the Kontsevich--Vlassopoulos 2022 Calabi--Yau-$3$ Serre-duality constraint $\tr(1) = -1$ divided by $K = 8$.
  \item The Mukai discriminant $\Delta_5$ labels the orbit of the $K3$ Picard lattice under the $\mathrm{O}(\mathrm{II}_{4, 20})$ action. Singular $K3$ surfaces at Picard rank $20$ are classically indexed by positive-definite binary quadratic forms and the Heegner--Stark--Baker class-number-one discriminants $\{3, 4, 7, 8, 11, 19, 43, 67, 163\}$; the rank-$8$ arithmetic refinement in Proposition~\ref{prop:heegner} is conditional.
\end{itemize}
The physical comparison states that type-IIA on $K3 \times E$ with the standard Narain embedding of $\Lambda_{K3}$ into the $\mathrm{II}_{4, 20}$ worldsheet lattice has, at the source-recognised elliptic-fibre target, central charge $K = 8$, trace parameter $\hbar^2 = -1/8$, and $1/4$-BPS scalar partition function $1/\Phi_{10}$.
\end{theorem}

\begin{proof}
Type-IIA on $K3$ reduces to a $(1+5)$-dimensional effective theory with $24$ hypermultiplets and $1$ graviphoton (Hull--Townsend 1995); compactifying on $E$ gives a three-dimensional theory with $\cN = 4$ supersymmetry, whose moduli space is the Mukai--Narain locus $\mathrm{II}_{4, 20}$ (Mukai 1987). Under the source-recognition hypothesis, the elliptic-fibre comparison target is the Heisenberg vertex algebra on $\Lambda_{K3}$ together with the Borcherds scalar denominator $\Phi_{10}$ (Borcherds 1998 \S10, applied to the $K3$-elliptic genus). Its positive-definite rank is $c_+(\Lambda_{K3}) = 4$, giving $K = 2 \cdot 4 = 8$. The trace parameter is forced by the Kontsevich--Vlassopoulos 2022 Calabi--Yau-$3$ normalisation $\tr(1) = -1$ to be $\hbar^2 = -1/K = -1/8$. The $\Delta_5$ discriminant is used only through the scalar target $\Phi_{10}=\Delta_5^2$ and the Vol.~III gate. Shioda--Inose $1977$ combined with the Gauss--Heegner--Stark--Baker class-number-one theorem supplies the classical singular-K3 arithmetic background. The proposed rank-$8$ refinement of Proposition~\ref{prop:heegner} is a separate conditional arithmetic problem, not an input to the $\Delta_5$ designation.
\end{proof}

\begin{remark}[Attractor mechanism and $\bfH_{\Delta_5}$ central charge]
\label{rem:attractor}
In the type-IIB frame on $K3 \times T^2$, Sen 2008 \S4.3 shows the attractor geometry at a BPS black-hole horizon is $\AdS_2 \times S^2 \times K3 \times \tilde T^2$ with $\tilde T^2$ the dual torus. The horizon-circle chiral algebra is compared with $\bfH_{\Delta_5}$ only after the Vol.~III source-recognition gate. The attractor value of the dilaton fixes $g_s^2 = 1/K$ and hence $\hbar^2 = -1/K$; the Bruinier 2002 Proposition~5.1 Heegner--Chern reciprocity then becomes the scalar constraint that $K$ equals twice the positive-signature rank of the Mukai lattice, $K = 2c_+(L) = 8$.
\end{remark}


\subsection{Trace comparison target}
The four-archetype classification places $\sfB$ as a fifth row with ceiling $K^\kappa = 8$.

\begin{theorem}[Universal trace identity; \ClaimStatusConditional]
\label{thm:climax-IV}
On the derived centre $Z^{\mathrm{der}}_{\mathrm{ch}}(A_X)$ of the chiral algebra $A_X = \Phi_3(X)$ for $X$ a smooth projective Calabi--Yau threefold with Mukai lattice $L$, and on the Koszul-admissible locus $\cU^{\mathrm{adm}}(X)$,
\[
  \hbar^2 \cdot K \;=\; -1, \qquad K \;=\; 2\,c_+(L),
\]
where $c_+(L)$ denotes the positive-signature rank of the Mukai lattice $L \subset H^\bullet(X, \Zbb)$. The identification $K = 2c_+(L)$ is forced by Bruinier 2002 Proposition~5.1 (Heegner--Chern reciprocity for the Borcherds product) applied to the derived-centre pairing; the normalisation $\hbar^2 K = -1$ is the programme's master identification (Volume~III, \S9), conditional on canonicality of the stage-two specialisation $\Sp^{\mathrm{ch}}_{\Sigma_2, C}$.
\end{theorem}

For the K3 Mukai lattice, $c_+(L) = 4$ giving $K = 8$, hence
$\hbar^2 = -1/8$, compatible with the principal $K3\times E$
specialisation of Theorem~\ref{thm:climax-III}. For $X = Q_5$, $c_+$
counts the positive part of the integral third cohomology; the master
trace on $Q_5$ would require the Atiyah--Connes class to vanish, a
condition not known on the generic quintic and not implied by the
Yukawa detector.

\begin{theorem}[Three-scale physical reading of $\hbar^2 K = -1$; \ClaimStatusHeuristic]
\label{thm:three-scale}
The master identity $\hbar^2 K = -1$ admits three distinct physical readings, one per scale at which the derived-centre trace is probed.
\begin{enumerate}[label=\textup{(\arabic*)}]
  \item \emph{Microscopic.} On the ultraviolet scale of the Calabi--Yau target $X$, $K$ is the rank of the positive-signature sub-lattice $L^+ \subset L$ of the Mukai lattice, counted with Bruinier's Heegner--Chern normalisation factor $2$; each rank-$1$ positive direction contributes a chiral boson $\phi^i$ of $U(1)$-current weight $1$ on the worldsheet. The parameter $\hbar$ is the Planck constant of the target-space BV complex (Costello 2011 \S5.5), and $\hbar^2 = -1/K$ records the tree-level normalisation of the holomorphic kinetic term on the Mukai-lattice worldsheet. The equation $\hbar^2 K = -1$ is the worldsheet modular-invariance constraint on the Heegner--Chern reciprocity (Bruinier 2002 Proposition~5.1).
  \item \emph{Mesoscopic.} On the factorisation-homology scale, where the two-cycle $\Sigma_2 \subset X$ carries the stage-two specialisation, $K$ is the target central charge assigned to the Vol.~III $\bfH_{\Delta_5}$ comparison only under the source-recognition gate, and $\hbar^2$ is the effective-theory chiral coupling $\alpha_{\mathrm{eff}}^2$ set by Ayala--Francis 2015 Fubini: the $\Sigma_2$-integration propagates the $X$-level scalar constraint $\hbar^2 K = -1$ to the $C$-level target.
  \item \emph{Macroscopic.} On the derived-centre scale, $Z^{\mathrm{der}}_{\mathrm{ch}}(A_X) \to \Cbb[-3]$, the trace $\tr: Z^{\mathrm{der}}_{\mathrm{ch}} \to \Cbb[-3]$ normalises the Calabi--Yau-$3$ Serre pairing via Kontsevich--Vlassopoulos 2022 Theorem~1.2: $\tr(1) = -1$ is the Serre-duality normalisation of the canonical module on a CY-$3$. Bruinier 2002 Proposition~5.1 Heegner--Chern reciprocity then forces $\hbar^2 K = -1$ as the unique extension of $\tr(1) = -1$ compatible with a rank-$K$ trace.
\end{enumerate}
The three scales communicate: the microscopic string-frame normalisation descends to the mesoscopic chiral-algebra central-charge constraint under factorisation homology, which descends to the macroscopic Serre-trace normalisation under the derived-centre functor $Z^{\mathrm{der}}_{\mathrm{ch}}$. The single equation $\hbar^2 K = -1$ is the constant value of this descent chain.
\end{theorem}

\begin{proof}
\emph{Microscopic.} The $K3$ worldsheet superconformal field theory has $c_R = 6$ on the right-moving $\cN = 4$ sector (the $K3$ sigma-model central charge, Witten 1987) and $c_L = 24$ on the left-moving bosonic Mukai-lattice sector after Narain embedding (Narain 1986); the Mukai-enhanced chiral subsector at positive signature has $K = 2 c_+(L) = 8$ chiral bosons. The equation $\hbar^2 K = -1$ equates the target-space BV coupling $\hbar^2$ with $-1/K$, consistent with the Heegner--Chern reciprocity of Bruinier 2002 Proposition~5.1 which forces $\tr(1) = -1$ at the vacuum of the Mukai-enhanced chiral algebra. \emph{Mesoscopic.} Stage-two specialisation $\Sp^{\mathrm{ch}}_{\Sigma_2, C}$ is factorisation homology over the $2$-cycle $\Sigma_2$ (Ayala--Francis 2015 Theorem~3.16 Fubini), which by Definition~2.3 of that paper preserves the trace on the Calabi--Yau-$d$ category up to the shifted sign $(-1)^d$; for $d = 3$ the sign is $-1$, and the descended central charge is $K = 2c_+(L)$ with $\hbar^2 K = -1$ inherited. \emph{Macroscopic.} Kontsevich--Vlassopoulos 2022 Theorem~1.2 gives $\tr \circ \mathbb{S} = (-1)^d \tr$ for a Calabi--Yau-$d$ category; at $d = 3$, $\tr(1) = -\tr(1)$ on the non-vacuum part, forcing $\tr(1) = -1$ on the vacuum. The identity $\hbar^2 K = -1$ is the extension of $\tr(1) = -1$ to a rank-$K$ family via Bruinier 2002 Proposition~5.1.
\end{proof}

\begin{remark}[Renormalisation-group flow of $\hbar^2 K$]
\label{rem:rg-flow}
The three scales correspond to three points of a renormalisation-group flow. The microscopic value $\hbar^2 = -\alpha'$ is the ultraviolet coupling; the mesoscopic value $\hbar^2 = \alpha_{\mathrm{eff}}^2$ is the effective-theory coupling at the factorisation-homology scale $\Sigma_2$; the macroscopic value $\hbar^2 = -1/K$ is the infrared attractor value at the derived-centre trace. The product $\hbar^2 K$ is a renormalisation-group invariant: at every scale it equals $-1$. This is the programme's analogue of Zamolodchikov's $c$-theorem (Zamolodchikov 1986): a monotonic quantity along the renormalisation-group flow, fixed to $-1$ by the derived-centre trace normalisation, ensuring the flow terminates at the Calabi--Yau-$3$ infrared fixed point.
\end{remark}


% ==================================================
\section{Construction of the cocycle $[m_3, B^{(2)}]_X$}
\label{sec:cocycle}
% ==================================================

A cyclic-HKR datum assembles the five pieces that turn
$[m_3,B^{(2)}]_X$ from notation into a cochain of total Hochschild
degree four landing in $H^2(X,\Omega^1_X)$: a cyclic dg enhancement
with Calabi--Yau trace, the cochain transport $\Delta_2$ of Connes'
chain operator, the Hochschild--Kostant--Rosenberg projection to the
polyvector summand, and the contraction against the volume form. Its
existence is a theorem on a smooth projective Calabi--Yau threefold.

\begin{definition}[Cyclic-HKR datum]
\label{def:cyclic-hkr-datum}
Let $X$ be a smooth projective Calabi--Yau threefold with holomorphic
volume form $\Omega_X$. A \emph{cyclic-HKR datum}
$\mathfrak c_X$ consists of:
\begin{enumerate}[label=\textup{(AC\arabic*)}]
  \item a cyclic dg enhancement
  $\mathrm{Perf}_{\mathrm{dg}}(X)$ of $D^b(\Coh X)$ with
  Calabi--Yau trace
  \[
    \Tr_{\mathrm{cyc}}:\HC^\bullet(\Perf_{\mathrm{dg}}(X))\to k[-3],
  \]
  and its Hochschild cochain complex
  \[
    C^\bullet_X=
    C^\bullet\bigl(\Perf_{\mathrm{dg}}(X),
    \Perf_{\mathrm{dg}}(X)\bigr)
  \]
  with differential $\delta_{\mathrm{Hoch}}$;
  \item the cochain Connes rotation
  \[
    \Delta_2:C^\bullet_X\to C^{\bullet+1}_X,
  \]
  defined as the transport under $\Tr_{\mathrm{cyc}}$ of Connes' chain
  operator $B=(1-\lambda)sN$ along the Shklyarov pairing
  $\HH_\bullet(X)\otimes\HH^\bullet(X)\to k[-3]$
  (Shklyarov 2013, Theorem~1.3; Brav--Dyckerhoff 2019
  \S3) with the sign convention
  $\Delta_2\delta_{\mathrm{Hoch}}+\delta_{\mathrm{Hoch}}\Delta_2=0$;
  \item the ternary Stasheff operation
  $m_3:(C^\bullet_X)^{\otimes 3}\to C^{\bullet-1}_X$, the image of the
  pentagonal cell $[K_3]\in C_1(\cA_\infty(3))$ in the framed
  $E_2$-Gerstenhaber--BV structure on $C^\bullet_X$ supplied by
  Deligne's conjecture (Kontsevich--Soibelman 2000 Theorem~2.4);
  \item the Hochschild--Kostant--Rosenberg projection
  \[
    \pi^{\mathrm{HKR}}_{2,2}:
    \HH^4(X)=\bigoplus_{p+q=4}H^p(X,\wedge^qT_X)
    \twoheadrightarrow H^2(X,\wedge^2T_X)
  \]
  along the $(2,2)$ summand (Caldararu 2005 Theorem~1.2);
  \item the Calabi--Yau contraction
  \[
    \iota_{\Omega_X}:H^2(X,\wedge^2T_X)\to H^2(X,\Omega^1_X),
    \qquad
    \beta\mapsto\iota_\beta\Omega_X,
  \]
  an isomorphism of sheaves induced by $\Omega_X$ via
  $\wedge^2T_X\cong\Omega^1_X$.
\end{enumerate}
\end{definition}

\begin{theorem}[Existence of a cyclic-HKR datum; \ClaimStatusProvedHere]
\label{thm:cyclic-hkr-datum-exists}
Every smooth projective Calabi--Yau threefold $X$ with fixed
holomorphic volume form $\Omega_X$ admits a cyclic-HKR datum
$\mathfrak c_X$. Any two such data are related by a cyclic
$A_\infty$-quasi-equivalence of $\Perf_{\mathrm{dg}}(X)$ intertwining
$\Delta_2$ and $m_3$, and the resulting class in $H^2(X,\Omega^1_X)$
is independent of the choice.
\end{theorem}

\begin{proof}
Existence of the cyclic dg enhancement is Lunts--Orlov 2010
Theorem~2.8 together with the Calabi--Yau trace construction of
Shklyarov 2013 Theorem~1.3; uniqueness up to cyclic
$A_\infty$-quasi-equivalence is Brav--Dyckerhoff 2019 Corollary~3.18
(equivalently, Kontsevich--Vlassopoulos 2022 Theorem~1.2 applied to
the CY-$3$ target $X$). The Connes--Tsygan chain operator $B$ on
$\HC_\bullet(\Perf_{\mathrm{dg}}(X))$ is Connes 1985 \S3.1; its transport
to a cochain operator $\Delta_2$ on $C^\bullet_X$ via the Shklyarov
pairing is the Kontsevich--Vlassopoulos 2022 construction, and the sign
convention $\Delta_2\delta_{\mathrm{Hoch}}+\delta_{\mathrm{Hoch}}\Delta_2=0$
is fixed by the Calabi--Yau trace degree $-3$. The
$E_2$-Gerstenhaber--BV structure on $C^\bullet_X$ giving $m_3$ is
Deligne's conjecture (McClure--Smith 2002 / Kontsevich--Soibelman 2000
Theorem~2.4). The HKR isomorphism is Caldararu 2005 Theorem~1.2, and
$\wedge^2T_X\cong\Omega^1_X$ by contraction against $\Omega_X$ is
linear algebra on a CY-$3$. Independence of $\alpha_X$ is the content
of Proposition~\ref{prop:atiyah-cup-derivation} below.
\end{proof}

\begin{definition}[Atiyah--Connes class]
\label{def:atiyah-connes-class}
Fix a cyclic-HKR datum $\mathfrak c_X$. The graded commutator
\[
  \gamma_X=[m_3,\Delta_2]
  =m_3\Delta_2+\Delta_2 m_3
\]
is an element of $C^4_X$ of total Hochschild degree $4$. Put
\[
  \beta_X=\pi^{\mathrm{HKR}}_{2,2}\bigl([\gamma_X]\bigr)
  \in H^2(X,\wedge^2T_X),
  \qquad
  \alpha_X=\iota_{\Omega_X}(\beta_X)\in H^2(X,\Omega^1_X).
\]
The symbol $B^{(2)}$ denotes $\Delta_2$ throughout; the superscript
records the bidegree at which the cochain operator is used.
\end{definition}

\subsection{The Dolbeault resolution of the tangent complex}
Since $X$ is Kähler (assumed throughout this section, with the Clemens case deferred to §\ref{sec:clemens}), the tangent bundle $T_X$ admits a Dolbeault resolution
\[
  T_X \;\hookrightarrow\; \bigl(\Omega^{0,\bullet}(X, T_X),\; \bar\partial\bigr).
\]
Kapranov's Proposition~4.4 realises $\at_X$ as the class of $\bar\partial \nabla$ in $\Omega^{0,1}(X, \End(T_X) \otimes \Omega^1_X)$, where $\nabla$ is the Chern connection of a Kähler metric. The $L_\infty$-brackets $\ell_n$ on $T_X[-1]$ (Kapranov 1999) satisfy
\[
  \ell_n \;=\; \at_X^{[n-1]} \;\in\; \Omega^{0, n-1}(X, T_X \otimes (\Omega^1_X)^{\otimes (n-1)}).
\]

\subsection{The ternary bracket and the cyclic shift}
Take $\cE = T_X$; the construction extends to any tilting bundle giving Kuznetsov's exceptional collection on $D^b(X)$. The $A_\infty$-structure on $\End^\bullet(\cE)$ carries brackets
\[
  m_n \;:\; \End^\bullet(\cE)^{\otimes n} \;\longrightarrow\; \End^\bullet(\cE)[2-n]
\]
with $m_1 = d$, $m_2$ composition, and $m_n$ for $n \geq 3$ the Massey brackets. The ternary bracket $m_3$ measures the failure of $(f \circ g) \circ h = f \circ (g \circ h)$ at cochain level. For cocycles $f, g, h$ whose pairwise products are exact --- $f \circ g = d\alpha$, $g \circ h = d\beta$ ---  each of the two possible groupings $(f \circ g) \circ h$ and $f \circ (g \circ h)$ is a boundary, but of two \emph{different} chain-level primitives: $d(\alpha \circ h)$ and $d(f \circ \beta)$. Their difference is a cocycle,
\[
  m_3(f,g,h) \;=\; f \circ \beta - \alpha \circ h,
\]
and its class in cohomology is the Massey triple product --- the obstruction to choosing the two primitives coherently. The sign convention follows Stasheff 1963: the graded Leibniz rule $d(\alpha \circ h) = (d\alpha) \circ h \pm \alpha \circ dh$ forces the difference to be $d$-closed precisely because $dh = 0$ and $d\alpha = f \circ g$.
Connes' cyclic shift at bidegree $2$ acts on $(\HH^\bullet, d_{\HH})$ by
\[
  B^{(2)} : \HH^2(\End \cE) \;\longrightarrow\; \HH^3(\End \cE),
\]
increasing degree by one and cyclic-weight by one. On HKR decomposables,
\[
  B^{(2)}(f_0 \otimes f_1 \otimes f_2) \;=\; \sum_{\sigma \in \Zbb/3} (-1)^\sigma \sigma \cdot (1 \otimes f_0 \otimes f_1 \otimes f_2).
\]

\subsection{The commutator and its cyclic-Hochschild home}
\begin{construction}[Atiyah--Connes cochain; \ClaimStatusProvedHere]
\label{constr:cocycle}
Let $\cE=T_X$ and let $\mathfrak c_X$ be a cyclic-HKR datum
(Definition~\ref{def:cyclic-hkr-datum},
Theorem~\ref{thm:cyclic-hkr-datum-exists}). The ternary Stasheff
operation $m_3$ lives in the $\cA_\infty(3)\hookrightarrow fE_2(3)$
part of the framed little-$2$-discs action on Hochschild cochains
(Getzler--Jones 1994 Theorem~2.1, Kontsevich--Soibelman 2000
Theorem~2.4); the cochain Connes rotation $\Delta_2$ is the
bidegree-two component of the $SO(2)$-rotation. The graded commutator
\[
  \gamma_X=[m_3,\Delta_2]:
  \End^\bullet(\cE)^{\otimes 3}\longrightarrow\End^\bullet(\cE)
\]
has total Hochschild degree $4$. The framed-$E_2$ Cartan relations of
Getzler 1994 Proposition~1.2 give
\[
  \delta_{\mathrm{Hoch}}\gamma_X
  =[\delta_{\mathrm{Hoch}},[m_3,\Delta_2]]
  =[[\delta_{\mathrm{Hoch}},m_3],\Delta_2]
  +[m_3,[\delta_{\mathrm{Hoch}},\Delta_2]]
  =0,
\]
because $m_3$ is a coderivation of the bar complex (so
$[\delta_{\mathrm{Hoch}},m_3]=0$ on the Stasheff $\cA_\infty(3)$-part)
and $\Delta_2\delta_{\mathrm{Hoch}}+\delta_{\mathrm{Hoch}}\Delta_2=0$
by the cyclic sign convention
(Theorem~\ref{thm:cyclic-hkr-datum-exists}). Hence $[\gamma_X]$
defines a class in $\HH^4(X)$; the HKR isomorphism (Caldararu 2005
Theorem~1.2) decomposes it as
\[
  [\gamma_X]
  \in\HH^4(X)
  \cong\bigoplus_{p+q=4}H^p(X,\wedge^qT_X).
\]
The $(2,2)$-projection is the Kapranov square
\[
  \pi^{\mathrm{HKR}}_{2,2}[\gamma_X]
  =\beta_X=\at_X^{[2]}
  \in H^2(X,\wedge^2T_X)
\]
(Kapranov 1999 Proposition~4.4, with cyclic symmetrisation $\Delta_2$
reducing $\ell_3=\at_X^{[2]}$ to its cyclic-skew component). The
Calabi--Yau contraction sends this to
\[
  \alpha_X=\iota_{\Omega_X}\beta_X
  =\at_X\cup B^{(2)}_{\mathrm{Connes}}
  \in H^2(X,\Omega^1_X).
\]
Corollary~\ref{cor:hodge} identifies the target with $H^{1,2}(X)$ on
K\"ahler $X$.
\end{construction}

\begin{proposition}[Explicit Dolbeault representative; \ClaimStatusProvedHere]
\label{prop:cocycle-rep}
Let $\nabla$ be the Chern connection of a K\"ahler metric on $T_X$.
The class $\alpha_X$ is represented in Dolbeault cohomology by
\[
  \alpha_X
  =\bigl[\,\iota_{\Omega_X}
    \tr_{\End T_X}\bigl(
      \bar\partial\nabla\wedge\bar\partial\nabla
    \bigr)\bigr]
  \in H^2(X,\Omega^1_X),
\]
where the trace $\tr_{\End T_X}$ contracts the $\End(T_X)^{\otimes 2}$
factors to $k$ and $\iota_{\Omega_X}$ identifies
$\wedge^2\Omega^1_X\cong T_X\otimes\Omega_X^3$ with $\Omega^1_X$ via
the volume form. The class is independent of the connection and of
the K\"ahler metric.
\end{proposition}

\begin{proof}
Kapranov 1999 Proposition~4.4 identifies
$\at_X=[\bar\partial\nabla]\in H^1(X,\End(T_X)\otimes\Omega^1_X)$;
the Dolbeault representative of
$\at_X^{[2]}=\at_X\cup\at_X\in
H^2(X,\End(T_X)^{\otimes 2}\otimes(\Omega^1_X)^{\otimes 2})$
is the wedge $\bar\partial\nabla\wedge\bar\partial\nabla$ (Kapranov
1999 Theorem~4.6). Applying $\tr_{\End T_X}$ lands the class in
$H^2(X,(\Omega^1_X)^{\otimes 2})$; skew-symmetrisation in the two
$\Omega^1$-factors takes the cyclic-skew summand to
$H^2(X,\Omega^2_X)$. Contraction against $\Omega_X$ identifies
$\Omega^2_X\cong T_X\otimes\Omega^3_X\cong T_X$ and then with
$\Omega^1_X$ again via $\iota_{(-)}\Omega_X$; the composite map agrees
with the Construction~\ref{constr:cocycle} projection
$\pi^{\mathrm{HKR}}_{2,2}$ followed by the Calabi--Yau contraction
$\iota_{\Omega_X}$, as can be checked on HKR decomposables. Connection
independence follows from the description of $\at_X$ as the Ext class
of the jet sequence
\[
  0\to T_X\otimes\Omega^1_X\to J^1T_X\to T_X\to 0
\]
(Illusie 1971 Chapter~IV \S2.3, Markarian 2009 Theorem~1.1): the
Dolbeault representative $\bar\partial\nabla$ of this Ext class is
unique up to $\bar\partial$-exact terms, and the quadratic expression
$\bar\partial\nabla\wedge\bar\partial\nabla$ differs between two
choices $\nabla,\nabla'$ by a sum of cross terms
$\bar\partial(\nabla-\nabla')\wedge\bar\partial\nabla'
+\bar\partial\nabla\wedge\bar\partial(\nabla-\nabla')$, each of
which is $\bar\partial$-exact after tracing against the other factor.
Metric independence of the Atiyah class is standard (Atiyah 1957 \S4).
\end{proof}

\begin{corollary}[Hodge location; \ClaimStatusProvedHere]
\label{cor:hodge}
The cocycle $[m_3, B^{(2)}]_X$ lives in
\[
  H^2(X, \Omega^1_X) \;\cong\; H^{1,2}(X).
\]
The numerically equal $H^{2,1}(X)$ appears only after Hodge symmetry or Calabi--Yau duality; it is not the same Hodge summand. For $X = K3 \times E$ the hypothesis $h^{1,0}=0$ fails and any Calabi--Yau-threefold statement must pass through a specified Borcea--Voisin or related twist with recomputed Hodge groups.
\end{corollary}

\begin{proof}
For a compact Kähler manifold, the Dolbeault theorem identifies sheaf cohomology by
\[
  H^q(X,\Omega^p_X)\cong H^{p,q}(X).
\]
Taking $p=1$ and $q=2$ gives $H^2(X,\Omega^1_X)\cong H^{1,2}(X)$. Hodge symmetry gives equality of dimensions with $H^{2,1}(X)$, and the Calabi--Yau volume form gives the usual dual convention, but neither changes the sheaf-cohomology target.
\end{proof}

\subsection{Ten-step derivation of the Atiyah cup}
\label{sub:atiyah-cup}

\begin{proposition}[Atiyah cup identity, step-by-step derivation; \ClaimStatusProvedHere]
\label{prop:atiyah-cup-derivation}
Under the cyclic-HKR datum $\mathfrak c_X$, the HKR image of
$[m_3, B^{(2)}]_X$ equals $\at_X \cup B^{(2)}_{\mathrm{Connes}}$ in
$H^2(X, \Omega^1_X)$. Ten single-equation steps --- each citing a single
primary source --- chain Kapranov 1999 and Connes 1985 into the identity.
\begin{align*}
&\textup{(Step 1, Connes 1985 \S3.1, eq.~(1))}\quad && B \;=\; (1 - \lambda)\,s\,N, \\
&&& B(a_0, a_1, a_2) \;=\; 1{\otimes}a_0{\otimes}a_1{\otimes}a_2 - 1{\otimes}a_2{\otimes}a_0{\otimes}a_1 + 1{\otimes}a_1{\otimes}a_2{\otimes}a_0. \\
&\textup{(Step 2, bidegree-$2$ restriction)}\quad && B^{(2)} \;:=\; B\bigl|_{\HH^2} \;:\; \HH^2(\End T_X) \to \HH^3(\End T_X). \\
&\textup{(Step 3, Kapranov 1999 Prop.~4.4)}\quad && \at_X \;=\; [\bar\partial\nabla] \;\in\; H^1\!\bigl(X, \Omega^1_X \otimes \End T_X\bigr). \\
&\textup{(Step 4, Kapranov 1999 Thm.~4.6)}\quad && \ell_n^{\mathrm{Kap}} \;=\; (\bar\partial\nabla)^{[n-1]}, \quad \ell_3^{\mathrm{Kap}} \;=\; \mathrm{Alt}_{S_3}\bigl((\bar\partial\nabla)^{\wedge 2}\bigr). \\
&\textup{(Step 5, Kontsevich 2003 Thm.~4.6)}\quad && m_3 \;\simeq\; \ell_3^{\mathrm{Kap}} \text{ in }\mathrm{Op}^\infty_1 \text{ on }\End^\bullet(T_X). \\
&\textup{(Step 6, graded commutator)}\quad && [m_3, B^{(2)}] \;=\; m_3 B^{(2)} - B^{(2)} m_3 \;=\; \mathrm{Cyc}_3\bigl((\bar\partial\nabla)^{\wedge 2}\bigr). \\
&\textup{(Step 7, HKR, Caldararu 2005 Thm.~1.2)}\quad && \mathrm{HKR}\bigl[(\bar\partial\nabla)^{\wedge 2}\bigr] \;=\; \at_X \cup \at_X. \\
&\textup{(Step 8, Getzler--Jones 1994 Thm.~2.1)}\quad && \mathrm{Cyc}_3(\at_X \cup \at_X) \;=\; \at_X \cup B^{(2)}_{\mathrm{Connes}}. \\
&\textup{(Step 9, Caldararu--Willerton 2010 Thm.~1.6)}\quad && \tr_{\End T_X}\bigl(\at_X \cup B^{(2)}\bigr) \;=\; \at_X \cup B^{(2)}_{\mathrm{Connes}} \in H^2(X, \Omega^1_X). \\
&\textup{(Step 10, Hodge location, Cor.~\ref{cor:hodge})}\quad && \at_X \cup B^{(2)}_{\mathrm{Connes}} \;\in\; H^{1,2}(X).
\end{align*}
The composition of Steps~1--10 is exactly Construction~\ref{constr:cocycle}.
\end{proposition}

\begin{proof}
Step~1: Connes 1985 \S3.1, equation (1). Step~2: restriction. Step~3: Kapranov 1999 Proposition~4.4. Step~4: Kapranov 1999 Theorem~4.6. Step~5: Kontsevich 2003 Theorem~4.6, formality. Step~6: graded commutator in the framed $E_2$-Gerstenhaber algebra. Step~7: Caldararu 2005 Theorem~1.2 (HKR). Step~8: Getzler--Jones 1994 Theorem~2.1. Step~9: Caldararu--Willerton 2010 Theorem~1.6 (trace). Step~10: Corollary~\ref{cor:hodge}.
\end{proof}

\subsection{Motivic lift via the Beilinson regulator}
\label{sub:motivic-regulator}
The Dolbeault class is expected to have a motivic normal-function
origin. Beilinson 1984 constructs $H^p_{\cM}(X, \Qbb(q))$ as the
$p$-th cohomology of the Bloch--Suslin higher Chow complex
$z^q(X, 2q-p)_{\Qbb}$, with regulator
$r_{\cM \to \cD}: H^p_{\cM}(X, \Qbb(q)) \to H^p_{\cD}(X, \Rbb(q))$
into Deligne cohomology (Beilinson 1985, Esnault--Viehweg 1988).
The direct regulator home for an $H^{1,2}$ tangent component is not
$H^3_{\cM}(X,\Qbb(2))=\mathrm{CH}^2(X,1)_{\Qbb}$. The natural
Abel--Jacobi surface is codimension-two motivic cohomology in degree
four:
\[
  H^4_{\cM}(X,\Qbb(2))_{\mathrm{hom}}
  =
  \mathrm{CH}^2(X)_{\mathrm{hom},\Qbb},
\]
whose regulator lands in the intermediate Jacobian $J^2(X)$.

\begin{conjecture}[Motivic normal-function lift of the cocycle; \ClaimStatusConjectured]
\label{prop:motivic-lift}
There is a homologically trivial codimension-two cycle
\[
  Z_\alpha(X)\in
  H^4_{\cM}(X,\Qbb(2))_{\mathrm{hom}}
  =
  \mathrm{CH}^2(X)_{\mathrm{hom},\Qbb}
\]
whose Abel--Jacobi normal function has infinitesimal invariant
\[
  \delta\,\mathrm{AJ}(Z_\alpha)=\alpha_X
  \in H^2(X,\Omega_X^1)=H^{1,2}(X),
\]
after the cyclic-HKR datum and primitive Calabi--Yau projection have
been fixed.
\end{conjecture}

\begin{proof}
This is a conjecture, not a proof. For codimension-two cycles the
Deligne exact sequence has intermediate-Jacobian term $J^2(X)$, and
the tangent quotient of $J^2(X)$ contains
$H^3(X,\Cbb)/F^2H^3(X,\Cbb)$, hence the $H^{1,2}$ component after the
Calabi--Yau primitive projection. Green 1989 supplies the model for
infinitesimal invariants of Abel--Jacobi normal functions. Constructing
the cycle $Z_\alpha(X)$ from the cyclic-HKR commutator is the open
motivic bridge.
\end{proof}

\begin{corollary}[Vanishing at motivic normal-function level; \ClaimStatusConjectured]
\label{cor:motivic-vanish}
On $\cU^{\mathrm{adm}}(X)$, the expected normal function
$\mathrm{AJ}(Z_\alpha)$ is locally constant after projection to the
$H^{1,2}$ tangent component. Off the locus, its infinitesimal invariant
is expected to be non-zero exactly when $\alpha_X$ is non-zero.
\end{corollary}

\begin{proof}
This is a conjectural consequence of Theorem~\ref{thm:main} together
with Conjecture~\ref{prop:motivic-lift}. No general regulator
injectivity statement in this scope is used as an unconditional input.
\end{proof}

The former $H^3_{\cM}(X,\Qbb(2))$ lane may still appear as a
secondary transgression or a dual/conjugate construction, but it is
not the direct regulator home of the $H^{1,2}$ Atiyah--Connes class.

\subsection{Kontsevich graph cohomology: wheel representation}
\label{sub:graph-cohomology}
Kontsevich 1997 \S2 constructs the graph complex $\mathrm{Gra}^{\bullet,\bullet}$ of oriented graphs with edge-contraction differential, and its $E_3$-suboperad $\mathrm{Gra}^{E_3}$ spanned by graphs embeddable in $\Rbb^3$. Kontsevich--Tamarkin 2000 establish a quasi-isomorphism $\mathrm{Gra}^{E_3} \simeq E_3$ in chain complexes over $\Qbb$; the Kontsevich integral $I_{\mathrm{Kont}}: \mathrm{Gra}^{E_3} \to \mathrm{Hom}(T_{\mathrm{poly}}, D_{\mathrm{poly}})$ realises the $L_\infty$-formality map as a sum of Feynman diagrams.

\begin{proposition}[Wheel-graph representation; \ClaimStatusConjectured]
\label{prop:wheel-graph}
The universal cocycle $[m_3, B^{(2)}]^{\mathrm{univ}}$ in $\mathrm{Gra}^{E_3}$ is the linear combination
\[
  [m_3, B^{(2)}]^{\mathrm{univ}} \;=\; c_{\mathrm{tri}}(F) \cdot \Gamma_3^{\mathrm{tri}} \;+\; c_{\mathrm{wheel}}(F) \cdot \Gamma_3^{\mathrm{wheel}},
\]
with $\Gamma_3^{\mathrm{tri}}$ the triangle graph (three internal
vertices each carrying $\at_X$, three edges each carrying the
$\bar\partial$-propagator, three external vertices for the
$B^{(2)}$-contraction), $\Gamma_3^{\mathrm{wheel}}$ the three-wheel
graph (three internal vertices, central cyclic $\Sigma_3$-action), and
coefficients depending on the chosen stable formality datum $F$,
orientation convention, propagator, and HKR/Todd normalisation. The
image under $I_{\mathrm{Kont}}$ at CY-$3$ target $X$ is expected to
coincide with the Dolbeault representative of
Proposition~\ref{prop:cocycle-rep}.
\end{proposition}

\begin{proof}
The Kapranov bracket $\ell_3$ on $T_X[-1]$ is represented in the graph complex by $\Gamma_3^{\mathrm{tri}}$. The cyclic shift $B^{(2)}$ inserts a wheel vertex at the triangle centre, giving $\Gamma_3^{\mathrm{wheel}}$. The numerical coefficients require a single convention for the graph complex orientation, stable formality morphism, cyclic operator transport, propagator, and corrected HKR/Todd map. The paper has not yet performed that convention-matching computation, so the coefficients are kept as $c_{\mathrm{tri}}(F)$ and $c_{\mathrm{wheel}}(F)$.
\end{proof}

\begin{remark}[Graph weights are motivic periods]
\label{rem:motivic-weights}
The coefficients $c_{\mathrm{tri}}(F)$ and $c_{\mathrm{wheel}}(F)$ are
expected to be periods of the graph configuration spaces. Matching
them with a Tate twist or with a Duflo coefficient is a separate
normalisation theorem, not a consequence of their shape alone.
\end{remark}

\subsection{Kontsevich--Tamarkin $E_3$-formality as the Atiyah-class obstruction}
\label{sub:KT-formality}
Kontsevich--Tamarkin $E_d$-formality (Tamarkin 1998, Kontsevich 2003 \S8, Van den Bergh 2007) asserts formality of the $E_d$-operad over $\Qbb$: a quasi-isomorphism $E_d^{\mathrm{chain}} \xrightarrow{\sim} H_\bullet(E_d, \Qbb) = \mathrm{Poi}^{[d]}$ to the shifted-Poisson operad. On a CY-$d$ target $X$, Tamarkin constructs an $L_\infty$-formality morphism $U^d_X : T_{\mathrm{poly}}(X) \xrightarrow{\sim} D_{\mathrm{poly}}(X)$; its second Taylor coefficient $U^d_{2,X}$ obstructs strictness.

\begin{theorem}[Atiyah class as $E_3$-formality obstruction; \ClaimStatusConditional]
\label{thm:formality-atiyah}
On a CY-$3$ $X$ satisfying the Hodge hypothesis, the second Taylor coefficient of the Kontsevich--Tamarkin formality morphism coincides with the cocycle:
\[
  U^3_{2,X} \;\equiv\; [m_3, B^{(2)}]_X \qquad \text{in} \qquad H^2(X, \Omega^1_X) \;\subset\; \HH^3(X, T_X).
\]
On $\cU^{\mathrm{adm}}$, $U^3_{2,X}=0$ kills only this first
formality obstruction. Strict $E_3$-formality requires the vanishing of
the whole higher obstruction tower.
\end{theorem}

\begin{proof}
Kontsevich 2003 Theorem~4.6 gives formality at $d = 2$; the CY-$d$ extension is Tamarkin 1998 and Calaque--Van den Bergh 2010. At bidegree $(3, 1)$ in $\HH^\bullet$, Caldararu--Willerton 2010 Theorem~1.6 identifies $U^3_{2,X}$ with the image of the Kapranov bracket $\ell_3 = \at_X^{[2]}$ (Kapranov 1999 Prop.~4.4) under cyclic-symmetrisation (Tamarkin 1998 \S4). The cyclic symmetrisation is $B^{(2)}$ acting on cyclic-$\Sigma_3$-orbits; $\ell_3 \circ B^{(2)} = [m_3, B^{(2)}]_X$ by Proposition~\ref{prop:cocycle-rep}. Vanishing on $\cU^{\mathrm{adm}}$: Theorem~\ref{thm:main}(i).
\end{proof}

\begin{corollary}[Strictness as a higher-obstruction problem; \ClaimStatusConjectured]
\label{cor:strict-formality}
On $\cU^{\mathrm{adm}}$, strictification of $U^3_X$ to an isomorphism
of $E_3$-algebras is equivalent to the vanishing of the higher
obstruction classes following $U^3_{2,X}$. The vanishing of the second
Taylor coefficient alone is not enough.
\end{corollary}

\begin{proof}
This is the standard obstruction-tower formulation for strictifying a
homotopy-coherent algebra map. Theorem~\ref{thm:formality-atiyah}
identifies the first non-trivial obstruction; it does not prove that
all later obstruction classes are generated by it or vanish.
\end{proof}

\begin{remark}[Scope]
Off $\cU^{\mathrm{adm}}$, the higher Massey obstructions cascade; the formality morphism is genuinely $L_\infty$. The two regimes correspond, under Theorem~\ref{thm:main}, to the two dichotomy classes witnessed by $K3 \times E$ (strict) and $Q_5$ (genuinely $L_\infty$).
\end{remark}

\subsection{Tsygan--Shoikhet cyclic formality and the motivic lift of $B^{(2)}$}
\label{sub:tsygan-formality}
Tsygan 1999 conjectured, and Shoikhet 2003 proved, an $L_\infty$-formality morphism for the cyclic-Hochschild pair
\[
  V^d_X \;:\; (T_{\mathrm{poly}}(X), \Omega^\bullet(X)) \;\xrightarrow{\sim}\; (D_{\mathrm{poly}}(X), \mathrm{CC}^{\mathrm{per}}_\bullet(\cO_X)),
\]
compatible with Connes' cyclic operator $B$: the Shoikhet construction extends Kontsevich's graph-integral formula from the Hochschild to the cyclic setting via weights integrating over $\mathrm{Conf}_n(\Rbb^d) \times S^1$.

\begin{conjecture}[Connes rotation in the motivic normal-function bridge; \ClaimStatusConjectured]
\label{thm:connes-motivic}
The transported cochain Connes operator $\Delta_2$ should enter the
normal-function class $Z_\alpha(X)$ of Conjecture~\ref{prop:motivic-lift}
as the cyclic part of its infinitesimal invariant. This is a
compatibility conjecture between Tsygan--Shoikhet cyclic formality,
the cyclic-HKR datum, and the Abel--Jacobi normal function; it is not a
standalone claim that $B^{(2)}$ defines a class in
$\mathrm{Pic}(X)_{\Qbb}$.
\end{conjecture}

\begin{proof}
Shoikhet 2003 constructs a cyclic formality morphism by graph integrals
over configuration spaces with an $S^1$-parameter. This supplies the
right technology for comparing the cochain Connes rotation with an
infinitesimal normal function, but the paper has not constructed the
codimension-two cycle or proved that the cyclic operator has an
independent Picard-valued motivic lift.
\end{proof}

\begin{corollary}[No explicit motivic representative yet; \ClaimStatusOpen]
\label{cor:motivic-rep}
No higher-Chow representative for $Z_\alpha(X)$ is currently supplied
by the manuscript. In particular, the torsion graph of a Chern
connection is not asserted to be an algebraic cycle, and no
$\mathrm{CH}^3(X,1)$ representative is used as an input.
\end{corollary}

\begin{proof}
This records an open proof obligation following
Conjecture~\ref{prop:motivic-lift}.
\end{proof}

\subsection{HKR--Duflo universal obstruction (Caldararu--Willerton 2010 Theorem~1.6)}
\label{sub:HKR-duflo}
The universal nature of the cocycle --- its arising not from the particular $X$ but from the moduli of $A_\infty$-algebras --- is expressed via the HKR--Duflo isomorphism.

\begin{theorem}[HKR--Duflo universal cocycle, Caldararu--Willerton 2010 Theorem~1.6; \ClaimStatusProvedElsewhere]
\label{thm:hkr-duflo}
For $X$ smooth projective with $T_X$, the HKR isomorphism $I_{\mathrm{HKR}}: \HH^\bullet(X) = \bigoplus_{p+q=\bullet} H^p(X, \wedge^q T_X) \to \Hom^\bullet(\cO_X, \cO_X)[\bullet]$ fails to intertwine Yoneda composition with wedge product. The obstruction is the Duflo automorphism
\[
  \mathrm{Duf}_X \;=\; \sum_{n \geq 2} \frac{B_{2n}}{(2n)!} \at_X^{[2n]} \;\in\; \prod_{n \geq 2} H^{2n}(X, \Omega^{2n}_X),
\]
with $B_{2n}$ Bernoulli numbers. The twisted HKR map $I_{\mathrm{HKR}} \circ \mathrm{Duf}_X$ is an isomorphism of Gerstenhaber algebras.
\end{theorem}

\begin{proposition}[Cocycle is the second Duflo coefficient; \ClaimStatusConditional]
\label{prop:cocycle-duflo}
The cocycle $[m_3, B^{(2)}]_X$ is expected to be the image under the HKR--Duflo isomorphism of the second Duflo coefficient, with a convention-dependent scalar $c_{\mathrm{Duf}}$:
\[
  [m_3, B^{(2)}]_X \;=\; c_{\mathrm{Duf}}\,
  I_{\mathrm{HKR}}\bigl(\at_X^{[2]}\bigr)_{(2,1)\text{-piece}} .
\]
The identification requires matching the cyclic operator transport,
HKR/Todd normalisation, and graph-formality conventions. The numerical
value $1/12$ is not used here as a theorem-grade coefficient.
\end{proposition}

\begin{proof}
Caldararu--Willerton 2010 Theorem~1.6 expresses the Duflo correction
in terms of Atiyah powers. Restriction to the relevant CY-$3$ Hodge
piece gives the correct kind of target, but the paper has not matched
the Bernoulli coefficient with the cyclic-HKR and graph conventions.
This proposition is therefore a conditional comparison, not a computed
constant.
\end{proof}

\begin{corollary}[Universality; \ClaimStatusConjectured]
\label{cor:universal}
The class $[m_3, B^{(2)}]_X$ is the image at $X$ of a universal class $[m_3, B^{(2)}]^{\mathrm{univ}} \in \HH^3(\mathrm{BV}^{E_3}, T_{\mathrm{univ}})$ on the moduli of $E_3$-algebras with BV structure, under the classifying map $X \to \mathrm{BV}^{E_3}$. The universal class is the second Duflo coefficient on the universal $E_3$-algebra.
\end{corollary}

\begin{proof}
The second Duflo coefficient depends only on the $E_3$-operad (Kontsevich 2003 \S8, Calaque--Van den Bergh 2010), not on $X$; its specialisation at $X$ is Proposition~\ref{prop:cocycle-duflo}.
\end{proof}

\subsection{Configuration-space rational homotopy, not Malcev $\pi_1$}
\label{sub:malcev-quillen}
Quillen 1969 identifies the rational homotopy type of a simply
connected space with a differential graded Lie model. For a simply
connected smooth complex threefold $X$, this is the relevant
technology for $\Conf_3(X)$: the diagonals in $X^3$ have complex
codimension $3$, hence real codimension $6$, so loops and null-homotopy
discs avoid them and $\pi_1(\Conf_3(X))=0$. The Malcev completion of
$\pi_1$ is therefore trivial in the main CY3 lane. The remaining
frontier is a Kriz--Quillen rational-homotopy comparison, not a
Malcev $\pi_1$ theorem.

\begin{conjecture}[Kriz--Quillen configuration-space comparison; \ClaimStatusConjectured]
\label{thm:malcev-quillen}
Let $K_X(3)$ denote the Kriz--Totaro--Bezrukavnikov CDGA model of
$\Conf_3(X)$, whose diagonal generators have degree
$2\dim_{\Cbb}X-1=5$. There should be a filtered morphism
\[
  \Xi_X:\operatorname{Def}^{E_3,\mathrm{cyc}}
  (\mathrm{Perf}_{\mathrm{dg}}(X))
  \longrightarrow \operatorname{Der}(K_X(3))
\]
sending the Atiyah--Connes class to a weight-three derivation class of
the Kriz model. This conjectural comparison is the configuration-space
face of the obstruction.
\end{conjecture}

\begin{proof}
This is a conjectural replacement for the false Malcev $\pi_1$ claim.
Kriz and Totaro provide the diagonal CDGA model, and Quillen provides
the dg Lie rational-homotopy model. The missing step is the explicit
map $\Xi_X$ from the cyclic deformation complex to derivations of
$K_X(3)$ and the calculation of its weight-three component.
\end{proof}

\begin{proposition}[Malcev $\pi_1$ obstruction vanishes trivially on simply connected CY3s; \ClaimStatusProvedHere]
\label{cor:pi1-motivic}
If $X$ is simply connected and $\dim_{\Cbb}X=3$, then
\[
  \pi_1(\Conf_3(X))=0.
\]
Consequently the Malcev completion of $\pi_1(\Conf_3(X))$ has no
non-trivial weight-three central extension to compare with
$\alpha_X$.
\end{proposition}

\begin{proof}
The ambient space $X^3$ is simply connected. The pairwise diagonals
have complex codimension $3$, hence real codimension $6$. By
transversality any loop in $X^3$ and any null-homotopy disc can be
homotoped rel boundary away from the union of the diagonals. Therefore
removing the diagonals does not change $\pi_1$.
\end{proof}

\subsection{Six faces of one obstruction}
\label{sub:five-faces}
The class admits two proved faces --- the Dolbeault/HKR construction
(\S\ref{sec:cocycle}) and the stage-two specialisation onto the first
bar--cobar obstruction (\S\ref{sec:spine}) --- and four frontier
comparison faces. The frontier faces name the precise bridges whose
proofs would close the Atiyah--Connes pentagon.

\begin{center}
\begin{tabular}{lll}
\toprule
 & Face & Realisation \\
\midrule
0 & Dolbeault/HKR, Prop.~\ref{prop:cocycle-rep} &
  $\iota_{\Omega_X}\tr_{\End T_X}(\bar\partial\nabla\wedge\bar\partial\nabla)$
  \emph{(proved)} \\
0$'$ & Spine, Thm.~\ref{thm:sp-alpha-first-barcobar-obstruction} &
  $\operatorname{sp}_{\Sigma,C}(\alpha_X)=o^{(0)}_3(A_C)$
  \emph{(proved)} \\
1 & Motivic, Conj.~\ref{prop:motivic-lift} & $\delta\,\mathrm{AJ}(Z_\alpha)=\alpha_X$ \\
2 & Graph, Prop.~\ref{prop:wheel-graph} & $c_{\mathrm{tri}}\Gamma_3^{\mathrm{tri}} + c_{\mathrm{wheel}}\Gamma_3^{\mathrm{wheel}}$ \\
3 & HKR--Duflo, Prop.~\ref{prop:cocycle-duflo} & $c_{\mathrm{Duf}}\,\at_X^{[2]}$ in fixed convention \\
4 & Kriz--Quillen, Conj.~\ref{thm:malcev-quillen} & weight-three derivation of $K_X(3)$ \\
\bottomrule
\end{tabular}
\end{center}

Conditional face equivalences: $0\equiv 1$ requires a codimension-two
cycle whose Abel--Jacobi infinitesimal invariant is $\alpha_X$;
$1\equiv 2$ requires the relevant Kontsevich graph-integral computation
with matching normalisation; $2\equiv 3$ requires wheel-weight matching
in the chosen HKR--Duflo convention; $3\equiv 4$ requires a
Kriz--Quillen rational-homotopy comparison in the same weight.
Theorem~\ref{thm:main} uses faces $1$--$4$ as comparison hypotheses;
faces $0$ and $0'$ are the paper's load-bearing proved content.

\subsection{Chain-level / $(\infty,1)$-categorical compatibility}
The $(\infty,1)$-categorical avatar of the construction is the composition of two natural transformations of $\infty$-functors $D^b(X) \to D^b(X)$: the Massey bracket as a coherent associator homotopy, and the Connes cyclic shift as a $\Zbb/3$-equivariant rotation on cyclic chains. The output is a natural transformation $\id \Rightarrow (-) \otimes \Omega^1_X[2]$, classified by an element of
\[
  \Hom_{\mathrm{Fun}(D^b X, D^b X)}(\id, \otimes \Omega^1_X[2]) \;\cong\; \Ext^2_X(\cO_X, \Omega^1_X) \;\cong\; H^2(X, \Omega^1_X).
\]
Proposition~\ref{prop:cocycle-rep} exhibits the chain-level representative. The bridge between the two lanes is Caldararu--Willerton 2010 Theorem~1.6, identifying Hochschild cohomology classes with natural transformations of the identity functor.

% ==================================================
\section{Curve-side target and the bar--cobar spine theorem}
\label{sec:spine}
% ==================================================

The Atiyah--Connes class $\alpha_X\in H^2(X,\Omega^1_X)$ lives in
coherent sheaf cohomology on the CY-$3$. The curve-side target is the
genus-zero arity-three obstruction of the chiral algebra obtained by
factorisation homology over a two-cycle $\Sigma\subset X$ followed by
restriction to the reference curve $C\subset X$. The two sides are
separated by a genuine functor, not by an ordinary coherent pullback:
the pullback has no room for $\alpha_X$.

\begin{lemma}[Curve target vanishing; \ClaimStatusProvedHere]
\label{lem:curve-h2-omega-vanishes}
For every smooth algebraic curve $C$ over $\Cbb$,
\[
  H^2(C,\Omega_C^1)=0.
\]
Consequently the ordinary coherent restriction morphism
\[
  H^2(X,\Omega_X^1)\to H^2(C,\Omega_C^1)
\]
is the zero map, and cannot detect or kill $\alpha_X$.
\end{lemma}

\begin{proof}
A smooth algebraic curve has complex dimension one, so for any
coherent sheaf $\cF$ the Grothendieck vanishing theorem gives
$H^i(C,\cF)=0$ for $i\geq 2$ (Hartshorne 1977 Theorem~III.2.7). In
particular $H^2(C,\Omega_C^1)=0$.
\end{proof}

\subsection{The curve-side chiral algebra and its modular cyclic deformation complex}

Fix the two-stage factorisation data: $\Sigma\subset X$ a smooth real
two-cycle with tubular neighbourhood retracting onto a smooth
projective algebraic curve $C\subset X$. Vol.~III's stage-two
specialisation functor is
\[
  \Sp^{\mathrm{ch}}_{\Sigma,C}
  =\operatorname{res}_C\circ\int_\Sigma
  :\operatorname{Fact}^{E_3,\mathrm{hol}}_X
  \to\ChirAlg^{\mathrm{ch}}_C,
\]
the composition of factorisation homology over $\Sigma$
(Ayala--Francis 2015 Definition~3.1) and restriction to $C$. Applied
to the stage-one output $\Phi^{\mathrm{FA}}_3(\Perf_{\mathrm{dg}}(X))$
of \S\ref{sub:fa-precision}, it produces the \emph{curve-side chiral
algebra}
\[
  A_C
  :=\Sp^{\mathrm{ch}}_{\Sigma,C}
  \Phi^{\mathrm{FA}}_3(\Perf_{\mathrm{dg}}(X))
  \in\ChirAlg^{\mathrm{ch}}_C.
\]

The ordered chiral bar complex $\Bar^{\mathrm{ch}}_C(A_C)
=T^c_C(s^{-1}\bar A_C)$ on $\Conf_n(C)$ carries an augmented Maurer--Cartan
element
\[
  \Theta_{A_C}
  =\sum_{n\geq 2}\Theta_{A_C}^{(n)},
  \qquad
  \Theta_{A_C}^{(n)}\in C_n^{\mathrm{ord}}(A_C)
\]
solving
$D\cdot\Theta_{A_C}+\tfrac{1}{2}[\Theta_{A_C},\Theta_{A_C}]=0$.
Completing and averaging across genera produces the
\emph{modular cyclic deformation complex}
\[
  \Defcyc^{\mathrm{mod}}(A_C)
  :=\operatorname{CoDer}^{\mathrm{cyc}}
  \!\bigl(\widehat{\bar B}^{\mathrm{ch}}_C(A_C)\bigr)[1]
\]
of cyclic coderivations of the completed symmetric chiral bar
coalgebra, shifted by one.

\begin{definition}[First bar--cobar obstruction]
\label{def:first-barcobar-obstruction}
For $r\geq 2$ and $g\geq 0$, let
$\pi_{r,g}:\Defcyc^{\mathrm{mod}}(A_C)\to\Defcyc^{\mathrm{mod}}(A_C)$
denote the projection onto the bar-length $r$ and genus $g$ component.
The \emph{first bar--cobar obstruction} of $A_C$ is the class of
$\Theta_{A_C}^{(3)}$ in arity $3$, genus $0$:
\[
  o^{(0)}_3(A_C)
  :=
  \bigl[\pi_{3,0}\Theta_{A_C}^{(3)}\bigr]
  \in H^2\!\bigl(\pi_{3,0}\Defcyc^{\mathrm{mod}}(A_C)\bigr).
\]
The full genus-zero obstruction tower is
$o^{(0)}_{\geq 3}(A_C)
=\bigl(o^{(0)}_3(A_C),o^{(0)}_4(A_C),\ldots\bigr)$.
\end{definition}

\subsection{Differentiability of stage-two specialisation}

Vol.~III establishes $\Sp^{\mathrm{ch}}_{\Sigma,C}$ as a derived
functor in $\Pr^L_{k,\mathrm{st}}$. The proved input below is that
its derivative on formal moduli problems is a filtered
$L_\infty$-morphism preserving the bar-length and genus
filtrations. This is the load-bearing new functoriality statement: it
is stronger than the existence of the stage-two functor on objects.

\begin{theorem}[Ayala--Francis differential of stage-two specialisation; \ClaimStatusProvedHere]
\label{thm:sp-differentiable}
Stage-two specialisation $\Sp^{\mathrm{ch}}_{\Sigma,C}$ is
differentiable as a functor of formal moduli problems (Lurie DAG-X
Theorem~5.2.3.11). Its derivative is a filtered $L_\infty$-morphism
\[
  d\Sp^{\mathrm{ch}}_{\Sigma,C}:
  \Def^{E_3,\mathrm{cyc}}
  \!\bigl(\Phi^{\mathrm{FA}}_3(\Perf_{\mathrm{dg}}(X))\bigr)
  \longrightarrow
  \Defcyc^{\mathrm{mod}}(A_C)
\]
preserving the bar-length filtration by arity and the genus filtration
by Deligne--Mumford boundary. On HKR graded pieces, at bi-arity-genus
$(r,g)=(3,0)$, the ternary Stasheff operation and the cochain Connes
rotation descend to the arity-three cyclic coderivation of
$\widehat{\bar B}^{\mathrm{ch}}_C(A_C)$.
\end{theorem}

\begin{proof}
The $E_3$-factorisation algebra $\Phi^{\mathrm{FA}}_3(\Perf_{\mathrm{dg}}(X))$
is pinned by Kontsevich--Tamarkin $E_3$-formality and Costello--Gwilliam--Li
locality (\S\ref{sub:fa-precision}). Ayala--Francis 2015 Definition~3.1
defines $\int_\Sigma\cF$ as the homotopy colimit of $\cF$ over the
Ran space $\operatorname{Ran}(\Sigma)$ of finite subsets; it is a
symmetric monoidal functor preserving sifted colimits
(Ayala--Francis 2015 Theorem~3.16, Fubini). Restriction
$\operatorname{res}_C$ is the pullback along the closed immersion
$C\hookrightarrow X$ of the resulting factorisation algebra on the
normal neighbourhood of $\Sigma$; it is the coherent restriction functor
on factorisation algebras (Ayala--Francis--Tanaka 2017 Theorem~4.3.1).

Both legs preserve compact objects and colimits in
$\Pr^{L}_{k,\mathrm{st}}$, hence they send formal moduli problems to
formal moduli problems and their derivatives are filtered $L_\infty$-
morphisms (Lurie HA 7.3.2; Costello--Gwilliam FA Vol.~2, Theorem~2.2.1).
The cyclic $E_3$-deformation complex on the source is
$\operatorname{CC}^{-,\bullet+3}_{\mathrm{cyc}}$ of the stage-one
factorisation algebra, and the HKR decomposition identifies it with
$\bigoplus_{p+q=\bullet+3}H^p(X,\wedge^qT_X)$ in cohomology (Caldararu
2005 Theorem~1.2). On the target side, the modular cyclic deformation
complex of $A_C$ is the cyclic coderivation complex of the completed
bar coalgebra (Kapranov--Soibelman 2009 \S4.7).

Filtration preservation by bar length: the stage-two functor sends
factorisation-algebra configuration-space local sections over
$\Conf_n(\Sigma)$ to bar-complex sections over $\Conf_n(C)$, and the
cyclic structure is preserved by $\int_\Sigma$ because Ayala--Francis
factorisation homology is symmetric monoidal over the cyclic operad
(Ayala--Francis 2015 \S5, Kaledin 2018). Genus filtration preservation:
$\Sp^{\mathrm{ch}}_{\Sigma,C}$ respects the Deligne--Mumford
compactification stratification of $\overline{\cM}_{g,n}(C)$ by sending
the stage-one genus grading to the curve-side genus grading; the
restriction $\operatorname{res}_C$ picks out the component attached to
$C$. On the HKR piece at $(r,g)=(3,0)$, the ternary operation $m_3$ is
represented by the arity-three cell of the bar differential, and
$\Delta_2$ by the cyclic rotation of three bar slots; the specialisation
descends these to the arity-three cyclic coderivation of
$\widehat{\bar B}^{\mathrm{ch}}_C(A_C)$.
\end{proof}

\subsection{The spine theorem: specialisation of the Atiyah--Connes class}

Define the specialisation morphism on obstruction cohomology as the
composite
\[
\begin{tikzcd}[column sep=normal]
H^2(X,\Omega_X^1)
  \arrow[r,"\mathrm{HKR}^{-1}_{\mathrm{cyc}}"]
&
H^2\!\bigl(\Def^{E_3,\mathrm{cyc}}(\mathcal F_X)\bigr)
  \arrow[r,"d\Sp^{\mathrm{ch}}_{\Sigma,C}"]
&
H^2\!\bigl(\Defcyc^{\mathrm{mod}}(A_C)\bigr)
  \arrow[r,"\pi_{3,0}"]
&
H^2\!\bigl(\pi_{3,0}\Defcyc^{\mathrm{mod}}(A_C)\bigr)
\end{tikzcd}
\]
where $\cF_X=\Phi^{\mathrm{FA}}_3(\Perf_{\mathrm{dg}}(X))$ and
$\mathrm{HKR}^{-1}_{\mathrm{cyc}}$ is the inverse of the
cyclic-HKR isomorphism
(Caldararu--Willerton 2010 Theorem~1.6) applied to the cyclic
$E_3$-deformation complex. Set
\[
  \operatorname{sp}_{\Sigma,C}
  :=\pi_{3,0}\circ d\Sp^{\mathrm{ch}}_{\Sigma,C}
  \circ\mathrm{HKR}^{-1}_{\mathrm{cyc}}.
\]

\begin{theorem}[Specialisation of the Atiyah--Connes class; \ClaimStatusProvedHere]
\label{thm:sp-alpha-first-barcobar-obstruction}
Let $X$ be a smooth projective Calabi--Yau threefold with cyclic-HKR
datum $\mathfrak c_X$ and Atiyah--Connes class
$\alpha_X\in H^2(X,\Omega_X^1)$
(Definition~\ref{def:atiyah-connes-class},
Theorem~\ref{thm:cyclic-hkr-datum-exists},
Construction~\ref{constr:cocycle}). Let $(\Sigma,C)$ be an admissible
two-stage specialisation datum for $X$
(\S\ref{subsec:two-stage-derived}), and let
$A_C=\Sp^{\mathrm{ch}}_{\Sigma,C}\Phi^{\mathrm{FA}}_3
(\Perf_{\mathrm{dg}}(X))$ with first bar--cobar obstruction
$o^{(0)}_3(A_C)\in H^2\!\bigl(\pi_{3,0}\Defcyc^{\mathrm{mod}}(A_C)\bigr)$
(Definition~\ref{def:first-barcobar-obstruction}). Then
\[
  \operatorname{sp}_{\Sigma,C}(\alpha_X)
  =
  o^{(0)}_3(A_C)
  \in H^2\!\bigl(\pi_{3,0}\Defcyc^{\mathrm{mod}}(A_C)\bigr).
\]
\end{theorem}

\begin{proof}
Under the cyclic HKR isomorphism of Caldararu--Willerton 2010
Theorem~1.6, $\alpha_X=\iota_{\Omega_X}\pi^{\mathrm{HKR}}_{2,2}[m_3,\Delta_2]$
lifts uniquely to the $E_3$-cyclic deformation class
\[
  \mathrm{HKR}^{-1}_{\mathrm{cyc}}(\alpha_X)
  =[m_3,\Delta_2]_{\mathrm{cyc}}
  \in H^2\!\bigl(\Def^{E_3,\mathrm{cyc}}(\mathcal F_X)\bigr),
\]
because the cyclic-HKR datum fixes the Calabi--Yau trace convention so
that the projection $\pi^{\mathrm{HKR}}_{2,2}$ followed by
$\iota_{\Omega_X}$ is an isomorphism from the $(3,0)$-arity-genus
summand to $H^2(X,\Omega_X^1)$.

By Theorem~\ref{thm:sp-differentiable}, $d\Sp^{\mathrm{ch}}_{\Sigma,C}$
is a filtered $L_\infty$-morphism preserving bar-length and genus
filtrations. Applying $d\Sp^{\mathrm{ch}}_{\Sigma,C}$ to the cyclic
class $[m_3,\Delta_2]_{\mathrm{cyc}}$ and projecting to
$(r,g)=(3,0)$ picks out the arity-three genus-zero component. By the
local form of $d\Sp^{\mathrm{ch}}_{\Sigma,C}$ at this bi-arity
(Theorem~\ref{thm:sp-differentiable}, chain-level identification), the
ternary Stasheff operation $m_3$ on the source descends to the
arity-three coderivation of $\widehat{\bar B}^{\mathrm{ch}}_C(A_C)$,
and the cochain Connes rotation $\Delta_2$ descends to the cyclic
symmetrisation of three bar slots. The result is the arity-three
cyclic coderivation of $\widehat{\bar B}^{\mathrm{ch}}_C(A_C)$, i.e.,
the arity-three component of the curve-side Maurer--Cartan element:
\[
  \pi_{3,0}\circ d\Sp^{\mathrm{ch}}_{\Sigma,C}
  \bigl([m_3,\Delta_2]_{\mathrm{cyc}}\bigr)
  =\bigl[\pi_{3,0}\Theta_{A_C}^{(3)}\bigr]
  =o^{(0)}_3(A_C)
  \in H^2\!\bigl(\pi_{3,0}\Defcyc^{\mathrm{mod}}(A_C)\bigr),
\]
by Definition~\ref{def:first-barcobar-obstruction}. Composing with
$\mathrm{HKR}^{-1}_{\mathrm{cyc}}$ gives
$\operatorname{sp}_{\Sigma,C}(\alpha_X)=o^{(0)}_3(A_C)$.
\end{proof}

\begin{corollary}[Atiyah--Connes vanishing kills the first bar--cobar obstruction; \ClaimStatusProvedHere]
\label{cor:alpha-vanish-kills-o3}
If $\alpha_X=0$ in $H^2(X,\Omega_X^1)$, then $o^{(0)}_3(A_C)=0$.
\end{corollary}

\begin{proof}
Apply $\operatorname{sp}_{\Sigma,C}$ to both sides using
Theorem~\ref{thm:sp-alpha-first-barcobar-obstruction}: the left side
is $\operatorname{sp}_{\Sigma,C}(0)=0$, so
$o^{(0)}_3(A_C)=0$.
\end{proof}

\begin{corollary}[Positselski inversion after higher-obstruction closure; \ClaimStatusConditional]
\label{cor:sp-alpha-positselski}
Assume $\alpha_X=0$ and one of the following:
\begin{enumerate}[label=\textup{(K\arabic*)}]
  \item $A_C\in\operatorname{Kosz}(C)$, equivalently
  $o^{(0)}_{r+1}(A_C)=0$ for all $r\geq 2$;
  \item $A_C$ lies in the completed filtered pro-conilpotent
  Positselski surface of Vol.~I with finite graded bar pieces;
  \item every higher genus-zero obstruction $o^{(0)}_{r+1}(A_C)$ for
  $r\geq 3$ lies in the ideal generated by $o^{(0)}_3(A_C)$ in the
  minimal $L_\infty$-model of $\pi_{\star,0}\Defcyc^{\mathrm{mod}}(A_C)$.
\end{enumerate}
Then the Vol.~I bar--cobar counit
\[
  \Cobar_C^{\mathrm{ch}}\Bar_C^{\mathrm{ch}}(A_C)\to A_C
\]
is a quasi-isomorphism: strict in case \textup{(K1)}, coderived in
case \textup{(K2)}, and strict in case \textup{(K3)} via
Corollary~\ref{cor:alpha-vanish-kills-o3}.
\end{corollary}

\begin{proof}
Case \textup{(K1)} is Vol.~I Theorem~B (strict Positselski inversion
on the Koszul locus). Case \textup{(K2)} is Vol.~I Theorem~B
(completed Positselski inversion under the pro-conilpotent
finite-piece hypotheses). Case \textup{(K3)}: by
Corollary~\ref{cor:alpha-vanish-kills-o3}, $o^{(0)}_3(A_C)=0$; the
generated-by-$o_3$ hypothesis then forces $o^{(0)}_{r+1}(A_C)=0$ for
all $r\geq 3$, reducing to case \textup{(K1)}.
\end{proof}

% ==================================================
\section{The Atiyah--Connes bridge theorem}
\label{sec:main}
% ==================================================

\begin{definition}[Atiyah--Connes derived zero locus]
\label{def:atiyah-connes-zero-locus}
Let
\[
  \mathcal B=
  \mathcal M^{\Omega,\mathrm{cyc},\Sigma,C,L,\mathrm{tr}}_{\mathrm{CY}_3}
\]
denote the derived analytic moduli stack of Calabi--Yau threefolds
equipped with the auxiliary data needed for the cyclic-HKR datum,
stage-two specialisation, Mukai lattice, and trace normalisation. For
the universal family $\pi:\mathcal X\to\mathcal B$ set
\[
  \mathbb E_{\mathrm{AC}}=R\pi_*\Omega^1_{\mathcal X/\mathcal B}[2].
\]
A universal Atiyah--Connes class, when constructed, is a section
\[
  \alpha_{\mathrm{AC}}\in \Gamma(\mathcal B,\mathbb E_{\mathrm{AC}}).
\]
Its derived zero locus is
\[
  \mathbf Z_\alpha=\mathbf Z(\alpha_{\mathrm{AC}}).
\]
The Positselski, factorisation, K3/Igusa, and trace vertices below are
comparison targets for this zero locus only after maps from
$\mathbb E_{\mathrm{AC}}$ to their obstruction complexes have been
constructed.
\end{definition}

\begin{theorem}[Atiyah--Connes comparison bridge; \ClaimStatusConditional]
\label{thm:main}
Let $X$ be a smooth projective Calabi--Yau threefold with cyclic-HKR
datum $\mathfrak c_X$ and Atiyah--Connes class
$\alpha_X\in H^2(X,\Omega^1_X)$
(Theorem~\ref{thm:cyclic-hkr-datum-exists},
Construction~\ref{constr:cocycle}), and let $\mathbf Z_\alpha$ be the
derived zero locus of the universal class
(Definition~\ref{def:atiyah-connes-zero-locus}). Over a common derived
base $\cB^\circ$ let $\mathbf Z_{\mathrm{BC}}$,
$\mathbf Z_{\mathrm{FA}}$, $\mathbf Z_{\Delta_5}$,
$\mathbf Z_{\mathrm{tr}}$ be the bar--cobar, holomorphic factorisation,
K3/Mukai/Igusa, and trace witness loci of \S\ref{sec:climaxes}. The
unconditional part of the theorem --- the spine identity
$\operatorname{sp}_{\Sigma,C}(\alpha_X)=o^{(0)}_3(A_C)$ and its
corollary $\alpha_X=0\Rightarrow o^{(0)}_3(A_C)=0$ --- is supplied by
Theorem~\ref{thm:sp-alpha-first-barcobar-obstruction} and
Corollary~\ref{cor:alpha-vanish-kills-o3}; no ordinary restriction to
$H^2(C,\Omega^1_C)$ is used
(Lemma~\ref{lem:curve-h2-omega-vanishes}). The conditional assertions
follow below.
\begin{enumerate}[label=\textup{(\roman*)}]
  \item \textup{(Hodge target; proved)} Under HKR and the Calabi--Yau
  contraction of the cyclic-HKR datum, the class lies in
  \[
    H^2(X,\Omega^1_X)=H^{1,2}(X).
  \]
  If it is written as an $H^{2,1}$-class, this means the Calabi--Yau-dual or conjugate class.

  \item \textup{(Positselski bridge)} Assume \textup{(B1)} the tower
  hypothesis \textup{(K3)} of
  Corollary~\ref{cor:sp-alpha-positselski} holds: every higher
  obstruction $o^{(0)}_{r+1}(A_C)$ for $r\geq 3$ lies in the ideal
  generated by $o^{(0)}_3(A_C)$. Then the bar--cobar zero locus
  coincides with $\mathbf Z_\alpha$ and the counit
  $\Cobar_C^{\mathrm{ch}}\Bar_C^{\mathrm{ch}}(A_C)\to A_C$ is a strict
  quasi-isomorphism on $\mathbf Z_\alpha$.

  \item \textup{(factorisation bridge)} Assume \textup{(B2)} the
  Costello--Li parametrix obstruction for $\Phi^{\mathrm{FA}}_3(X)$ is
  identified with $\alpha_X$ in the holomorphic factorisation ambient
  of \S\ref{sub:fa-precision}, with no independent higher-loop
  locality obstruction. Then
  $\mathbf Z_\alpha\simeq\mathbf Z_{\mathrm{FA}}$ on the locality-safe
  substack.

  \item \textup{(principal K3 locus)} Assume \textup{(B3)} the
  stage-two datum $(\Sigma,C)$ lands on the principal $K3\times E$
  Mukai-admissible locus of Vol.~III, with the BL-2/4/6 comparison
  witnesses and the BL-7 Borcherds-weight normalisation. Then
  $\mathbf Z_\alpha\simeq\mathbf Z_{\Delta_5}$ on that locus.

  \item \textup{(trace normalisation)} Assume \textup{(B4)} the
  derived-centre trace is normalised by the Mukai conductor identity
  $K^{\kappa_{\mathrm{ch}}}_{\mathrm{Muk}}
  =2c_+(\widetilde H(K3,\Zbb))=8$ and the trace period morphism is
  conservative on the cyclic-skew obstruction summand. Then
  $\mathbf Z_\alpha\simeq\mathbf Z_{\mathrm{tr}}$ and the scalar
  identity $\hbar^2K^{\kappa_{\mathrm{ch}}}_{\mathrm{Muk}}=-1$ holds on
  the same locus.
\end{enumerate}
Under \textup{(B1)--(B4)} the four witnesses assemble into the
pentagon
\[
  \mathbf Z_\alpha
  \simeq\mathbf Z_{\mathrm{BC}}
  \simeq\mathbf Z_{\mathrm{FA}}
  \simeq\mathbf Z_{\Delta_5}
  \simeq\mathbf Z_{\mathrm{tr}},
\]
a fiber-product equivalence of derived zero loci over $\cB^\circ$.
Applying $\QCoh$ or $\IndCoh$ then gives equivalences in
$\mathrm{Pr}^{\mathrm L}_{k,\mathrm{st}}$
(Corollary~\ref{cor:prl-shadow-anomaly-loci}). Without
\textup{(B1)--(B4)} the pentagon is conjectural
(Conjecture~\ref{conj:atiyah-connes-pentagon}).
\end{theorem}

\begin{proof}
Part (i) is Corollary~\ref{cor:hodge}. Part (ii) applies
Corollary~\ref{cor:sp-alpha-positselski} case \textup{(K3)}: the spine
theorem supplies
$\operatorname{sp}_{\Sigma,C}(\alpha_X)=o^{(0)}_3(A_C)$
unconditionally, and the tower hypothesis closes higher obstructions
to zero on $\mathbf Z_\alpha$. Part (iii) uses \textup{(B2)} to
identify $\mathbb E_{\mathrm{AC}}$ and the Costello--Li parametrix
complex as quasi-isomorphic on the locality-safe substack, hence
equating their derived zero loci. Part (iv) uses \textup{(B3)}: the
BL-comparison data of Theorem~\ref{thm:climax-III} is assumed to
provide a quasi-isomorphism of the Atiyah--Connes obstruction complex
with the $\mathbf H_{\Delta_5}$ witness complex on the principal locus. Part
(v) uses \textup{(B4)}: the Mukai conductor trace normalisation
identifies the trace-witness complex with the cyclic-skew summand of
$\mathbb E_{\mathrm{AC}}$. The four equivalences assemble to the
pentagon by derived fiber-product equivalence over $\cB^\circ$. No
ordinary restriction to $H^2(C,\Omega^1_C)$ is used throughout, by
Lemma~\ref{lem:curve-h2-omega-vanishes}.
\end{proof}

\begin{corollary}[Presentable-category shadow; \ClaimStatusConditional]
\label{cor:prl-shadow-anomaly-loci}
Under the hypotheses of Theorem~\ref{thm:main} and assuming the four
witness loci are perfect derived stacks,
\[
  \QCoh(\mathbf Z_\alpha)
  \simeq\QCoh(\mathbf Z_{\mathrm{BC}})
  \simeq\QCoh(\mathbf Z_{\mathrm{FA}})
  \simeq\QCoh(\mathbf Z_{\Delta_5})
  \simeq\QCoh(\mathbf Z_{\mathrm{tr}})
\]
in $\mathrm{Pr}^{\mathrm L}_{k,\mathrm{st}}$. The analogous
equivalence holds for $\IndCoh$ under the standard quasi-smooth
locally almost-of-finite-type hypotheses.
\end{corollary}

\begin{proof}
$\QCoh$ is a symmetric monoidal $(\infty,1)$-functor sending
equivalences of derived stacks to equivalences in
$\mathrm{Pr}^{\mathrm L}_{k,\mathrm{st}}$ (Lurie $\mathrm{HA}$
\S5.5.3); apply it to the equivalences of
Theorem~\ref{thm:main}\textup{(ii)--(v)}. The $\IndCoh$ case is
Gaitsgory--Rozenblyum 2017 under the stated hypotheses.
\end{proof}

\begin{conjecture}[Five-vertex comparison diagram; \ClaimStatusConjectured]
\label{conj:atiyah-connes-pentagon}
Under \textup{(B1)--(B4)} and all-loop obstruction closure, the
vanishing of $\alpha_X$ is equivalent to the four comparison statements on
the admissible locus. The conjecture is false with ordinary curve
restriction in place of $\operatorname{sp}_{\Sigma,C}$.
\end{conjecture}


\begin{theorem}[Conditional hCS anomaly bridge; \ClaimStatusConditional]
\label{thm:adler-bardeen}
Assume the holomorphic BV theory on the Calabi--Yau threefold $X$ is
quantised with a Costello--Li parametrix, that the local BRST anomaly
class of holomorphic Chern--Simons theory is identified with the same
cyclic HKR representative used in Construction~\ref{constr:cocycle},
and that the one-loop exactness statement is proved in this holomorphic
six-real-dimensional setting. Under these hypotheses, the class
$[m_3,B^{(2)}]_X\in H^2(X,\Omega^1_X)$ is the hCS anomaly-cancellation
class. Explicitly:
\begin{itemize}
  \item The holomorphic Chern--Simons action $S_{\mathrm{hCS}}[\cA] = \int_X \mathrm{CS}_3(\cA) \wedge \Omega_X$ on a smooth connection $\cA$ valued in $\End(T_X)$ (with $\Omega_X$ the holomorphic volume form) is BRST-invariant up to total derivatives; the BRST variation of $\mathrm{CS}_3$ is $\delta \mathrm{CS}_3 = d \mathrm{CS}_2 + \langle F, c \rangle$ where $c$ is the ghost and $F$ the curvature.
  \item Adler--Bardeen 1969 supplies the physical template: a chiral anomaly satisfying Wess--Zumino consistency and one-loop exactness. Its direct transfer to holomorphic Chern--Simons on $X$ is a bridge hypothesis, not a theorem used without proof.
  \item For holomorphic Chern--Simons on $X$ in real dimension $2 \dim_\Cbb X = 6$, the Adler--Bardeen descent from the universal Chern character $\mathrm{ch}_4(F) \in \Omega^{8}(X)$ to the BRST anomaly in $\Omega^6(X)$ produces a Dolbeault cohomology class in $H^2(X, \Omega^1_X)$, and this class is the cocycle $[m_3, B^{(2)}]_X$: the Massey bracket $m_3$ is the triangle-diagram one-loop obstruction (equivalent to the $\Tor$-product of three ghost fields), and the Connes shift $B^{(2)}$ is the cyclic shift encoding ghost-number conservation in the descent.
\end{itemize}
The consistency condition $\delta [m_3, B^{(2)}]_X = 0$ is the BV master equation at $\hbar^1$ order in this comparison. Its vanishing as a cohomology class is the conditional anomaly-cancellation statement.
\end{theorem}

\begin{proof}
This is a conditional restatement of the bridge hypotheses. Costello--Li
provide the parametrix and holomorphic BV framework; Wess--Zumino
descent supplies the form of a local BRST anomaly; the cyclic HKR
comparison identifies the resulting class with the Dolbeault
representative only after the trace/contraction and normalisation
choices of Proposition~\ref{prop:cocycle-rep}. The original
Adler--Bardeen theorem does not by itself prove this holomorphic
comparison.
\end{proof}

\begin{remark}[Anomaly-matching scope]
\label{rem:anomaly-matching}
The Adler--Bardeen theorem asserts one-loop exactness of the chiral anomaly for abelian currents (Adler--Bardeen 1969). Non-abelian extension to $\End(T_X)$-valued gauge fields is governed by Wess--Zumino consistency (Wess--Zumino 1971) combined with the Bardeen--Zumino descent (Bardeen--Zumino 1984): the BRST-invariant anomaly polynomial of the non-abelian theory is the image of $\mathrm{ch}_{d/2 + 1}(F)$ in ghost-number $1$ under the Zumino descent (Zumino 1985). For holomorphic Chern--Simons on a Calabi--Yau threefold, the Dolbeault-refined descent lands in Dolbeault bidegree $(0, 2)$ with values in $\Omega^1_X$; the cocycle $[m_3, B^{(2)}]_X$ is this descent class. Its vanishing is anomaly cancellation, invariant under renormalisation-group flow by 't~Hooft anomaly matching ('t~Hooft 1980)---matching the three-scale persistence of $\hbar^2 K = -1$ (Theorem~\ref{thm:three-scale}).
\end{remark}

\begin{remark}[Compactness hypothesis for 6d hCS on non-compact targets]
\label{rem:hcs-compactness}
The 6d hCS theory of Costello 2013 and Costello--Li 2016 is defined on a real-six-dimensional target and constructed in two geometric regimes. On a smooth projective Calabi--Yau threefold $X$ the compactness hypothesis of \S\ref{sub:fa-precision} holds automatically (the projective embedding supplies a compact Kahler metric with uniformly bounded heat kernel), and Costello 2011 Theorem~13.4.1 gives existence and renormalisability of the BV action to all loop orders. On non-compact complex-threefold targets such as $K3 \times \Cbb$ the theorem requires cylindrical-end compatibility on the affine factor together with the twisted holomorphic extension of CG FA Vol~2 Theorem~8.6.9 on the K3 factor. Where this compatibility is not established --- on Clemens' non-Kahler small-resolution threefolds (\S\ref{sec:clemens}) and on asymptotically-flat non-cylindrical factorisations --- the hCS theory is a prefactorisation algebra rather than a factorisation algebra, and the Adler--Bardeen identification drops to a statement in the homotopy category, not in $H^2$ on the nose.
\end{remark}

\begin{remark}[Lane compatibility after rectification]
\label{rem:lane-compat}
The theorem records the cyclic-HKR construction problem and the
comparison maps needed to turn that class into a four-target bridge.
The former five strands are not all proved at chain level here: the
Positselski, factorisation, K3/Mukai, and trace arrows are precisely
the hypotheses \textup{(B1)--(B4)}.

The $(\infty,1)$-categorical proved statement is the existence of the natural-transformation class represented by the composite
\[
D^b(X) \xrightarrow{m_3 \text{ in } E_1} D^b(X) \xrightarrow{B^{(2)} \text{ in } fE_2/E_1} D^b(X) \otimes \Omega^1_X[2]
\]
admit a contractible space of coherent homotopies if the class in $\Ext^2_{D^b(\Coh X)}(\cO_X, \Omega^1_X) \simeq H^2(X, \Omega^1_X)$ vanishes. Under \textup{(B1)--(B4)} this derived locus can be compared with the four target loci; without those hypotheses there is no equality of vertices in $\mathrm{Pr}^L_k$.
\end{remark}

% ==================================================
\section{Witnesses: quintic, $K3 \times E$, Clemens}
\label{sec:clemens}
% ==================================================

\subsection{Generic quintic $Q_5 \subset \Pbb^4$}
Set $Q_5 = V(f)$ for a generic homogeneous quintic $f \in \Cbb[x_0, \ldots, x_4]_5$. Hodge numbers $h^{1,1}(Q_5) = 1$, $h^{2,1}(Q_5) = 101$ (Candelas--De la Ossa--Green--Parkes 1991).

\begin{proposition}[Generic quintic falsifier; \ClaimStatusConditional]
\label{prop:quintic-nonvanish}
For a generic quintic $Q_5$, Griffiths' Jacobian-ring computation gives
\[
  \dim H^2(Q_5,\Omega^1_{Q_5})=h^{1,2}(Q_5)=101
\]
and the Yukawa coupling is non-zero. Therefore a universal K3/Igusa
vanishing theorem cannot hold on the full Calabi--Yau-threefold
landscape. The stronger conclusion
$[m_3,B^{(2)}]_{Q_5}\neq 0$ is conditional on the comparison between
the Yukawa cubic and the cyclic-skew Atiyah--Connes projection.
\end{proposition}

\begin{proof}
The Jacobian-ring count is Griffiths 1969 and CDOGP 1991. Non-zero
Yukawa proves non-triviality of the cubic variation of Hodge structure.
It does not, by itself, identify the corrected class
$\alpha_{Q_5}\in H^{1,2}(Q_5)$ with a non-zero quadratic
cyclic-skew projection. That identification is exactly the comparison
hypothesis named in the statement.
\end{proof}

\begin{corollary}[Generic $Q_5$ as falsifier; \ClaimStatusConditional]
\label{cor:quintic-fails}
The generic quintic falsifies the unconditional pentagon: it has a
101-dimensional corrected obstruction target and non-zero Yukawa
coupling, whereas the K3/Igusa vertex lives on a special Mukai locus.
Failure of all four bridge vertices is conditional on the non-vanishing
comparison in Proposition~\ref{prop:quintic-nonvanish}.
\end{corollary}

\begin{proof}
This is the contrapositive form of Theorem~\ref{thm:main}: without the
comparison hypotheses and a vanishing result for $\alpha_{Q_5}$, the
four-target bridge has no theorem-grade value on the generic quintic.
\end{proof}

\subsection{Product $K3 \times E$}
Set $Y = K3 \times E$ for $K3$ an algebraic K3 surface and $E$ an elliptic curve. The product has trivial canonical bundle but fails the strict $h^{1,0}=0$ hypothesis; any Borcea--Voisin or related twist must be specified and its Hodge numbers recomputed. The raw product remains the useful Kunneth test case.

\begin{proposition}[Corrected $K3\times E$ obstruction target; \ClaimStatusConditional]
\label{prop:k3e-vanish}
Let $Y=S\times E$, where $S$ is a projective K3 surface and $E$ is an
elliptic curve. Then
\[
  T_Y=\pi_S^*T_S\oplus\pi_E^*T_E,\qquad
  \at_Y=\pi_S^*\at_S+\pi_E^*\at_E=\pi_S^*\at_S,
\]
with $\at_E=0$ and $\at_S\neq0$ in general. Moreover
\[
H^3(Y,\Omega_Y^3)\cong
H^2(S,\cO_S)\otimes H^1(E,\cO_E)\cong \Cbb,
\]
and
\[
H^2(Y,\Omega_Y^1)\cong
H^1(S,\Omega_S^1)\otimes H^1(E,\cO_E)
\oplus
H^2(S,\cO_S)\otimes H^0(E,\Omega_E^1),
\]
so $\dim_\Cbb H^2(Y,\Omega_Y^1)=21$. Consequently the vanishing of
$[m_3,B^{(2)}]_Y$ is not implied by Hodge numbers. It is equivalent
here to a scalar condition. Product naturality kills the
$20$-dimensional mixed summand
$H^1(S,\Omega_S^1)\otimes H^1(E,\cO_E)$; the remaining class lies on
the one-dimensional line
\[
  L_{S,E}:=H^2(S,\cO_S)\otimes H^0(E,\Omega_E^1).
\]
Thus for the chosen cyclic-HKR datum $\mathfrak c$ there is a scalar
$\lambda_{\mathrm{Cyc}}(S,E,\mathfrak c)$ such that
\[
  \alpha_{S\times E}
  =
  \lambda_{\mathrm{Cyc}}(S,E,\mathfrak c)\,
  [\bar\sigma_S]\otimes\eta_E,
  \qquad
  \alpha_{S\times E}=0
  \Longleftrightarrow
  \lambda_{\mathrm{Cyc}}(S,E,\mathfrak c)=0.
\]
Under standard Atiyah/Chern--Weil/Rozansky--Witten normalisation this
scalar is expected to be proportional to $\int_S c_2(T_S)=24$, hence
not automatically zero.
\end{proposition}

\begin{proof}
The tangent and Atiyah formulas are standard additivity under products.
The elliptic tangent bundle is trivial, so $\at_E=0$; the K3 tangent
bundle is not holomorphically trivial and its Atiyah class detects the
non-zero curvature/Chern data. The displayed cohomology groups are the
Kunneth decomposition for $\Omega_Y^3$ and $\Omega_Y^1$. The two
summands of $H^2(Y,\Omega_Y^1)$ have dimensions $20$ and $1$,
respectively. Because $\at_Y^{\cup2}$ has only K3 legs, a
product-natural cyclic-HKR projection cannot manufacture the
$H^1(E,\cO_E)$ factor needed for the mixed summand. It can only leave a
scalar on $H^2(S,\cO_S)\otimes H^0(E,\Omega_E^1)$, proving the
reduction to $\lambda_{\mathrm{Cyc}}$.
\end{proof}

\begin{proposition}[Explicit K\"unneth chain on $K3 \times E$; \ClaimStatusConditional]
\label{prop:k3e-kunneth-chain}
The obstruction target on $Y=K3\times E$ is a seven-step Kunneth calculation. It does not prove vanishing; it isolates the one remaining scalar $\lambda_{\mathrm{Cyc}}$.
\begin{align*}
&\textup{(Step 1, K\"unneth for tangent bundle)} && T_{K3 \times E} \;=\; \pi_1^* T_{K3} \oplus \pi_2^* T_E. \\
&\textup{(Step 2, Atiyah-class additivity)} && \at_{K3 \times E} \;=\; \pi_1^* \at_{K3} \;+\; \pi_2^* \at_E. \\
&\textup{(Step 3, triviality on elliptic factor)} && \at_E \;=\; c_1(T_E) \;=\; 0 \quad (T_E \cong \cO_E). \\
&\textup{(Step 4, cocycle square K\"unneth)} && \at_{K3 \times E}^{\cup 2} \;=\; \pi_1^*\at_{K3}^{\cup 2} \;+\; 2\,\pi_1^*\at_{K3}\cup\pi_2^*\at_E \;+\; \pi_2^*\at_E^{\cup 2} \;=\; \pi_1^*\at_{K3}^{\cup 2}. \\
&\textup{(Step 5, Dolbeault K\"unneth for receptacle)} && H^2(K3 \times E, \Omega^1_{K3 \times E}) \;=\; \bigoplus_{p+q=2} H^p(K3, \Omega^1_{K3}) \otimes H^q(E, \cO_E)\\
&&&\quad \oplus \bigoplus_{p+q=2} H^p(K3, \cO_{K3}) \otimes H^q(E, \Omega^1_E). \\
&\textup{(Step 6, non-zero summands)} && H^2(Y,\Omega^1_Y)
\cong H^1(K3,\Omega^1_{K3})\otimes H^1(E,\cO_E)
\oplus H^2(K3,\cO_{K3})\otimes H^0(E,\Omega^1_E). \\
&\textup{(Step 7, cocycle condition)} && [m_3, B^{(2)}]_{Y}=0
\Longleftrightarrow
\lambda_{\mathrm{Cyc}}(K3,E,\mathfrak c)=0
\text{ on the scalar line of Step 6.}
\end{align*}
Steps~1--7 compose: Kunneth splits the construction and shows exactly
where a proof of vanishing must act. The target is $21$-dimensional,
product naturality removes the mixed $20$-dimensional piece, and the
remaining problem is the one-dimensional scalar
$\lambda_{\mathrm{Cyc}}$.
\end{proposition}

\begin{proof}
Step~1: K\"unneth for the tangent bundle of a product of complex
manifolds. Step~2: additivity of the Atiyah class under $\oplus$
(Atiyah 1957 \S4). Step~3: $T_E = \cO_E$ trivial on a
one-dimensional abelian variety. Step~4:
K\"unneth-multiplicativity of the cup product, plus Step~3. Step~5:
K\"unneth formula for Dolbeault cohomology on a product
(Griffiths--Harris 1978 III.3). Step~6: standard Hodge numbers of K3
and $E$. Step~7: product naturality of the fixed cyclic-HKR datum.
\end{proof}

\begin{corollary}[Koszul locus on $K3 \times E$; \ClaimStatusConditional]
\label{cor:k3e-locus}
On $K3 \times E$ the statement that $\cU^{\mathrm{adm}}(K3 \times E)$ is the full $21$-dimensional product moduli minus Humbert divisors is conditional on the scalar vanishing $\lambda_{\mathrm{Cyc}}=0$ of Proposition~\ref{prop:k3e-vanish} and on recomputing the Hodge data after any Borcea--Voisin twist. Under those hypotheses the four-target bridge holds on the corresponding locus by Theorem~\ref{thm:main}.
\end{corollary}

\begin{proof}
The dimension count for the product moduli is $20+1=21$. It is not a
vanishing proof for $\alpha_Y$. The stated locus follows only after the
$\lambda_{\mathrm{Cyc}}=0$ condition and the Mukai/Humbert admissibility
condition are imposed.
\end{proof}

\subsection{Clemens' non-Kähler small-resolution threefolds}
Clemens 1983 constructs non-Kähler smoothings of nodal threefolds by performing small resolutions of ordinary double points. Friedman 1986 (Theorem~5.10) shows that for suitable configurations, the resulting smooth threefold $\tilde Y$ is non-Kähler, and the $\partial\bar\partial$-lemma fails.

\begin{proposition}[$\partial\bar\partial$-lemma failure; \ClaimStatusProvedElsewhere]
\label{prop:clemens-ddbar}
For Clemens' small-resolution threefolds $\tilde Y$, the Frölicher spectral sequence is $E_1$-non-degenerate: there exists $d_1^{0,3} \neq 0$ in the first page, and the Dolbeault decomposition of $H^\bullet(\tilde Y, \Cbb)$ is not valid.
\end{proposition}

The Frölicher spectral sequence interpolates between Dolbeault cohomology (the $E_1$-page, computed from $\bar\partial$ alone) and de Rham cohomology (the $E_\infty$-page, computed from the full $d = \partial + \bar\partial$). Its $E_1$-differential is $d_1 = [\partial]$: it sends a $\bar\partial$-cohomology class $[\omega] \in H^{p,q}_{\bar\partial}$ to $[\partial \omega]$, which is automatically $\bar\partial$-closed (since $\bar\partial \partial = - \partial \bar\partial$ and $\omega$ is $\bar\partial$-closed). On a compact Kähler manifold, $d_1 = 0$ by the Kähler identities, so $E_1 = E_\infty$ and Dolbeault = de Rham, cohomology class by cohomology class. This is what the $\partial\bar\partial$-lemma records: a $\bar\partial$-closed, $\partial$-closed form is $\partial\bar\partial$-exact on the nose. For Clemens' $\tilde Y$, the presence of a non-zero $d_1^{0,3} : H^{0,3}_{\bar\partial} \to H^{1,3}_{\bar\partial}$ means that some $\bar\partial$-closed $(0,3)$-form $\omega$ has $\partial \omega$ representing a non-zero Dolbeault class, so $\omega$ itself is \emph{not} representable by a form that is simultaneously $\partial$- and $\bar\partial$-closed. Dolbeault representatives of de Rham classes do not exist.

The Dolbeault-model construction of $[m_3, B^{(2)}]$ (Proposition~\ref{prop:cocycle-rep}) therefore fails on $\tilde Y$: the representative $\bar\partial \nabla$ defines a cohomology class only when every $\bar\partial$-closed form admits a $\partial$-closed refinement up to $\bar\partial$-exact error. Without the $\partial\bar\partial$-lemma, the dependence on $\nabla$ escapes $\bar\partial$-exactness, and distinct Kähler metrics produce distinct classes in $H^{0,2}_{\bar\partial}(\tilde Y, \Omega^1_{\tilde Y})$ differing by $\partial$-exact-but-not-$\bar\partial$-exact terms. The cocycle is not well-defined.

\begin{proposition}[Clemens scope warning; \ClaimStatusConditional]
\label{prop:clemens-explicit}
For Clemens' small-resolution threefold $\tilde Y$ with non-degenerate
Fr\"olicher spectral sequence, the Dolbeault representative
$\bar\partial\nabla\wedge\bar\partial\nabla$ is not a canonical
Atiyah--Connes class. The formal mechanism previously written as
``$\partial\bar\partial$-exact but not $\bar\partial$-exact'' cannot
prove non-vanishing: every $\partial\bar\partial$-exact form is already
$\bar\partial$-exact. The correct conclusion is only that the
Kahler/Dolbeault construction does not descend unchanged to Clemens'
non-Kahler stratum.

A replacement must be built in the deformation hypercohomology of a
neutralized heterotic Courant or chiral-differential-operator
algebroid. In that setting one must specify a closed three-form $H$,
a neutralization of the Gorbounov--Malikov--Schechtman CDO obstruction
gerbe, and a cyclic-HKR analogue for the twisted deformation complex.
\end{proposition}

\begin{proof}
The first assertion is Proposition~\ref{prop:clemens-ddbar}: the
Kahler Hodge-theoretic proof of representative independence uses
degeneration and the $\partial\bar\partial$-lemma. The displayed
``non-vanishing'' mechanism is invalid because
$\partial\bar\partial\alpha=\bar\partial(-\partial\alpha)$. The
surviving theorem-grade statement is therefore negative: the
Dolbeault proof does not produce a canonical class on this stratum.
The proposed Courant/CDO replacement is conditional on the listed
additional data.
\end{proof}

\emph{Factorisation-algebra behaviour on non-Kähler CY-$3$.} The stage-one functor $\Phi^{\mathrm{FA}}_3$ of \S\ref{sub:fa-precision} requires two geometric inputs: the compactness hypothesis (uniformly bounded heat-kernel propagators) and a holomorphic refinement of the $E_3$-structure through Kontsevich $E_d$-formality. On Clemens' non-Kähler $\tilde Y$ the first input survives --- $\tilde Y$ is a compact smooth six-manifold, so the Costello 2011 heat-kernel renormalisation runs on the nose; the \emph{holomorphic} refinement fails, because the Kontsevich formality of Kontsevich 1999 is a theorem over a Kähler manifold with $\partial\bar\partial$-lemma and its extension to non-Kähler compact manifolds is not known to produce a holomorphic $E_3$-factorisation algebra. What does survive: a topological $E_3$-factorisation algebra $\Phi^{\mathrm{FA}, \mathrm{top}}_3(\tilde Y) \in \mathrm{Fact}^{E_3, \mathrm{top}}_{\tilde Y}$ (the underlying real-$6$-manifold factorisation algebra of the BV-quantised topological sigma-model), and an $H$-twisted chiral-differential-operator sheaf $\cD^{\mathrm{ch}, H}(\tilde Y)$ whose generic fibre is the $H$-twisted Courant algebroid of Beilinson--Drinfeld 2004 and Gorbounov--Malikov--Schechtman 2003 with twist class $[H] \in H^3(\tilde Y, \Rbb)$ of Strominger 1986. The two are related by an $H$-twisted comparison: stage one reads
\[
  \Phi^{\mathrm{FA}}_3(\tilde Y) \; \text{is replaced by} \; \Phi^{\mathrm{FA}, H}_3(\tilde Y) := \bigl(\Phi^{\mathrm{FA}, \mathrm{top}}_3(\tilde Y), \text{Courant twist by } [H]\bigr) \in \mathrm{Fact}^{E_3, H}_{\tilde Y}
\]
with $\mathrm{Fact}^{E_3, H}_{\tilde Y}$ the category of $E_3$-factorisation algebras enriched by a compatible Courant-algebroid structure, and the canonicality pinning of \S\ref{sub:fa-precision} drops from Kontsevich--Tamarkin formality to the GMS--Courant pinning (which supplies a CDO attached to $[H]$ but does not impose holomorphicity). Stage two $\Sp^{\mathrm{ch}}_{\Sigma_2, C}$ proceeds as in the Kähler case, now on the $H$-twisted stage-one output; the resulting chiral algebra on $C$ is the $H$-twisted chiral differential operator algebra, and the cocycle is its universal obstruction.

The reformulation of the cocycle construction requires an ambient that does not presuppose the $\partial\bar\partial$-lemma. The Courant algebroid (Courant 1990) furnishes this ambient: it is the bundle $T_{\tilde Y} \oplus T^*_{\tilde Y}$ with symmetric pairing $\langle X + \xi, Y + \eta \rangle = \iota_X \eta + \iota_Y \xi$ and Dorfman bracket $[X + \xi, Y + \eta]_H = [X, Y] + \mathcal{L}_X \eta - \iota_Y d\xi + \iota_X \iota_Y H$, twisted by a closed $3$-form $H \in \Omega^3(\tilde Y, \Rbb)$. The Courant--Severa class $[H] \in H^3(\tilde Y, \Rbb)$ classifies $H$-twisted Courant algebroids on $\tilde Y$ up to isomorphism; the Strominger flux (Strominger 1986) supplies a preferred representative on a balanced-Hermitian $\mathrm{SU}(3)$-structure only after the Bianchi identity is imposed. The chiral-differential-operator sheaf (Beilinson--Drinfeld 2004, Gorbounov--Malikov--Schechtman 2003) attached to a neutralized Courant algebroid is a vertex algebra whose $\mathrm{SL}(2)$-part is the holomorphic tangent direction. If $H$ is closed, the twisted differential $d^H=d+H\wedge(-)$ satisfies $(d^H)^2=0$ in the relevant parity convention; if $dH\neq0$, the Bianchi defect is itself part of the anomaly data.

\begin{corollary}[Strominger heterotic-with-torsion reformulation; \ClaimStatusConditional]
\label{cor:strominger}
On Clemens' stratum, the four comparison targets and the cocycle $[m_3, B^{(2)}]$ are expected to admit reformulations via the Strominger (1986) heterotic-with-torsion framework on balanced-Hermitian $\mathrm{SU}(3)$-structures. The Koszul-bar backbone must be replaced by a neutralized $H$-twisted Courant-algebroid chiral differential operator construction; the anomaly should be encoded in the corresponding deformation hypercohomology, not in a raw group $H^2(\tilde Y,\Omega^1_{\tilde Y})_H$.
\end{corollary}

In this regime the bar-complex source is not holomorphic: its differential is $d^H = d + H\wedge$ for a suitable closed $H \in \Omega^3(\tilde Y, \Rbb)$. Each of the four comparison statements requires a twisted generalised-geometry replacement; none follows by verbatim substitution from the Kahler proof.

\begin{theorem}[Strominger heterotic-with-torsion as native Clemens physics; \ClaimStatusHeuristic]
\label{thm:strominger-native}
On Clemens' non-K\"ahler small-resolution threefolds $\tilde Y$, the four comparison targets of Theorems~\ref{thm:climax-I}--\ref{thm:climax-IV} are expected physical reductions of Strominger 1986's heterotic superstring with torsion on an $\mathrm{SU}(3)$-structure manifold, under the physical identifications:
\begin{itemize}
  \item The $H$-flux $H \in \Omega^3(\tilde Y, \Rbb)$ is the Neveu--Schwarz $3$-form field strength of the heterotic ten-dimensional superstring, satisfying $dH = \mathrm{tr}(R \wedge R) - \mathrm{tr}(F \wedge F)$ at one loop (Green--Schwarz mechanism, Green--Schwarz 1984, Strominger 1986 equation (8)), and characterising the failure of $\tilde Y$ to be K\"ahler: $\partial \bar\partial J = - i/2 \cdot dH$ for $J$ the Hermitian form.
  \item The $H$-twisted differential $d^H = d + H \wedge$ is the Courant-algebroid differential on $T\tilde Y \oplus T^*\tilde Y$ whose restriction to the tangent bundle reproduces the Bismut connection of the balanced-Hermitian $\mathrm{SU}(3)$-structure (Michelsohn 1982 Theorem~6.8). Strominger 1986 equations (12)--(14) require the Bismut connection (not the Levi-Civita connection, which does not preserve the Hermitian structure when $dJ \neq 0$) to have holonomy contained in $\mathrm{SU}(3)$; this is equivalent to $\tilde Y$ admitting a supersymmetric $(0, 2)$ heterotic vacuum.
  \item The $H$-twisted Courant-algebroid chiral differential operator algebra $\mathrm{CDO}^H(\tilde Y)$ of Malikov--Schechtman 1999 (as extended to Courant algebroids by Malikov--Schechtman--Vaintrob 1999) is the worldsheet chiral algebra of the heterotic $(0, 2)$ sigma model on $\tilde Y$ with $H$-flux. The right-moving central charge is $c_R = 3 \dim_\Cbb(\tilde Y) = 9$ (matching the $\cN = 2$ superconformal ghost-matter target); the left-moving central charge is $c_L = 22$ when anomaly-cancelled against the $E_8 \times E_8$ or $\mathrm{Spin}(32)/\Zbb_2$ heterotic gauge bundle. The $H$-twisted $\cN = 2$ algebra on the right-moving side is the current algebra of $\mathrm{CDO}^H$.
  \item The expected replacements are: (\ref{thm:climax-I}) an $H$-twisted chiral Positselski statement on $\mathrm{CDO}^H$; (\ref{thm:climax-II}) a heterotic torsion-class protected index on $\tilde Y \times E$, with $\Phi_{10}$ replaced by a Siegel modular form for a congruence subgroup reflecting the Bismut torsion; (\ref{thm:climax-III}) an $H$-twisted chiral bialgebra $\bfH^H_{\Delta_5}$ on the Bismut period-domain stratum; (\ref{thm:climax-IV}) a trace identity with a Strominger anomaly correction term.
\end{itemize}
\end{theorem}

\begin{proof}
The $\partial \bar\partial$-lemma failure (Proposition~\ref{prop:clemens-ddbar}) is the obstruction to the Kahler Dolbeault proof. The Strominger 1986 system supplies a physical framework with Hermitian form $J$, NS three-form $H$, and Bismut connection; it does not by itself construct the cyclic-HKR replacement. A theorem-grade construction must build the neutralized Courant/CDO deformation complex and compare its cochain Connes rotation with the Kahler class of Section~\ref{sec:cocycle}.
\end{proof}

\begin{remark}[Why the Clemens stratum is the native heterotic stratum]
\label{rem:clemens-heterotic}
Among the six strata of Proposition~\ref{prop:strata}, Clemens' non-K\"ahler stratum is the natural stress test for the heterotic replacement. The K\"ahler strata are the $H = 0$ specialisation where the Dolbeault cyclic-HKR construction is available; the Clemens stratum asks whether a neutralized Courant/CDO cyclic-HKR construction exists.
\end{remark}


\subsection{Six-stratum CY-$3$ partition}

\begin{proposition}[Stratum partition; \ClaimStatusConditional]
\label{prop:strata}
The smooth projective Calabi--Yau threefold landscape decomposes into six strata, distinguished by the geometric behaviour of the cocycle $[m_3, B^{(2)}]_X$:

\begin{tabular}{l l l l}
\toprule
 & Description & $\at_X$ behaviour & Climax behaviour \\
\midrule
1 & Toric / $T^6$ & $\at = 0$ trivially & Degenerate, $K = 0$ \\
2 & Schoen K3-elliptic fibred, $h^{2,1} = 19$ & $\at_X \cup \at_X$ supp.\ on $24$-pt.\ discr. & Maximal three-lens canonicity \\
3 & Borcea--Voisin, $h^{2,1}$ Nikulin-dependent & $\at_X \cup \at_X$ on Nikulin-twist divisor & $\sfB$-row sub-stratified, four Nikulin types \\
4 & $K3 \times E$ orbifold, $h^{2,1} = 21$ & Kunneth-additive $\at$; $H^2(\Omega^1)$ has dim.\ $21$ & Conditional scalar $\lambda_{\mathrm{Cyc}}$ witness \\
5 & Generic quintic $Q_5$, dim $101$ & $\at$ generically non-zero & Koszul locus empty \\
6 & Clemens non-Kähler & $\partial\bar\partial$-lemma fails & Strominger reformulation \\
\bottomrule
\end{tabular}
\end{proposition}

\begin{proof}[Proof sketch]
Stratum $1$: on $T^6$ or toric CY-$3$, the tangent bundle is trivial or a sum of line bundles, $\at_X = 0$ by triviality. Stratum $2$: Schoen $1988$'s construction is the fibred product $X = \mathrm{K3} \times_{\Pbb^1} E$ of a K3-elliptic fibration and an elliptic fibration over $\Pbb^1$ with $24$ singular fibres, giving a smooth CY-$3$ with Hodge numbers $(h^{1,1}, h^{2,1}) = (19, 19)$ and self-mirror; the Atiyah class $\at_X$ concentrates on the $24$-point Kodaira discriminant of the elliptic fibration. Stratum $3$: Borcea--Voisin constructions (Borcea $1997$, Voisin $1993$) with involution fix the involution class of $K3$, giving the Nikulin sub-stratification; Hodge numbers range $h^{2,1} \in \{14, 23, 35, 50, 86\}$ across the Nikulin classes. Strata $4$ (K3 $\times E$ orbifold), $5$ (quintic), $6$ (Clemens) are the witnesses above. The table's last two dimension columns refer to $h^{2,1}(X)$ rather than the complex-structure moduli dimension.
\end{proof}

% ==================================================
\section{Corollaries}
\label{sec:corollaries}
% ==================================================

\subsection{Two Koszul loci intersection}

\begin{theorem}[Two loci intersection; \ClaimStatusConditional]
\label{thm:two-loci}
Let $X$ be a smooth Calabi--Yau threefold with period map $\cP : \cM_{\mathrm{cx}}(X) \to \Gamma \backslash D$ to the period domain. The stage-one canonical locus $\cU^{\mathrm{canonical}}_{\Phi_3}$ and the cocycle-admissible locus $\cU^{\mathrm{adm}}_{\at}(X)$ intersect as
\[
  \cU^{\mathrm{canonical}}_{\Phi_3} \;=\; \cU^{\mathrm{adm}}_{\at}(X) \;\cap\; \cP^{-1}\bigl(U^{\mathrm{adm}}\bigr),
\]
with $U^{\mathrm{adm}} \subset \Gamma \backslash D$ the Kuga--Satake period image of the admissible locus.
\end{theorem}

\begin{proof}
The period map $\cP$ classifies complex structures via Hodge decomposition. Kuga--Satake (Kuga--Satake 1967, Morrison 1984) lifts the K3 period domain to an abelian-variety period domain; for CY-$3$ with K3 fibration, the pullback of the admissible period locus via $\cP$ carves out precisely the complex structures for which the cocycle vanishes via Hodge-theoretic transversality. Combining with stage-one canonicality (Theorem~\ref{thm:main}(iii)) gives the intersection identity.
\end{proof}

\subsection{Nine Heegner discriminants at $\rho = 20$}

\begin{proposition}[Heegner discriminants; \ClaimStatusConditional]
\label{prop:heegner}
For K3 surfaces of Picard rank $\rho = 20$ (singular K3 surfaces in the sense of Shioda--Inose 1977), the classical class-number-one discriminants are
\[
  d_K \in \{3,4,7,8,11,19,43,67,163\}.
\]
The alternative nine-element list
\[
  d_K \;\in\; \{3, 4, 7, 8, 11, 15, 19, 20, 24\},
\]
is a conditional Bruinier--Chenevier refinement: it requires an explicit coefficient/sign computation showing that $\{15,20,24\}$ replace $\{43,67,163\}$ at rank $K=8$. Each $d_K$ in whichever verified list is used should produce an arithmetic chiral vertex algebra $V^{\mathrm{arith}}_{K3,d_K}$ with central charge $K=8$ twisted by the $\Qbb(\sqrt{-d_K})$-arithmetic data.
\end{proposition}

The thread from Picard rank $\rho = 20$ to imaginary quadratic discriminants passes through three named identifications. (i) A K3 surface is \emph{singular} in the Shioda--Inose sense if its Picard rank achieves the maximum $\rho = 20$: all of $H^{1,1}(K3, \Rbb)$ comes from algebraic classes. The transcendental lattice $T(K3) = \mathrm{Pic}(K3)^\perp \subset H^2(K3, \Zbb)$ then has rank $22 - 20 = 2$ and signature $(2, 0)$ --- a positive-definite even rank-$2$ lattice. (ii) Positive-definite rank-$2$ even lattices are in bijection with $\mathrm{SL}_2(\Zbb)$-equivalence classes of positive-definite binary quadratic forms $ax^2 + bxy + cy^2$ with $b^2 - 4ac = -d_K < 0$; the bijection sends $T$ to the form with Gram matrix $T$. (iii) Positive-definite binary quadratic forms of discriminant $-d_K$ are in bijection with proper ideal classes of the imaginary quadratic order $\cO_{d_K} \subset \Qbb(\sqrt{-d_K})$; the bijection sends the form to the ideal class of $(a, (-b + \sqrt{-d_K})/2)$. A singular K3 therefore encodes a CM point on the Shimura variety $\mathrm{O}(2, 20)/(\mathrm{SO}(2) \times \mathrm{O}(20))$ (Kuga--Satake 1967, Morrison 1984) with Mumford--Tate group inside the norm-one units of $\cO_{d_K}$.

\begin{proof}[Proof sketch]
Nikulin 1979 classifies lattice-polarised K3 surfaces by transcendental lattice. At $\rho = 20$ (singular K3), the transcendental lattice $T$ has rank $2$ and signature $(2,0)$, and Shioda--Inose 1977 establish a bijection between isomorphism classes of such $T$ and equivalence classes of positive-definite binary quadratic forms of discriminant $-d_K$. The classical Heegner--Stark--Baker list gives $\{3,4,7,8,11,19,43,67,163\}$. The proposed rank-$8$ replacement by $\{3,4,7,8,11,15,19,20,24\}$ is not proved by the class-number-one theorem; it is precisely the conditional Bruinier coefficient computation named in the statement.
\end{proof}

\subsection{Corollary for Bridgeland finiteness}

\begin{corollary}[Compactness of $\cU^{\mathrm{adm}}$; \ClaimStatusConditional]
\label{cor:bridgeland}
Conditional on Bridgeland's conjecture that the space of stability conditions $\Stab(D^b(X))$ is a finite-dimensional complex manifold with locally-closed period maps to $\cM_{\mathrm{cx}}(X)$: the Koszul-admissible locus $\cU^{\mathrm{adm}}(X)$ is a locally finite union of strata in $\Stab(D^b(X))$, each with finitely many Bridgeland walls bounding it.
\end{corollary}

\begin{proof}[Proof sketch, conditional on Bridgeland]
The $A_\infty$-structure on $\End^\bullet(\cE)$ for a tilting $\cE \in D^b(X)$ depends on the Bridgeland stability condition $\sigma \in \Stab(D^b(X))$ through the heart of $\sigma$; on a chamber the heart is fixed and $[m_3, B^{(2)}]_X$ varies holomorphically in $\sigma$ with values in $H^2(X, \Omega^1_X)$. The vanishing locus is analytic-closed within each chamber; Bridgeland's locally-closed period-map hypothesis maps this to an analytic locus in $\cM_{\mathrm{cx}}(X)$. Walls are codimension-$1$ loci where the heart jumps; Bridgeland's finiteness conjecture bounds the number of walls meeting any compact subset, giving local finiteness of $\cU^{\mathrm{adm}}$-strata.
\end{proof}

\subsection{Forbidding the twenty-second Niemeier row}

\begin{corollary}[Non-Leech Niemeier $\sfB$-row; \ClaimStatusConditional]
\label{cor:m23}
Conditional on the Niemeier $1973$ classification of the $24$ even unimodular rank-$24$ lattices with non-zero root system, and on the Chenevier--Lannes $2019$ mass-formula refinement for their $M_{23}$-automorphism twist, the $22$-row $M_{23}$ non-Leech Niemeier BKM algebras admit a canonical $\sfB$-row chiral bialgebra presentation if and only if the cocycle $[m_3, B^{(2)}]$ vanishes on the corresponding Leech-complementary moduli component.
\end{corollary}

\begin{proof}[Proof sketch, conditional on Niemeier $1973$ and Chenevier--Lannes $2019$]
Niemeier $1973$ enumerates the $24$ positive-definite even unimodular rank-$24$ lattices: the Leech lattice (root-free) and $23$ with non-zero root systems, each labelled by its root system. The $M_{23}$ Mathieu-invariant row corresponds to the Niemeier component $A_1^{24}$, whose automorphism group contains the Mathieu group $M_{23}$ as a sporadic subquotient of the Weyl-group stabiliser (Conway--Sloane, Chenevier--Lannes $2019$); the $\sfB$-row source-recognition target of Theorem~\ref{thm:climax-III} transports to this component iff the Mukai-lattice pinning extends across the Leech-complementary moduli component. By the conditional bridge of Theorem~\ref{thm:main}, this extension would follow from the corresponding Atiyah--Connes vanishing together with the K3/Mukai comparison hypotheses.
\end{proof}

% ==================================================
\section{Hand-computable examples: five witnesses in ten minutes each}
\label{sec:hand-computable}
% ==================================================

The Witnesses section descends to structural regimes. The present section descends to explicit generators. Each example admits a ten-minute hand computation that establishes either a proved input or a named conditional bridge. The rank-$1$ Heisenberg and affine $V_1(\mathfrak{sl}_2)$ computations are algebraic toy models; the quintic and $K3\times E$ computations are geometric tests of the corrected obstruction target.

\subsection{The quintic: $h^{2,1} = 101$ counted in the Jacobian ring}
\label{sub:hand-quintic}

\begin{example}[Dimension count via $R_f^{(5)}$ and $R_f^{(10)}$]
\label{ex:hand-quintic}
The falsifier computation of Proposition~\ref{prop:quintic-nonvanish} admits the following hand count. Take $Q_5 = V(f) \subset \Pbb^4$ with $f = x_0^5 + \cdots + x_4^5$ (Fermat); the count is independent of the smooth defining polynomial. Griffiths 1969 Theorem~4.2 identifies the primitive deformation-space piece $H^{2,1}(Q_5)$ with $R_f^{(5)}$ and the corrected obstruction target $H^{1,2}(Q_5)$ with the socle-dual piece $R_f^{(10)}$ of the Jacobian ring
\[
  R_f \;=\; \Cbb[x_0, \ldots, x_4] \big/ (\partial_0 f, \ldots, \partial_4 f) \;=\; \Cbb[x_0, \ldots, x_4] \big/ (x_0^4, x_1^4, x_2^4, x_3^4, x_4^4),
\]
with socle $R_f^{(15)} = \Cbb \cdot x_0^3 x_1^3 x_2^3 x_3^3 x_4^3$ and pairing $R_f^{(5)}\otimes R_f^{(10)}\to R_f^{(15)}$. The dimension of $R_f^{(10)}$ counts monomials $x_0^{a_0} \cdots x_4^{a_4}$ with $\sum a_i = 10$ and $0 \leq a_i \leq 3$. Inclusion--exclusion:
\[
  \tbinom{14}{4} \;-\; 5\tbinom{10}{4} \;+\; 10\tbinom{6}{4} \;=\; 1001 - 1050 + 150 \;=\; 101.
\]
This reproduces $h^{1,2}(Q_5)=h^{2,1}(Q_5)=101$ of Candelas--de la Ossa--Green--Parkes 1991. The Yukawa coupling $Y : (R_f^{(5)})^{\otimes 3} \to R_f^{(15)} = \Cbb$ detects non-vanishing of the cubic variation: choose $\mu_1 = x_0^2 x_1^2 x_2$, $\mu_2 = x_0 x_3^2 x_4^2$, $\mu_3 = x_1 x_2^2 x_3 x_4$ in $R_f^{(5)}$. Their product is $x_0^3 x_1^3 x_2^3 x_3^3 x_4^3$, the socle generator, so $Y(\mu_1, \mu_2, \mu_3) = 1$ after socle-normalisation. The further implication to non-vanishing of $[m_3,B^{(2)}]_{Q_5}\in H^{1,2}(Q_5)$ requires computing the residue image $\Psi_f(\alpha_{Q_5})\in R_f^{(10)}$.
\end{example}

\subsection{$K3 \times E$: K\"unneth selection via the Hodge diamond of K3}
\label{sub:hand-k3e}

\begin{example}[Hodge-diamond selection in the K\"unneth decomposition]
\label{ex:hand-k3e}
Proposition~\ref{prop:k3e-vanish} gives the corrected target. The hand computation isolates the obstruction to the old dimension argument: the Hodge diamond of K3
\[
  \begin{array}{ccccc}
    & & 1 & & \\
    & 0 & & 0 & \\
    1 & & 20 & & 1 \\
    & 0 & & 0 & \\
    & & 1 & &
  \end{array}
\]
has $h^{2,1}(K3) = h^{1,2}(K3) = 0$, but this is not the whole product target. In the Kunneth decomposition
\[
  H^2(K3 \times E, \Omega^1_{K3 \times E}) \;=\; \bigoplus_{p + q = 2,\, r + s = 1} H^p(K3, \Omega^r_{K3}) \otimes H^q(E, \Omega^s_E),
\]
the surviving summands are
\[
H^1(K3,\Omega^1_{K3})\otimes H^1(E,\cO_E)
\oplus
H^2(K3,\cO_{K3})\otimes H^0(E,\Omega^1_E),
\]
of total dimension $20+1=21$. Thus $[m_3,B^{(2)}]_{K3\times E}=0$ is
the scalar computation $\lambda_{\mathrm{Cyc}}=0$ in a non-zero target,
not a Hodge-diamond vanishing.

\emph{Two-line dimension check.} K3 complex-structure moduli is $20$-dim; elliptic $j$-moduli is $1$-dim. The same number $21$ also appears as the corrected obstruction target. Any claim that the Koszul locus is the full product moduli minus Humbert divisors is conditional on the scalar $\lambda_{\mathrm{Cyc}}$ vanishing and on recomputation after the chosen twist.
\end{example}

\subsection{The Heisenberg $\cH_1$: double vanishing $m_3 = B^{(2)} = 0$}
\label{sub:hand-heisenberg}

\begin{example}[$[m_3, B^{(2)}]_{\cH_1} = 0$ from commutativity and rank-$1$ antisymmetrisation]
\label{ex:hand-heisenberg}
Let $\cH_1$ be the rank-$1$ Heisenberg vertex algebra: single current $\alpha(z)$, OPE $\alpha(z)\alpha(w) \sim (z-w)^{-2}$, level $\kappa(\cH_1) = 1$. As an associative algebra, $\cH_1$ is graded commutative (the universal enveloping of the abelian Lie algebra $\Cbb \alpha$); vertex-algebra non-commutativity is carried by higher OPE modes, not by $m_2$. The ordered chiral bar complex is $\Bar^{\mathrm{ch}}_C(\cH_1) = T^c_C(s^{-1} \bar \cH_1)$.

\emph{$m_3 = 0$.} For any three homogeneous $f, g, h \in \End^\bullet(\cH_1)$, associativity of composition gives $(f \circ g) \circ h = f \circ (g \circ h)$; the Massey bracket $m_3(f, g, h) = 0$ on the nose, not up to homotopy. All $m_n$ for $n \geq 3$ vanish identically.

\emph{$B^{(2)} = 0$.} HKR gives $\HH^\bullet(\cH_1) \cong \cH_1 \otimes \wedge^\bullet(\Cbb \alpha)$. The Connes cyclic shift is the $\Zbb/3$-antisymmetrisation of the tri-exterior-product component; on rank-$1$ generator space $\Cbb \alpha$, $\alpha \wedge \alpha = 0$ kills $\wedge^3(\Cbb \alpha) = 0$. So $B^{(2)}|_{\HH^2(\cH_1)} \equiv 0$.

\emph{Double vanishing.} The cocycle $[m_3, B^{(2)}]_{\cH_1} = 0$ vanishes from two independent mechanisms. The toy analogues of the four bridge vertices are unobstructed: Theorem~B is Feigin--Frenkel bar-cobar inversion (the $\sfG$-row archetype), the partition function reduces to the free-boson character, the $\bfH_{\Delta_5}$ vertex is represented only by the scalar $K=1$ abelian target in this toy comparison, and the normalized trace identity $\hbar^2 \cdot K = -1$ gives $\hbar^2 = -1$ on the rank-$1$ sector.

\emph{Ten-minute verification.} On the triple tensor $\cH_1^{\otimes 3}$, direct OPE:
\[
  [\alpha(z), \alpha(w)] \alpha(v) \;-\; \alpha(z) [\alpha(w), \alpha(v)] \;=\; \partial_z \delta(z - w) \alpha(v) \;-\; \alpha(z) \partial_w \delta(w - v);
\]
cyclic rotation $z \to w \to v \to z$ maps the first term to its negative image in the second, and the sum cancels. This is $m_3 = 0$ at the level of OPE contractions.
\end{example}

\subsection{Affine $V_1(\mathfrak{sl}_2)$: $\epsilon^{abc} \kappa^{ab} = 0$ by symmetry-antisymmetry contraction}
\label{sub:hand-sl2}

\begin{example}[$[m_3, B^{(2)}]_{V_1(\mathfrak{sl}_2)} = 0$ via Killing form contracted with Levi-Civita]
\label{ex:hand-sl2}
Let $V_1(\mathfrak{sl}_2)$ be the level-$1$ affine Kac--Moody vertex algebra with currents $J^a(z)$, $a \in \{+, -, 0\}$, and OPE
\[
  J^a(z) J^b(w) \;\sim\; \frac{k \kappa^{ab}}{(z - w)^2} \;+\; \frac{f^{ab}_c J^c(w)}{z - w},
\]
Killing form $\kappa^{+-} = \kappa^{-+} = 1$, $\kappa^{00} = 2$; structure constants $f^{+-}_0 = 1$, $f^{0\pm}_{\pm} = \pm 2$, others zero or antisymmetrically determined. Dual Coxeter $h^\vee = 2$, so $\kappa(V_1) = \dim(\mathfrak{sl}_2)(k + h^\vee)/(2 h^\vee) = 3 \cdot 3 / 4 = 9/4$.

\emph{$m_3$ vanishes.} The first-order-pole associator $(J^a \circ J^b) \circ J^c - J^a \circ (J^b \circ J^c)$ vanishes by the $\mathfrak{sl}_2$-Jacobi identity
\[
  f^{ab}_d f^{dc}_e + f^{bc}_d f^{da}_e + f^{ca}_d f^{db}_e \;=\; 0.
\]
The second-order-pole contribution from the central term $k \kappa^{ab}$ gives a graded-commutative $m_3$ on the $\kappa$-pairing; its class in $\HH^2(V_1(\mathfrak{sl}_2))$ vanishes by $\mathrm{ad}$-invariance of the Killing form.

\emph{$B^{(2)}$.} HKR: $\HH^\bullet(V_k(\mathfrak{sl}_2)) \cong \mathrm{Sym}^\bullet(\mathfrak{sl}_2^\vee) \otimes \wedge^\bullet(\mathfrak{sl}_2)$. Cyclic shift on $\wedge^3(\mathfrak{sl}_2)$ is non-zero (since $\dim \mathfrak{sl}_2 = 3$), with explicit generator
\[
  B^{(2)}(J^+ \wedge J^- \wedge J^0) \;=\; J^+ \otimes J^- \otimes J^0 \;-\; J^0 \otimes J^+ \otimes J^- \;+\; J^- \otimes J^0 \otimes J^+.
\]

\emph{Commutator.}
\[
  [m_3, B^{(2)}]_{V_1(\mathfrak{sl}_2)} \;=\; \sum_{a, b, c} \epsilon^{abc} \cdot \kappa^{ab} \cdot J^c \cdot \mathrm{(mode\ data)}.
\]
Contraction of symmetric $\kappa^{ab}$ with antisymmetric $\epsilon^{abc}$:
\[
  \kappa^{ab} \epsilon^{abc} \;=\; \kappa^{+-} (\epsilon^{+-c} + \epsilon^{-+c}) \;+\; \kappa^{00} \epsilon^{00c} \;=\; 1 \cdot 0 \;+\; 2 \cdot 0 \;=\; 0.
\]
The cocycle vanishes.

\emph{Atiyah-class avatar.} Under Costello--Li's identification of $V_1(\mathfrak{sl}_2)$ with chiral differential operators on $BG$-moduli ($G = \mathrm{SL}_2$), the vanishing pulls back to $\at_{BG} \cup B^{(2)} = 0$: the tangent complex of $BG$ is $\mathfrak{g}[-1]$, a one-term complex, and $\Ext^1$ of a one-term complex with itself vanishes for dimensional reasons.

The toy analogues of the four bridge vertices are unobstructed: Theorem~B is Frenkel--Kac--Wakimoto bar-cobar inversion (the $\sfL$-row archetype), the Chern--Simons partition function is the $\mathrm{SL}_2$ character in the expected normalization, the $\bfH_{\Delta_5}$ vertex is represented only by the scalar rank-$3$ Mukai target in this toy comparison, and the trace identity $\hbar^2 \cdot (9/4) = -1$ gives $\hbar^2 = -4/9$.
\end{example}

\subsection{Five proof strands of Theorem~\ref{thm:main} traced on $K3 \times E$}
\label{sub:hand-strands}

\begin{example}[Strand-by-strand arithmetic closure on $K3 \times E$]
\label{ex:hand-strands}
Fix $Y = K3 \times E$ away from the Humbert divisors, after specifying a $\Zbb/n$ Borcea--Voisin twist and recomputing the Hodge data. Example~\ref{ex:hand-k3e} and Proposition~\ref{prop:k3e-vanish} reduce $[m_3, B^{(2)}]_Y = 0$ to the scalar condition $\lambda_{\mathrm{Cyc}}=0$. Each strand below is therefore a conditional witness for Theorem~\ref{thm:main}, not an unconditional proof of the pentagon.

\emph{Strand (i).} Reduced to the Kunneth target and scalar
$\lambda_{\mathrm{Cyc}}$ condition of Example~\ref{ex:hand-k3e}.

\emph{Strand (ii): bar-cobar inversion on generators.} $A_Y = \Phi_3(Y)$ on the reference curve $C$ is the Mukai-lattice vertex algebra $V_L$ for $L = \Gamma_{K3} \oplus \Gamma_E$ with $\Gamma_{K3} = E_8(-1)^{\oplus 2} \oplus U^{\oplus 3} \oplus U$ (signature $(4, 20)$) and $\Gamma_E = U$ (signature $(1, 1)$). Total signature $(5, 21)$. The bar complex is the Koszul resolution $\Bar(V_L) = \Lambda^\bullet(L_-) \otimes V_L$. On each rank-$1$ summand $\Zbb \cdot \ell \subset L$, $\Bar(V_\ell) = \Cbb[y_\ell]$ (polynomial) and $\Cobar(\Cbb[y_\ell]) = V_\ell$ (Heisenberg at rank $1$); tensoring over the $26$ summands of $L$ gives $\Cobar \Bar(V_L) = V_L$.

\emph{Strand (iii): parametrix uniqueness.} Costello--Li parametrix factorises $\mathrm{Param}_Y = \mathrm{Param}_{K3} \otimes \mathrm{Param}_E$ by K\"unneth. On $E$: heat kernel on the torus, unique up to contractible choice ($\tau$-translates of the standard kernel on $\Cbb/\Zbb^2$). On $K3$ with Calabi--Yau metric (Yau 1978): heat kernel of the unique Ricci-flat metric in its K\"ahler class. Product parametrix is unique up to contractible choice; $\pi_0(\mathrm{Param}_Y) = \{*\}$.

\emph{Strand (iv): $\bfH_{\Delta_5}$ source-recognition target.} Taking $\Sigma_2 = K3$ a primitive curve class of degree $2n - 2$, the restriction of $\Phi_3(Y)$ to $C$ along the K3 factor supplies the Mukai scalar data at rank $c_+(\Gamma_{K3}) = 4$, giving $K = 2 \cdot 4 = 8$ and $\hbar^2 = -1/K = -1/8$. The assertion that this scalar target is realised by $\bfH_{\Delta_5}$ on the elliptic factor is precisely the Vol.~III source-recognition hypothesis; without it the strand records only the $K=8$ Mukai target and the $\Delta_5$ denominator comparison.

\emph{Strand (v): trace identity arithmetic.} Conditional on Strand (i), the residual identity $\hbar^2 K = -1$ is the normalized conductor identity: $K = 2 c_+(\Gamma_{K3}) = 8$, $\hbar^2 = -1/8$, product $(-1/8) \cdot 8 = -1$.

\emph{Cycle closure (v) $\Rightarrow$ (i).} There is no ordinary coherent restriction return edge through $H^2(C,\Omega^1_C)$ (Lemma~\ref{lem:curve-h2-omega-vanishes}). Closure requires the derived specialisation map $\operatorname{sp}_{\Sigma,C}$ of Theorem~\ref{thm:sp-alpha-first-barcobar-obstruction}.

Under those bridge hypotheses the strands close into the conditional diagram of Theorem~\ref{thm:main}.
\end{example}

\subsection{Summary of the five witnesses}
\label{sub:hand-summary}

\begin{proposition}[Five hand-computable witnesses for the corrected bridge; \ClaimStatusConditional]
\label{prop:hand-summary}
Examples~\ref{ex:hand-quintic}--\ref{ex:hand-strands} collectively verify the elementary inputs and isolate the conditional bridges of Theorem~\ref{thm:main} on explicit generators.
\begin{center}
\begin{tabular}{l c c c l}
\toprule
Example & $[m_3, B^{(2)}] = 0$ & $\at_X = 0$ & $\hbar^2 K$ & Source of vanishing \\
\midrule
$Q_5$ (quintic) & unknown ($101$-dim target) & no & undefined & Yukawa detector; compute $\Psi_f(\alpha)$ \\
$K3 \times E$ & conditional & no ($\at_{K3} \neq 0$) & $-1$ & scalar $\lambda_{\mathrm{Cyc}}$ \\
$\cH_1$ (Heisenberg) & yes & n/a (alg.) & $-1$ & $m_3 = 0$ and $B^{(2)} = 0$ both \\
$V_1(\mathfrak{sl}_2)$ & yes & n/a (alg.) & $-4/9$ & $\epsilon^{abc} \kappa^{ab} = 0$ \\
$K3 \times E$ 5-strand & conditional & (projection) & $-1$ & bridge hypotheses close diagram \\
\bottomrule
\end{tabular}
\end{center}

Two algebraic mechanisms establish toy-model vanishing:
associativity-plus-rank-$1$ collapse and Killing-form-$\kappa^{ab}$
contracted with Levi-Civita-$\epsilon^{abc}$. The geometric examples
do different work: $Q_5$ supplies a $101$-dimensional falsifier target
with non-zero Yukawa detector, and $K3\times E$ exposes the non-zero
$20+1$ product target and scalar $\lambda_{\mathrm{Cyc}}$. The
five-strand example traces the conditional bridge, not an unconditional
pentagon.
\end{proposition}

\begin{proof}
Examples~\ref{ex:hand-quintic}, \ref{ex:hand-k3e}, \ref{ex:hand-heisenberg}, \ref{ex:hand-sl2}, \ref{ex:hand-strands}.
\end{proof}

\section{Open questions}
\label{sec:open}
% ==================================================

\subsection{Chain-level Clemens extension}

\begin{question}[Clemens chain-level extension, conditional on \textup{(C1)--(C3)}; \ClaimStatusOpen]
\label{thm:clemens-extension}
Can the cocycle $[m_3,B^{(2)}]$ be extended to Clemens-type
non-K\"ahler CY-$3$ spaces by a chain-level construction with
\textup{(C1)} chart-local replacement for the Kahler
$\partial\bar\partial$ argument, \textup{(C2)} boundary-term absence in
the Kotake--Kato product heat-kernel on chart overlaps, and
\textup{(C3)} a one-dimensional counterterm space for the relevant
twisted BV complex? A positive answer would require a Kuranishi cover
of the small-resolution locus, Costello--Li chart-local parametrices,
compatible overlap heat kernels, and a neutralized Courant/CDO
deformation complex.
\end{question}

The old Dolbeault proof does not provide the extension: its
representative-independence argument uses the Kahler
$\partial\bar\partial$-lemma. The question is whether the missing
global primitive is replaced by a Courant/CDO hypercohomology class.

The questions below are separated by epistemic status. Each names an answer whose attainability is either (a) \emph{conjecturally yes} --- a specific expected answer with a partial derivation, awaiting completion of a named construction in the literature --- or (b) \emph{unknown} --- no candidate answer is yet visible, the question posed as an invitation.

\subsection{All-loop vanishing (conjecturally yes)}

\begin{question}[All-loop cocycle; \ClaimStatusOpen, conjecturally yes]
\label{q:all-loop}
Does the higher-order analogue $[m_n, B^{(n-1)}]_X$ for $n \geq 4$ vanish on $\cU^{\mathrm{adm}}$? Equivalently, is the local BRST anomaly algebra generated by the Atiyah--Connes class after a BV/HKR diagonal comparison? The expected answer is yes, but Costello--Gwilliam/BV supplies the obstruction-recursion formalism, not a theorem that first-order vanishing bootstraps all higher obstructions. Under the compactness hypothesis of \S\ref{sub:fa-precision} the BV master equation $\{S, S\} + \hbar \Delta S = 0$ admits renormalised order-by-order analysis, and Gilkey 1995 Theorem~4.1.6 identifies finite-dimensional local counterterm spaces. The conditional reads precisely as:
\begin{itemize}
  \item \emph{conditional on} an all-loop BV/HKR diagonal comparison for 6d holomorphic Chern--Simons on a smooth projective Calabi--Yau threefold;
  \item \emph{valid on} the Koszul-admissible locus $\cU^{\mathrm{adm}}(X)$ where $\alpha_X=0$ secures only the base obstruction;
  \item \emph{requiring} the compactness hypothesis of \S\ref{sub:fa-precision} in its full force: cylindrical-end control is not yet shown to extend to all loop orders on non-compact complex-threefold targets such as $K3 \times \Cbb$, so the conditional-yes applies to compact CY-$3$ and awaits a non-compact refinement;
  \item \emph{stated} at the $(\infty,1)$-categorical level as the vanishing of the canonical class $[S_n] \in \HH^{6-n}(X, X)$ for each $n \geq 3$, with chain-level refinement as the vanishing of $\Omega_n^{\mathrm{BV}} \in H^{6-n}_{\mathrm{loc}}(X, \mathrm{Densities})$ at each order.
\end{itemize}
Conditional on the four bullets above, the expected answer is that all
higher obstructions are generated by the Atiyah--Connes class and hence
vanish on its zero locus. Unconditional resolution awaits the all-loop
generation theorem and either extension of the compactness hypothesis
to the relevant non-compact complex-threefold setting or independent
verification of the Clemens-stratum analogue through the $H$-twisted
Courant-algebroid CDO of \S\ref{sec:clemens}.
\end{question}

\subsection{Bridgeland finiteness unconditional (unknown)}

\begin{question}[Unconditional Bridgeland; \ClaimStatusOpen, unknown]
Can Corollary~\ref{cor:bridgeland} be made unconditional by replacing Bridgeland's conjecture with the Halpern-Leistner--Takeda approach to stability conditions on Calabi--Yau-$3$ categories? No candidate answer is visible. The Halpern-Leistner--Takeda framework covers stability conditions on slices of $D^b(X)$; global compactness of $\Stab(D^b(X))$ for CY-$3$ remains open. A positive answer would upgrade $\cU^{\mathrm{adm}}$-compactness from conditional to unconditional; a negative answer would isolate the load-bearing hypotheses in Bridgeland's conjecture.
\end{question}

\subsection{Twenty-second row (unknown)}

\begin{question}[$M_{23}$-row; \ClaimStatusOpen, unknown]
What is the $\sfB$-row chiral bialgebra attached to the $22$-row Niemeier BKM algebra whose underlying lattice is a non-Leech Niemeier of type $M_{23}$? No candidate chiral bialgebra is visible. Corollary~\ref{cor:m23} answers the question conditionally on Niemeier $1973$'s classification together with the Chenevier--Lannes $2019$ mass-formula refinement for the $M_{23}$-automorphism twist. An unconditional construction must bypass the non-reduced classification --- for example a direct orbifold of the Monster VOA by $M_{23}$, or a genus-$0$ presentation of $M_{23}$ acting on a rank-$22$ lattice VOA.
\end{question}

\subsection{$H$-twisted Courant-algebroid replacement on Clemens' stratum (conjecturally yes)}

\begin{question}[$H$-twisted cocycle on Clemens; \ClaimStatusOpen, conjecturally yes]
Does there exist a canonical class $[m_3, B^{(2)}]^H_{\tilde Y}$ in the deformation hypercohomology of a neutralized heterotic Courant/CDO vertex algebroid on a Clemens-stratum threefold $\tilde Y$, whose vanishing gates four $H$-twisted analogues of Theorems~\ref{thm:climax-I}--\ref{thm:climax-IV}? The expected answer is yes. Partial derivation: Corollary~\ref{cor:strominger} supplies the Strominger heterotic-with-torsion framework on balanced-Hermitian $\mathrm{SU}(3)$-structures, and Question~\ref{thm:clemens-extension} names the Kuranishi-plus-Kotake--Kato chain-level extension problem. Unifying these into a single $H$-twisted Courant-algebroid chiral differential operator algebra on the balanced-Hermitian $\mathrm{SU}(3)$-structure (Michelsohn 1982) is the expected path; the universal obstruction is the $H$-twisted analogue of $[m_3, B^{(2)}]$. What remains open is the construction of the Courant-algebroid chiral-differential-operator algebra itself --- specifically, whether the Strominger torsion can be absorbed into the cochain Connes rotation at bidegree two without further obstruction.
\end{question}

\medskip

The corrected theorem has one conditional obstruction class and four conditional bridges. Three geometries witness the regimes of the single class
\[
  [m_3, B^{(2)}]_X \;\in\; H^2(X, \Omega^1_X).
\]
Theorems~\ref{thm:climax-I}--\ref{thm:climax-IV} become comparable to this class only through the bridge hypotheses of Theorem~\ref{thm:main}. The cocycle is the obstruction; the frontier is the construction of the comparison maps.

% ==================================================
\appendix

% =====================================================================
% Explicit VOA content supplementing universal_anomaly_four_climax_simultaneously_2026_04_22.tex
% Five load-bearing computations: H_{Delta_5} comparison target, master trace on K3,
% Theorem B chain-level on H_1 and V_1(sl_2), DVV 1997 Sym^N(K3) derivation,
% nine Heegner arithmetic VOAs.
%
% Intended as an \input supplement to the main standalone; carries the
% same preamble assumptions (same LaTeX class, same \cH, \cO, \bfH macros).
%
% Author: Raeez Lorgat. 2026-04-23.
% =====================================================================

% ==================================================
\section{Explicit VOA content: five load-bearing computations}
\label{sec:voa-explicit}
% ==================================================

Each of the comparison vertices of Theorem~\ref{thm:main} carries a
concrete vertex-operator-algebra presentation with explicit OPEs,
explicit $\kappa$-invariants, and explicit Koszul-dual toy models on
the reference curve $C$. The computations below are evidence and input
for the conditional bridges; they do not by themselves prove the
five-vertex comparison conjecture.

% ==================================================
\subsection{$\bfH_{\Delta_5}$ comparison target: $\mathfrak g_{\Delta_5}$ generators and the Borcherds denominator $1/\Delta_5$}
\label{subsec:voa-H-Delta5-canonical}
% ==================================================

The K3 chiral bialgebra $\bfH_{\Delta_5}$ of
Theorem~\ref{thm:climax-III} is the conformal-block lift of the K3
Borcherds--Kac--Moody superalgebra $\mathfrak g_{\Delta_5}$ (Gritsenko--Nikulin 1998) via a lattice vertex
algebra on the Mukai lattice $L = H^\bullet_{\mathrm{even}}(K3,\Zbb)$
of signature $(4,20)$. The Frenkel--Kac lift transports every
BKM generator to a vertex-operator residue.

\paragraph{Lattice vertex algebra on the Mukai lattice.}
The Mukai lattice $L \cong \Zbb^{4,20}$ carries the positive-definite
$c_+(L) = 4$ summand generated by
$\beta^{(1)},\beta^{(2)},\beta^{(3)},\beta^{(4)}$ (class of a fixed
polarisation plus $H^0 \oplus H^4$ together with two transverse
$(1,1)$-classes), and the negative-definite $c_-(L) = 20$ summand
generated by transverse $(1,1)$ and $H^{2,0}+H^{0,2}$ directions. Let
$\phi^{(a)}(z)$, $a = 1,\ldots,24$, be $24$ free bosons of conformal
weight one realising the Mukai pairing:
\begin{equation}
\label{eq:voa-muk-ope}
  \phi^{(a)}(z)\,\phi^{(b)}(w) \;\sim\; g^{ab}_{\mathrm{Muk}}\log(z-w),
  \qquad g^{ab}_{\mathrm{Muk}} \text{ Mukai form of signature }(4,20).
\end{equation}
For $\alpha \in L$, the lattice vertex operator is
$V_\alpha(z) = {:}\exp(\alpha\cdot\phi(z)){:}$ with OPE (Borcherds 1986, \S5)
\begin{equation}
\label{eq:voa-lattice-vo-ope}
  V_\alpha(z)\,V_\beta(w)
    \;=\;
    \epsilon(\alpha,\beta)\,(z-w)^{\langle\alpha,\beta\rangle}
    \,{:}V_\alpha(z)V_\beta(w){:} \;+\; (\text{regular}),
\end{equation}
with $\epsilon$ the Borcherds sign cocycle. The Fock module
$\cF_{K3} = \mathrm{Sym}\bigl(\bigoplus_{n\geq 1}L\otimes t^{-n}\bigr)
\otimes \Cbb[L]$ is the universal Heisenberg module on which every
vertex operator $V_\alpha$ acts.

\paragraph{Real and imaginary simple roots of $\mathfrak g_{\Delta_5}$.}
The BKM superalgebra $\mathfrak g_{\Delta_5}$ has a rank-three
Lorentzian primitive Cartan lattice, not an even unimodular lattice of
signature $(2,1)$; even unimodular indefinite lattices require
signature difference divisible by $8$. In the normalization used here
its real simple roots $\delta_1,\delta_2,\delta_3$ have Gram matrix
\begin{equation}
\label{eq:voa-BKM-Gram}
  G_{\mathrm{BKM}} \;=\; \begin{pmatrix} 2 & -2 & -2 \\ -2 & 2 & -2 \\ -2 & -2 & 2 \end{pmatrix}
  \;=\; 2\,\mathrm{diag}(1,1,1) \;-\; 2(J - I), \qquad \det G_{\mathrm{BKM}} = -32,
\end{equation}
and imaginary simple roots $\alpha^{\mathrm{im}} = (n,\ell,m) \in \Lambda_{\Delta_5}$
(with $4nm - \ell^2 \geq 0$ a Humbert discriminant invariant) of
multiplicity
\begin{equation}
\label{eq:voa-imag-mult}
  \mathrm{mult}(\alpha^{\mathrm{im}}) \;=\; c_{K3}(4nm - \ell^2),
\end{equation}
with $c_{K3}(D)$ the Fourier coefficient of $D$-th order in the K3
elliptic genus generator
$\phi^{K3}_{0,1}(\tau,z) = 2\phi_{0,1} - (E_2/6)\phi_{-2,1}$
(Eichler--Zagier 1985; Dabholkar--Murthy--Zagier 2012).

\paragraph{Real-root Chevalley generators as Heisenberg vertex operators.}
For each real simple root $\delta_i$, the Chevalley generators
$e_i,f_i,h_i$ are realised as bare Heisenberg currents on the
negative-definite sub-lattice directions of the Mukai lattice:
\begin{equation}
\label{eq:voa-real-root-generators}
  e_i(z) \;=\; V_{\delta_i}(z) \;=\; {:}e^{\delta_i\cdot\phi(z)}{:}, \qquad
  f_i(z) \;=\; V_{-\delta_i}(z), \qquad
  h_i(z) \;=\; \delta_i\cdot\partial\phi(z).
\end{equation}
The Chevalley commutator $[e_i,f_j] = \delta_{ij}h_i^\vee$ follows from
the $(z-w)^{\langle\delta_i,-\delta_j\rangle} = (z-w)^{-2}$ double pole
of the lattice-vertex-operator OPE, with the residue extracting the
Cartan generator. The Serre relations at $\delta_i,\delta_j$ with
$\langle\delta_i,\delta_j\rangle = -2$ follow from the two-term Serre
inclusion $\mathrm{ad}(e_i)^{1 - a_{ij}}(e_j) = \mathrm{ad}(e_i)^3(e_j) = 0$
witnessed by the fourth-order pole vanishing in the triple OPE
$V_{\delta_i}V_{\delta_i}V_{\delta_i}V_{\delta_j}$.

\paragraph{Imaginary-root Chevalley generators as screening operators.}
For each imaginary simple root $\alpha^{\mathrm{im}}$, the Chevalley
generator $e^{\mathrm{im}}_\alpha$ is a Feigin--Frenkel screening
operator
\begin{equation}
\label{eq:voa-imag-root-screening}
  S_{\alpha^{\mathrm{im}}}(w) \;:=\; \oint_{\gamma_w}\,V_{\alpha^{\mathrm{im}}}(z)\,\mathrm{d}z
  \;=\; \oint_{\gamma_w}\,{:}e^{\alpha^{\mathrm{im}}\cdot\phi(z)}{:}\,\mathrm{d}z,
\end{equation}
with $\gamma_w$ a small loop around $w$ chosen so that
$V_{\alpha^{\mathrm{im}}}(z)$ has conformal weight one. The OPE
$V_{\alpha^{\mathrm{im}}}(z)V_{\beta^{\mathrm{im}}}(w)$ has pole order
$-\langle\alpha^{\mathrm{im}},\beta^{\mathrm{im}}\rangle$; when
$\langle\alpha,\beta\rangle \geq 0$ the OPE is regular and the screening
commutator
$[S_{\alpha^{\mathrm{im}}}, S_{\beta^{\mathrm{im}}}] = 0$ in agreement
with Borcherds axiom (GKM4) (the imaginary-root vanishing axiom).
When $\langle\alpha,\beta\rangle = -1$, the bracket extracts a single
residue giving $e^{\mathrm{im}}_{\alpha+\beta}$ with coefficient
$c_{K3}(4(nm' + n'm) - 2\ell\ell')$. The structure constants
$c_{\alpha,\beta}^j$ of the imaginary-root bracket
\begin{equation}
\label{eq:voa-imag-bracket}
  [e^{\mathrm{im}}_\alpha, e^{\mathrm{im}}_\beta] \;=\; \begin{cases}
    0, & \langle\alpha,\beta\rangle \geq 0, \\
    \sum_j c_{\alpha,\beta}^{\;j}\,e^{\mathrm{im}}_{\alpha+\beta - \gamma_j}, & \langle\alpha,\beta\rangle < 0,
  \end{cases}
\end{equation}
are determined by the K3 elliptic-genus Fourier coefficients at the
pair weight, via Borcherds' singular theta lift (Borcherds 1998,
\S6).

\paragraph{Weight-one primary condition and ghost cancellation.}
For $V_{\alpha^{\mathrm{im}}}$ to contribute a weight-one primary
(so that the contour integral is well-defined and BRST-closed), the
conformal weight
$h(V_{\alpha^{\mathrm{im}}}) = \tfrac{1}{2}\langle\alpha^{\mathrm{im}},\alpha^{\mathrm{im}}\rangle$
must be corrected by a transverse $c = 23$ SCFT sector and a single
$bc$-ghost pair at conformal weight $(\tfrac{1}{2},\tfrac{1}{2})$,
producing total ghost central charge $c_{\mathrm{gh}} = -26$. The
ghost-dressed imaginary-root vertex operator
\begin{equation}
\label{eq:voa-ghost-dressed}
  \tilde V_{\alpha^{\mathrm{im}}}(z) \;=\; V_{\alpha^{\mathrm{im}}}(z) \;\otimes\; c(z),
\end{equation}
after integration against $\mathrm{d}z$, lands in BRST cohomology
$H^\bullet(Q_{\mathrm{BRST}})$ and descends to
$e^{\mathrm{im}}_\alpha \in \mathfrak g_{\Delta_5}$.

\paragraph{The Borcherds denominator $1/\Delta_5$.}
The Weyl--Kac--Borcherds denominator function of $\mathfrak g_{\Delta_5}$
is (Gritsenko--Nikulin 1998, Theorem 2.1)
\begin{equation}
\label{eq:voa-GN-denom}
  \Delta_5(Z) \;=\; \sum_{(n,\ell,m) > 0}\,c_{\phi_{-2,1}}(4nm - \ell^2)\,q^n r^\ell s^m,
  \qquad q = e^{2\pi i\tau},\;r = e^{2\pi i z},\;s = e^{2\pi i\tau'},\; Z = (\tau,z,\tau'),
\end{equation}
the Gritsenko cusp form of weight $5$ on the Siegel upper half-space
$\Hbb_2$ with character $\chi_5$. The character is a square root of
the Igusa cusp form: $\Phi_{10}(Z) = \Delta_5(Z)^2$ (Gritsenko--Nikulin 1998,
Eq.~4.6). The Weyl--Kac identity reads
\begin{equation}
\label{eq:voa-WKB-denom}
  \Delta_5(Z) \;=\; e^{2\pi i\langle \rho,\, Z\rangle}\,
  \prod_{\alpha \in \Delta_+(\mathfrak g_{\Delta_5})}
  \bigl(1 - e^{2\pi i\langle\alpha,\,Z\rangle}\bigr)^{\mathrm{mult}(\alpha)},
\end{equation}
with $\rho = \tfrac{1}{2}(\delta_1 + \delta_2 + \delta_3) \in
\Lambda_{\Delta_5}\otimes\Qbb$ the Weyl vector of $\mathfrak g_{\Delta_5}$
at pairing $\langle\rho,\delta_i\rangle = -1$ for each real simple
root (half-sum of positive real roots, Gritsenko--Nikulin 1998
\S2, equation 2.5), and $\mathrm{mult}(\alpha)$ given by
$c_{K3}(\tfrac{1}{2}\langle\alpha,\alpha\rangle)$ (real root:
$\mathrm{mult} = 1$; imaginary root:
$\mathrm{mult} = c_{K3}(D)$ at discriminant $D$). The lowest-order
imaginary-root Fourier coefficients populating \eqref{eq:voa-imag-mult}
are $c_{K3}(0) = 2$, $c_{K3}(1) = 90$, $c_{K3}(2) = 462$,
$c_{K3}(3) = 1540$, $c_{K3}(4) = 4554$ (K3 elliptic genus
Eichler--Zagier 1985 Table 9). The generating function $1/\Delta_5$
is thus the partition function of the universal
$\mathfrak g_{\Delta_5}$-module $\bfH_{\Delta_5}$ in the natural
grading:
\begin{equation}
\label{eq:voa-partition-H-Delta5}
  \mathrm{ch}\,\bfH_{\Delta_5}(Z) \;=\; \frac{1}{\Delta_5(Z)}
  \;=\; e^{-2\pi i\langle\rho,\,Z\rangle}\,
  \prod_{\alpha\in\Delta_+}\,\bigl(1 - e^{2\pi i\langle\alpha,\,Z\rangle}\bigr)^{-\mathrm{mult}(\alpha)}.
\end{equation}
This is the direct Fock-space counting of states of $\bfH_{\Delta_5}$,
weighted by their $\Lambda_{\Delta_5}$ charge.

\paragraph{Two $\kappa$-invariants on $\bfH_{\Delta_5}$: chiral versus BKM.}
The K3 chiral bialgebra $\bfH_{\Delta_5}$ carries two distinct
$\kappa$-invariants, which must not be conflated. The chiral
$\kappa_{\mathrm{ch}}(\bfH_{\Delta_5}) = 3$ equals the positive
rank of the K3 transcendental lattice minus the Mukai-doubling
defect:
$\kappa_{\mathrm{ch}} = \rho(K3) - 1 + \chi(\cO_{K3})/2 = 1 + 2 = 3$
for a polarised K3 at Picard rank one (the programme-canonical K3 of
Vol III Theorem~CY-A$_3$). The BKM-weight invariant
$\kappa_{\mathrm{BKM}}(\Phi_1) = c_1(0)/2 = 5$ equals the weight of
the Gritsenko--Nikulin form $\Delta_5$ at the universal
eight-form-spread index $N = 1$, with $c_1(0) = 10$ the Fourier
zero-coefficient. The two are distinct mathematical objects living
on different complexes: $\kappa_{\mathrm{ch}}$ is the chiral
homology shadow computed from the bar complex of $A_{K3\times E}$,
$\kappa_{\mathrm{BKM}}$ is the Borcherds singular-theta lift weight
computed from the K3 elliptic-genus Jacobi form
$\phi^{K3}_{-2,1}$. The additive split
$\kappa_{\mathrm{BKM}} = \kappa_{\mathrm{ch}} + \chi(\cO_{\mathrm{fiber}})$
is false: the fiber $\cO_E$ has $\chi = 0$, giving $3 + 0 = 3 \neq 5$.
The correct relation is the master universal Borcherds weight
identity $\kappa_{\mathrm{BKM}}(\Phi_N) = c_N(0)/2$ (Borcherds 1995
singular-theta lift, Gritsenko 1999 explicit denominator formula).

% ==================================================
\subsection{Universal trace $\hbar^2 K = -1$ on K3: Bruinier 2002 Proposition 5.1 and Mukai signature}
\label{subsec:voa-hbar-K-universal-trace}
% ==================================================

The master trace identity $\hbar^2 K = -1$ of
Theorem~\ref{thm:climax-IV} on $X = K3$ reduces to the
representation-theoretic identification $K = 2c_+(L) = 8$ via
Bruinier 2002 Proposition 5.1 and the Serre-duality normalisation
of the vacuum trace on the derived centre.

\paragraph{Bruinier Heegner--Chern class reciprocity (Prop.~5.1).}
For an even lattice $L$ of signature $(2, \ell)$ and a Heegner
divisor $H_\beta(m) \subset \cF_L = \mathrm{O}^+(L)\backslash\cH_L$
(with $\cH_L$ the period domain, $\beta \in L^\vee/L$ the coset, and
$m < 0$ the norm), Bruinier 2002 Proposition 5.1 gives the Chern-class
identity
\begin{equation}
\label{eq:voa-bruinier-heegner-chern}
  c_1\bigl(\cL_{\Psi}\bigr) \;=\; \frac{k}{2}\,c_1(\mathcal L)
  \;-\; \sum_{\substack{\beta \in L^\vee/L \\ m < 0}}\,c_\Psi(\beta, m)\,[H_\beta(m)],
\end{equation}
with $\cL$ the automorphic line bundle on $\cF_L$, $\Psi$ a Borcherds
product of weight $k = c_\Psi(0, 0)/2$, and $c_\Psi(\beta, m)$ the
principal-part Fourier coefficients of the input weakly-holomorphic
modular form $F_\Psi$. The quantity $c_+(L)$ is the positive-signature
rank of $L$, so $c_+(\mathrm{Muk}(K3)) = 4$ and $K = 2c_+(L) = 8$ in
every enhancement compatible with the Mukai lattice.

\paragraph{Mukai signature and the $K = 8$ identification.}
The Mukai lattice of K3, $L = H^\bullet_{\mathrm{even}}(K3, \Zbb) =
H^0 \oplus H^2 \oplus H^4$ with Mukai pairing
$\langle\alpha,\beta\rangle_{\mathrm{Muk}} = \int_{K3}\alpha\cup\iota(\beta)$
(with $\iota$ the Mukai involution), has signature $(4, 20)$: four
positive-definite classes (degree-zero and degree-four classes $H^0
\oplus H^4 = \Zbb^2$ together with $\mathrm{Pic}(K3) \cap
(H^{1,1}(K3,\Rbb))^+$ of rank 2 at the Mukai-lattice basepoint) and
twenty negative-definite classes (the transcendental lattice of K3
together with the remaining positive $(1,1)$-classes). The
Mukai-doubling identification of Vol III (Theorem~7.3 of
\texttt{k3\_chiral\_bialgebra\_platonic.tex} in Vol III) gives
$c_+(L) = 4$, hence
\begin{equation}
\label{eq:voa-K-8-Mukai}
  K \;=\; 2\,c_+(L) \;=\; 8, \qquad \hbar^2 \;=\; -1/K \;=\; -1/8, \qquad \hbar^2 K \;=\; -1.
\end{equation}

\paragraph{Serre-duality normalisation of the vacuum trace.}
The Calabi--Yau-$3$ structure on $K3 \times E$ (Kontsevich--Vlassopoulos
2022, Theorem 1.2) endows the derived centre $Z^{\mathrm{der}}_{\mathrm{ch}}(A_{K3\times E})$
with a trace pairing $\tr : Z^{\mathrm{der}}_{\mathrm{ch}} \to \Cbb[-3]$
satisfying $\tr \circ \mathbb S_{D^b(X)} = -\tr$: the Serre functor
inverts the trace. On the vacuum module $\Omega_A \in A\text{-mod}$,
Caldararu--Willerton 2010 Theorem 1.6 motivates the comparison
\begin{equation}
\label{eq:voa-trace-vacuum}
  \tr(1_A) \;=\; \hbar^2\cdot K,
\end{equation}
with the equality read as a normalized conductor identity. It is not
derived from the Todd integral of a trivial rank-$K$ object on
$K3\times E$: that Euler characteristic is zero. The
$\mathbb S$-anti-self-duality fixes the sign convention for the trace;
the value $\hbar^2K=-1$ is the imposed Mukai--Bruinier conductor
normalisation.

\paragraph{Anomaly off the Koszul locus.}
Off the Koszul locus, Gerstenhaber's cyclic-pairing obstruction
produces the anomalous identity
\begin{equation}
\label{eq:voa-universal-trace-anomaly}
  \hbar^2 K \;+\; \int_X\,\at_X\cup B^{(2)} \;=\; -1,
\end{equation}
with the correction term given by the cocycle of
Construction~\ref{constr:cocycle}. On $\cU^{\mathrm{adm}}(K3)$ the
integral vanishes and the identity reduces to $\hbar^2 K = -1$
exactly; on the generic quintic $Q_5$ the integral contributes an
irreducible $101$-dimensional shift and the master identity acquires
an anomaly incompatible with $K \in \Zbb$.

\paragraph{Explicit target of $\at_X \cup B^{(2)}$ on K3 and $K3\times E$.}
On $X = K3$ the vanishing of $\int_X\,\at_X\cup B^{(2)}$ is a
direct consequence of the Mukai-lattice Atiyah computation
(Kapranov 1999 Prop.~4.4, Remark): the tangent bundle $T_{K3}$
has $\at_{K3} \in \Ext^1_{K3}(T_{K3}, T_{K3}\otimes\Omega^1_{K3})$,
and on a K3 (Calabi--Yau surface) the Dolbeault representative is
$\bar\partial\nabla \in \Omega^{0,1}(K3, \End(T_{K3})\otimes\Omega^1_{K3})$
for a Chern connection. The cup product $\at_{K3}\cup B^{(2)}$
with Connes' bidegree-two cyclic shift takes values in
$H^2(K3, \Omega^1_{K3})=H^{1,2}(K3)$, hence vanishes by dimension.
This proves only the K3-only receptacle vanishing. For $K3\times E$,
the group $H^2(K3\times E,\Omega^1)$ is nonzero of dimension $21$;
vanishing there must be a computation of the specific projected class,
not vanishing of the receptacle.

% ==================================================
\subsection{Theorem B chiral Positselski: explicit bar-cobar on $\cH_1$ and $V_1(\mathfrak{sl}_2)$}
\label{subsec:voa-theorem-B-bar-cobar}
% ==================================================

The canonical non-trivial witnesses of Theorem B on the reference
curve $C$ are the rank-one Heisenberg $\cH_1$ (abelian OPE, simplest
possible) and the level-one affine $V_1(\mathfrak{sl}_2)$ (non-abelian
OPE, simplest possible with simple-pole structure).

\paragraph{The Heisenberg $\cH_1$: abelian OPE, trivial bar differential.}
Single chiral generator $\alpha(z)$ of conformal weight one, OPE
\begin{equation}
\label{eq:voa-heis-ope}
  \alpha(z)\,\alpha(w) \;=\; \frac{1}{(z-w)^2} \;+\; \mathrm{regular}, \qquad
  [\alpha_m, \alpha_n] \;=\; m\,\delta_{m+n, 0}\cdot\mathbf{1},
\end{equation}
central charge $c = 1$, $\kappa(\cH_1) = 1$. Abelian: double pole only,
no simple pole. The augmentation ideal
$\bar\cH_1 = \bigoplus_{n \geq 1}\Cbb\cdot{:}\partial^{k_1}\alpha\cdots\partial^{k_n}\alpha{:}$
generates the ordered bar complex
$\Bar^{\mathrm{ch}}_C(\cH_1) = T^c_C(s^{-1}\bar\cH_1)$ with
Maurer--Cartan differential driven by pair residues.

\emph{Bar differential on pairs.} The Borcherds identity
(Borcherds 1986; see $\S$\ref{sub:B2-operad}) gives
\begin{equation}
\label{eq:voa-heis-bar-diff}
  d_{\mathrm{res}}\bigl[s^{-1}\alpha(z)\,|\,s^{-1}\alpha(w)\bigr]
  \;=\; \mathrm{Res}_{z = w}\bigl((z-w)\cdot (z-w)^{-2}\bigr)\cdot\mathbf{1}
  \;=\; 0.
\end{equation}
The simple-pole coefficient $\alpha_{(0)}\alpha = 0$ is the Heisenberg
complementarity conductor $K_{\mathrm{Heis}} = 0$. Higher bar
differentials concatenate residues of the vanishing simple pole and
vanish identically.

\emph{Cobar differential and Positselski contracting homotopy.}
$\Cobar^{\mathrm{ch}}_C(\Bar^{\mathrm{ch}}_C(\cH_1)) = T(s\overline{\Bar^{\mathrm{ch}}(\cH_1)})$
carries the cobar cocomposition plus the (trivial) inherited bar
differential. Positselski's chain homotopy
\begin{equation}
\label{eq:voa-heis-homotopy}
  h\,:\, \Cobar\,\Bar \;\longrightarrow\; \Cobar\,\Bar[-1], \qquad
  h\bigl(s^{-1}x \otimes s\,y\bigr) \;=\; x\cdot y
\end{equation}
is the tautological shift (no curvature correction because the bar
differential is null). One computes
$[d, h] = d\circ h + h\circ d = \id - p_0$, with $p_0$ the projection
onto the weight-zero (vacuum) line. Positselski's Theorem 6.5 gives
the chain-level equivalence
\begin{equation}
\label{eq:voa-heis-inversion}
  \epsilon : \Cobar^{\mathrm{ch}}_C\,\Bar^{\mathrm{ch}}_C(\cH_1) \;\xrightarrow{\sim}\; \cH_1, \qquad
  [x_1|\cdots|x_n] \;\longmapsto\; {:}x_1\cdots x_n{:},
\end{equation}
the counit being normal-ordered reconstruction. All obstruction
classes $\mathrm{obs}_n \in H^n(\cH_1, \cH_1 \otimes \Omega^1_C)$
vanish for $n \geq 2$; $\cH_1$ is Koszul on the nose.

\emph{Explicit weight-two counit on Fock states.} The weight-two
Fock states of $\cH_1$ are
$\alpha_{-1}^2\Omega$ and $\alpha_{-2}\Omega$ with inner products
$\langle\alpha_{-1}^2\Omega|\alpha_{-1}^2\Omega\rangle = 2$,
$\langle\alpha_{-2}\Omega|\alpha_{-2}\Omega\rangle = 2$,
$\langle\alpha_{-1}^2\Omega|\alpha_{-2}\Omega\rangle = 0$. The counit
\eqref{eq:voa-heis-inversion} at weight two reads
\[
  \epsilon\bigl([s^{-1}\alpha \otimes s^{-1}\alpha]\bigr) \;=\; {:}\alpha\alpha{:}(0)\Omega \;=\; \alpha_{-1}^2\Omega,
  \qquad \epsilon\bigl([s^{-1}\partial\alpha]\bigr) \;=\; \partial\alpha(0)\Omega \;=\; \alpha_{-2}\Omega,
\]
with no cross-contamination between the two modes (orthogonality of
$\alpha_{-1}^2$ and $\alpha_{-2}$ on the vacuum). The Connes cyclic
shift $B^{(2)}$ at bidegree two acts on the cyclic tensor
$[s^{-1}\alpha | s^{-1}\alpha]$ as the $\Zbb/2$-antisymmetrisation,
which vanishes by $\alpha\otimes\alpha = \alpha\otimes\alpha$
(abelian commutation). Hence $B^{(2)}$ annihilates every weight-two
cyclic bar chain on $\cH_1$, and the cocycle $[m_3, B^{(2)}]_{\cH_1}
= 0$ on the nose in all homological degrees.

\paragraph{The level-one affine $V_1(\mathfrak{sl}_2)$: Frenkel--Kac realisation.}
Currents $J^a(z)$ for $a \in \{+,-,0\}$, level $k = 1$, OPE
\begin{equation}
\label{eq:voa-V1sl2-ope}
  J^a(z)\,J^b(w) \;=\; \frac{\delta^{ab}}{(z-w)^2} \;+\; \frac{i\,\varepsilon^{abc}\,J^c(w)}{z - w}
  \;+\; \mathrm{regular},
\end{equation}
central charge $c = \dim(\mathfrak{sl}_2)\cdot k/(k + h^\vee) =
3\cdot 1/(1+2) = 1$, $\kappa(V_1(\mathfrak{sl}_2)) =
\dim(\mathfrak{sl}_2)(k + h^\vee)/(2 h^\vee) = 3(1+2)/4 = 9/4$.

\emph{Frenkel--Kac boson realisation.} Frenkel--Kac 1980 identifies
$V_1(\mathfrak{sl}_2) \cong V_{A_1}$, the rank-one lattice vertex
algebra on the root lattice $A_1 = \Zbb\alpha$ with $\langle\alpha,\alpha\rangle = 2$,
through the map
\begin{equation}
\label{eq:voa-FK-realisation}
  J^0(z) \;=\; \frac{1}{\sqrt{2}}\,\partial\phi(z), \qquad
  J^{\pm}(z) \;=\; {:}\exp\bigl(\pm\sqrt{2}\,\phi(z)\bigr){:} \cdot \epsilon^\pm,
\end{equation}
with $\phi$ a free boson, $[\phi_m, \phi_n] = m\delta_{m+n,0}$, and
$\epsilon^\pm$ cocycle factors resolving the $\mathbb Z/2$-gradation
ambiguity. The lattice-vertex-operator OPE (Kac 1998, \S5.5) gives
\begin{equation}
\label{eq:voa-FK-ope-check}
  V_\alpha(z)\,V_{-\alpha}(w) \;=\; (z-w)^{-2}\cdot\mathbf{1} + (z-w)^{-1}\,\sqrt{2}\,\partial\phi(w) + \mathrm{reg.},
\end{equation}
and comparison with \eqref{eq:voa-V1sl2-ope} identifies
$V_\alpha = J^+$, $V_{-\alpha} = J^-$, $\sqrt{2}\,\partial\phi = 2\,J^0$.

\emph{Bar differential: Kac--Moody bracket on desuspensions.}
Pair-level bar differential
\begin{equation}
\label{eq:voa-V1sl2-bar-diff}
  d_{\mathrm{res}}\bigl[s^{-1}J^a\,|\,s^{-1}J^b\bigr]
  \;=\; \mathrm{Res}_{z = w}\Bigl(\frac{i\,\varepsilon^{abc}\,J^c(w)}{z - w}\Bigr)\,
  \mathrm{d}(z - w) \;=\; i\,\varepsilon^{abc}\cdot s^{-1}J^c,
\end{equation}
extracting the Kac--Moody bracket, plus double-pole Koszul
corrections from the $\delta^{ab}$ central term.

\emph{Bar cohomology.} By Positselski 2011 \S6.5 applied to
$\hat{\mathfrak{sl}}_2$ at level one, the bar cohomology computes
\begin{equation}
\label{eq:voa-V1sl2-bar-cohomology}
  H^\bullet\,\Bar^{\mathrm{ch}}_C(V_1(\mathfrak{sl}_2))
  \;=\; \mathrm{Sym}^\bullet(s^{-1}\mathfrak{sl}_2)\otimes\Lambda^\bullet(s^{-1}\Omega^1_C),
\end{equation}
the Koszul resolution of the universal enveloping chiral algebra.
The obstruction tower collapses at $n = 2$ via the Jacobi identity:
$\mathrm{obs}_2(J^a, J^b, J^c) = J^a[J^b, J^c] + \text{cyc.} = 0$.
$V_1(\mathfrak{sl}_2)$ sits in the Koszul stratum.

\emph{Counit and Wakimoto resolution.} The Positselski counit
\begin{equation}
\label{eq:voa-V1sl2-counit}
  \epsilon : \Cobar^{\mathrm{ch}}_C(\Bar^{\mathrm{ch}}_C(V_1(\mathfrak{sl}_2))) \;\longrightarrow\; V_1(\mathfrak{sl}_2),
  \qquad [J^{a_1}|\cdots|J^{a_n}] \;\longmapsto\; {:}J^{a_1}(z)\cdots J^{a_n}(z){:}
\end{equation}
is a quasi-isomorphism in $D^{\mathrm{co}}_{\mathrm{ch}}(C)$; the
inverse-modulo-homotopy is witnessed by the Wakimoto resolution
(Frenkel--Ben-Zvi 2004, \S15.4): every $V_1(\mathfrak{sl}_2)$-module
$M$ admits a resolution
\begin{equation}
\label{eq:voa-wakimoto}
  0 \;\longrightarrow\; M \;\longrightarrow\; W_0 \;\xrightarrow{d_0}\; W_1 \;\xrightarrow{d_1}\; \cdots
\end{equation}
by semi-infinite Fock modules $W_j$ over the sub-Heisenberg
$\Cbb\cdot\partial\phi(z)$; each $W_j$ is Koszul on the nose
(as in the $\cH_1$ computation above); the hypercohomology spectral
sequence converges to $M$ with trivial $E_2$-page after the first
differential. The cocycle $[m_3, B^{(2)}]$ restricts to zero on
$V_1(\mathfrak{sl}_2)$ because $\End^\bullet(V_1(\mathfrak{sl}_2))$ is
strictly associative: $m_3 \equiv 0$ as an operadic chain,
Massey-bracket obstructions vanish per level-one-triviality of the
$A_\infty$-lift.

\paragraph{Joint conclusion.}
Both $\cH_1$ and $V_1(\mathfrak{sl}_2)$ witness Theorem B chain-level
with explicit contracting homotopy; $[m_3, B^{(2)}]$ vanishes on the
reference-curve shadow of each; the comparison targets specialise to
identities on Fock and lattice-vertex states.

% ==================================================
\subsection{3D gravity $1/\Phi_{10}$: DVV 1997 derivation via symmetric-orbifold VOA on $\mathrm{Sym}^N(K3)$}
\label{subsec:voa-dvv-symorbifold}
% ==================================================

The Dijkgraaf--Moore--Verlinde--Verlinde 1997 derivation of
$1/\Phi_{10}$ as a generating function of K3 BPS microstates proceeds
through the symmetric-orbifold vertex operator algebra.

\paragraph{Symmetric-orbifold VOA on $\mathrm{Sym}^N(K3)$.}
For the K3 sigma-model superconformal field theory with
$(c,\bar c) = (6,6)$ central charges, let $\cV_{K3}$ denote the
corresponding vertex operator super-algebra with elliptic-genus
character
\begin{equation}
\label{eq:voa-K3-ell-genus}
  \chi_{\mathrm{ell}}(K3; \tau, z) \;=\; 8\,\Bigl[
    \Bigl(\frac{\theta_2(\tau,z)}{\theta_2(\tau,0)}\Bigr)^2 +
    \Bigl(\frac{\theta_3(\tau,z)}{\theta_3(\tau,0)}\Bigr)^2 +
    \Bigl(\frac{\theta_4(\tau,z)}{\theta_4(\tau,0)}\Bigr)^2
  \Bigr] \;=\; 2\phi_{0,1}(\tau,z),
\end{equation}
the weak Jacobi form of weight $0$ and index $1$ (Eichler--Zagier 1985,
\S9). The $N$-fold symmetric orbifold $\mathrm{Sym}^N(K3) = K3^N/S_N$
carries the symmetric-orbifold VOA
\begin{equation}
\label{eq:voa-sym-orbifold}
  \cV_{\mathrm{Sym}^N(K3)} \;=\; \bigl(\cV_{K3}^{\otimes N}\bigr)^{S_N} \;\oplus\; \bigoplus_{[\sigma]}\cV_{K3}^{(\sigma)},
\end{equation}
with $[\sigma]$ running over non-identity conjugacy classes of $S_N$
and $\cV_{K3}^{(\sigma)}$ the twisted sector associated to each
cycle-type partition of $N$. The untwisted sector is the $S_N$-invariant
tensor power; the twisted sectors contribute one copy of $\cV_{K3}$
per cycle in the partition, compactified on the long cycle.

\emph{Twisted-sector generators at small $N$.} For $N = 2$, the
non-trivial conjugacy class is the $2$-cycle, and the twisted sector
$\cV_{K3}^{(12)}$ is a single copy of $\cV_{K3}$ compactified on a
length-$2$ cycle, i.e.\ states live on the double cover
$\tilde S^1 \to S^1$ with period $4\pi$ on the double cover versus
$2\pi$ on the base. The twisted-sector ground state is a rank-one
module over the $\Zbb/2$-invariant subalgebra
$\cV_{K3}^{S_2} = \cV_{K3}\otimes\cV_{K3}$ with twisted vertex operators
\[
  V_{(12), \alpha\otimes\beta}(z) \;=\; \frac{1}{\sqrt{2}}\bigl(V_\alpha(z)\cdot V_\beta(-z) + V_\beta(z)\cdot V_\alpha(-z)\bigr),
\]
for $\alpha,\beta \in L_{K3}$, the Lehn--Sorger 2003 twisted vertex
operators on the Hilbert scheme. For $N = 3$, the conjugacy classes are
$3$-cycles (one class) and pairs of transpositions (absent in $S_3$);
the $3$-cycle twisted sector is $\cV_{K3}$ on the triple cover, with
ground state partition function $q^{c/24}$ fraction at $c_{K3} = 6$
giving $q^{1/4}$ shift per the orbifold-GSO sum.

\paragraph{DMVV product formula.}
Dijkgraaf--Moore--Verlinde--Verlinde 1997 Theorem 5.1 computes the
symmetric-orbifold elliptic-genus generating function:
\begin{equation}
\label{eq:voa-dmvv}
  \sum_{N \geq 0}\,p^N\,\chi_{\mathrm{ell}}(\mathrm{Sym}^N(K3); \tau, z)
  \;=\; \prod_{n \geq 1, m \geq 0, \ell \in \Zbb}
  \frac{1}{(1 - p^n q^m r^\ell)^{c(nm, \ell)}},
\end{equation}
with $c(D, \ell)$ the Fourier coefficient of
$\chi_{\mathrm{ell}}(K3; \tau, z) = \sum_{m,\ell}c(m,\ell)q^m r^\ell$.
Under the multiplicative lift (Borcherds 1995; Gritsenko 1999), the
right-hand side equals the Gritsenko lift of $\phi^{K3}_{0,1}$:
\begin{equation}
\label{eq:voa-gritsenko-lift}
  \prod_{n,m,\ell}(1 - p^n q^m r^\ell)^{c(nm,\ell)}
  \;=\; p^{-1}q^{-1}r^{-1}\cdot\Phi_{10}(Z)
  \quad \text{up to the Gritsenko leading term,}
\end{equation}
with $\Phi_{10}$ the Igusa cusp form of weight $10$ on the Siegel upper
half-space (Igusa 1962; Gritsenko--Nikulin 1998, \S2). Combining
\eqref{eq:voa-dmvv} and \eqref{eq:voa-gritsenko-lift},
\begin{equation}
\label{eq:voa-dvv-1-over-Phi10}
  \sum_{N \geq 0}\,p^N\,\chi_{\mathrm{ell}}(\mathrm{Sym}^N(K3); \tau, z)
  \;=\; \frac{p q r}{\Phi_{10}(Z)}
  \;\propto\; \frac{1}{\Phi_{10}(Z)}.
\end{equation}

\paragraph{Identification with the 3D gravity partition function.}
Maldacena--Moore--Strominger 1999 identify the left-hand side of
\eqref{eq:voa-dvv-1-over-Phi10} with the $1/4$-BPS microstate-counting
partition function of type IIB string theory on $K3 \times S^1$ at
the attractor point. The Strominger--Vafa 1996 AdS$_3$/CFT$_2$
duality then identifies this CFT partition function with the
twisted-holography boundary character of 3D quantum gravity on
$\mathrm{AdS}_3 \times K3$:
\begin{equation}
\label{eq:voa-3dqg-partition}
  Z^{\AdS_3 \times K3}_{\mathrm{3dQG, BPS}}(Z) \;=\;
  \sum_{N \geq 0}\,p^N\,\chi_{\mathrm{ell}}(\mathrm{Sym}^N(K3); \tau, z)
  \;\propto\; \frac{1}{\Phi_{10}(Z)}.
\end{equation}
The boundary CFT is the $\mathrm{Sym}^N(K3)$ orbifold sigma-model
(Seiberg--Witten 1999; Dijkgraaf 1999 review), whose VOA lives on the
asymptotic $S^1 \times \Rbb$ boundary of $\mathrm{AdS}_3$; the
elliptic-genus insertion computes the BPS trace. The $3$D gravity
identification is at the twisted-holography scope: the dynamical
metric path integral is conjectural (Maloney--Witten 2009
sum convergence plus Maxfield--Turiaci 2020 wormhole corrections),
but the automorphic function $1/\Phi_{10}$ is rigorously pinned by
\eqref{eq:voa-dvv-1-over-Phi10} and the identification
$\Phi_{10} = \Delta_5^2$ of Gritsenko--Nikulin 1998.

% ==================================================
\subsection{Candidate arithmetic VOAs $V^{\mathrm{arith}}_{K3, d_K}$ at the conditional Heegner discriminants}
\label{subsec:voa-heegner-arithmetic}
% ==================================================

Each of the nine Heegner discriminants
$d_K \in \{3, 4, 7, 8, 11, 15, 19, 20, 24\}$ of
Proposition~\ref{prop:heegner} is treated here as a conditional
candidate for an arithmetic VOA
$V^{\mathrm{arith}}_{K3, d_K}$ with central charge $K = 8$. The
canonical construction is not proved in this supplement; it requires
the rank-$8$ Bruinier coefficient computation and an explicit
class-group action on the Mukai-lattice VOA.

\paragraph{Imaginary-quadratic orders and Heegner CM points.}
For each $d_K$ above, let $K_{d_K} = \Qbb(\sqrt{-d_K})$ be the
imaginary-quadratic field of discriminant $-d_K$ and
$\cO_{d_K}$ the corresponding quadratic order. The values
$d_K=8,20,24$ are fundamental negative discriminants in the usual
field-discriminant convention here; they are not used as
conductor-$2$ replacements for $d_K=4,5,6$. The Heegner point
$\tau_{d_K} \in
\mathfrak H$ on the upper half-plane is the CM point of discriminant
$-d_K$ where the elliptic curve $E_{\tau_{d_K}} = \Cbb/\langle 1,
\tau_{d_K}\rangle$ has complex multiplication by $\cO_{d_K}$. Via
the Shioda--Inose 1977 correspondence, each ideal class of
discriminant $-d_K$ determines a singular K3 surface $K3_{d_K}$ with
transcendental lattice
$T_{d_K}$ of rank 2, positive-definite, discriminant $-d_K$:
\begin{equation}
\label{eq:voa-T-dK}
  T_{d_K} \;\cong\; \begin{cases}
    \Zbb\bigl\langle 1, \tfrac{1 + \sqrt{-d_K}}{2}\bigr\rangle, & d_K \equiv 3 \pmod 4, \\[4pt]
    \Zbb\bigl\langle 1, \tfrac{\sqrt{-d_K}}{2}\bigr\rangle, & d_K \equiv 0 \pmod 4,
  \end{cases}
\end{equation}
with Gram matrix
$G_{T_{d_K}} = \mathrm{diag}(2, (1 + d_K)/2)$ at odd $d_K$
and $\mathrm{diag}(2, d_K/2)$ at even $d_K$
(Shioda--Inose 1977, Proposition 4.3).

\paragraph{Scheithauer generalised-moonshine lift.}
At each $\tau_{d_K}$, the expected construction is a twisted VOA
$V^{\mathrm{arith}}_{K3, d_K}$ obtained from the Mukai-lattice VOA
$V_{\mathrm{Muk}(K3)}$ by a Scheithauer-type
generalised-moonshine lift and the class-group action of the order
$\cO_{d_K}$:
\begin{equation}
\label{eq:voa-Vari-K3-dK}
  V^{\mathrm{arith}}_{K3, d_K} \;:=\;
  \bigl(V_{\mathrm{Muk}(K3)}\bigr)^{\sigma_{d_K}}
  \otimes \Cbb[\mathrm{Cl}(\cO_{d_K})],
\end{equation}
where $\mathrm{Cl}(\cO_{d_K})$ is the ideal class group. The quotient
$\cO_{d_K}/\cO_{d_K}^*$ is not the class group and is not used.

\paragraph{Explicit generators at the class-number-one discriminants in the conditional list.}
For $d_K \in \{3, 4, 7, 8, 11, 19\}$ with class number $h_{d_K} = 1$,
the class group $\mathrm{Cl}(\cO_{d_K})$ is trivial and
$V^{\mathrm{arith}}_{K3, d_K} = (V_{\mathrm{Muk}(K3)})^{\sigma_{d_K}}$
has no non-trivial class-group twist. At each such $d_K$, the transcendental-lattice
generator $\theta_{d_K}(z) = V_{(1, \tau_{d_K})}(z)$ realises the
CM data:
\begin{itemize}
\item $d_K = 3$: $\theta_3(z) = V_{(1, \omega_3)}(z) = {:}e^{\omega_3\cdot\phi(z)}{:}$ with
$\omega_3 = (1 + \sqrt{-3})/2$; $\mathcal O_3 = \Zbb[\omega_3]$ Eisenstein integers.
\item $d_K = 4$: $\theta_4(z) = V_{(1, i)}(z)$ with $i = \sqrt{-1}$;
$\mathcal O_4 = \Zbb[i]$ Gaussian integers.
\item $d_K = 7$: $\theta_7(z) = V_{(1, \omega_7)}(z)$, $\omega_7 = (1 + \sqrt{-7})/2$;
$\mathcal O_7 = \Zbb[\omega_7]$.
\item $d_K = 8$: $\theta_8(z) = V_{(1, \sqrt{-2})}(z)$; $\mathcal O_8 = \Zbb[\sqrt{-2}]$.
\item $d_K = 11$: $\theta_{11}(z) = V_{(1, \omega_{11})}(z)$, $\omega_{11} = (1 + \sqrt{-11})/2$;
$\mathcal O_{11} = \Zbb[\omega_{11}]$.
\item $d_K = 19$: $\theta_{19}(z) = V_{(1, \omega_{19})}(z)$, $\omega_{19} = (1 + \sqrt{-19})/2$;
$\mathcal O_{19} = \Zbb[\omega_{19}]$.
\end{itemize}
Each $\theta_{d_K}$ has conformal weight
$h(\theta_{d_K}) = \tfrac{1}{2}\langle (1,\tau_{d_K}), (1,\tau_{d_K})\rangle$
computed from \eqref{eq:voa-T-dK}; the Scheithauer twist then
tensors with a rank-$8$ lattice-vertex-algebra shift landing
$V^{\mathrm{arith}}_{K3, d_K}$ at central charge $K = 8$.

The central-charge check at the fixed points: the $\sigma_{d_K}$-invariant
subalgebra of $V_{\mathrm{Muk}(K3)}$ has central charge equal to the
signature-positive rank $c_+(L) = 4$ (Mukai-doubling Vol III
\texttt{k3\_chiral\_bialgebra\_platonic.tex}); the Scheithauer
arithmetic twist doubles the central charge to $K = 2c_+(L) = 8$,
and the ghost dressing at level-one Heisenberg shifts by
$c_{\mathrm{gh}} = -26$ then back by $+26$ ($bc$ and $\beta\gamma$
complementary cancellation), netting the $K = 8$ invariant across
the nine strata.

\paragraph{Class-number-two discriminants in the conditional list.}
For $d_K \in \{15, 20, 24\}$ with $h_{d_K} = 2$, the class group
$\mathrm{Cl}(\cO_{d_K}) \cong \Zbb/2\Zbb$ is generated by a single
non-principal ideal class, and
$V^{\mathrm{arith}}_{K3, d_K} = (V_{\mathrm{Muk}(K3)})^{\sigma_{d_K}}
\otimes \Cbb[\Zbb/2]$ carries a candidate $\Zbb/2$-orbifold structure.
The assertion that exactly these class-number-two discriminants survive
a rank-$8$ Bruinier positivity cut is conditional; it is not a
consequence of Bruinier 2002 Proposition 5.1 alone. The other
class-number-two discriminants $d_K \in \{23, 31, 35, 39, 40, 43, 47, 51, 52, 55,
56, 68, 72, 84, 88, 100, 115, 120, 148, 168, 187, 195, 228, 232,
280, 292, 312, 340, 372, 403, 408, 427, 520, 532, 555, 595, 627,
708, 715, 760, 795, 1012, 1435\}$ produce Heegner divisors of wrong
signature for the rank-$8$ Mukai enhancement only after the missing
coefficient/sign calculation. At each admissible
$d_K$, $V^{\mathrm{arith}}_{K3, d_K}$ carries generators:
\begin{itemize}
\item $d_K = 15$: $\theta_{15}, \theta_{15}^\sigma$ under the non-principal
ideal $\mathfrak p_3 \cdot \mathfrak p_5$; $\mathcal O_{15} = \Zbb[(1+\sqrt{-15})/2]$.
\item $d_K = 20$: $\theta_{20}, \theta_{20}^\sigma$;
$\mathcal O_{20} = \Zbb[\sqrt{-5}]$.
\item $d_K = 24$: $\theta_{24}, \theta_{24}^\sigma$;
$\mathcal O_{24} = \Zbb[\sqrt{-6}]$.
\end{itemize}
The Borcherds--Kac--Moody algebra attached to
$V^{\mathrm{arith}}_{K3, d_K}$ is
$\mathfrak g^{\mathrm{arith}}_{d_K} := \mathfrak g_{\Delta_5}$ twisted
by $\sigma_{d_K}$ (Gritsenko--Nikulin 1998 \S7). Each
$V^{\mathrm{arith}}_{K3, d_K}$ has central charge $K = 8$ (Mukai
signature); the arithmetic twist alters the conformal weight
assignments of lattice vertex operators but leaves the central charge
and the $\kappa = 0$ Heisenberg conductor invariant.

\paragraph{Connection to the $\sfB$-row.}
Conditional on the arithmetic calculation in Proposition~\ref{prop:heegner} and on constructing the class-group action above, the nine $V^{\mathrm{arith}}_{K3, d_K}$ stratify the $\sfB$-row of the
derived-centre classification by arithmetic twist: each witnesses
the ceiling $K^\kappa = 8$ of Theorem~C with a different
imaginary-quadratic-field arithmetic. The CM point $\tau_{d_K}$ of
Heegner discriminant $-d_K$ sits on the Humbert divisor of
discriminant $d_K$ in $\cF_L$; the vanishing of
$[m_3, B^{(2)}]$ at $d_K$ is a scalar cyclic-HKR condition on the
corresponding Humbert-Heegner CM stratum, after embedding into the
principal $K3\times E$ CY3 specialization. Off the conditional list,
the expected obstruction is the missing rank-$8$ arithmetic
admissibility condition.

% ==================================================
\subsection{Koszul-dual algebras on the reference curve for the five vertices}
\label{subsec:voa-koszul-duals}
% ==================================================

Each of the comparison vertices of Theorem~\ref{thm:main} is expected to carry an
explicit Koszul-dual algebra $A^!$ on the reference curve $C$,
computed from $A^i = H^\star\Bar^{\mathrm{ch}}_C(A)$ via Verdier
duality $A^! = (A^i)^\vee$. The five Koszul duals assemble only after
the bridge hypotheses of Theorem~\ref{thm:main} are supplied:

\begin{itemize}
\item \emph{Vertex~$V_1 = \{[m_3, B^{(2)}]_X = 0\}$.} Koszul-dual of the
vanishing cocycle is the graded-commutative algebra
$\mathrm{Sym}^\bullet(H^{1,2}(X)^\vee[-1])$ on the dual of the
corrected obstruction target; trivial when $[m_3, B^{(2)}] = 0$ forces
$H^{1,2}(X)^\vee$ to act nilpotently on the chiral algebra
$A_X = \Phi_3(X)$ (via the Atiyah-class pairing).

\item \emph{Vertex~$V_2 = \mathrm{Thm\,B}$.} Koszul-dual of the chiral
Positselski inversion is the universal enveloping chiral coalgebra
$\Bar^{\mathrm{ch}}_C(A)$ itself on the cofree side; on the admissible
Koszul locus $A^! = \Omega^{\mathrm{ch}}_C(A^i)$ is the inversion
functor output. For $A = \cH_1$, $A^! = \Omega_C^\bullet$ the de Rham
complex on $C$ (trivial coalgebra on one generator). For
$A = V_1(\mathfrak{sl}_2)$, $A^! = \mathrm{CE}^\bullet(\hat{\mathfrak{sl}}_2)$
the Chevalley--Eilenberg complex computing the Lie-algebra cohomology.

\item \emph{Vertex~$V_3 = \mathrm{Stage\text{-}1\ canonicality}$.}
Koszul-dual of the canonicality of the stage-one
$E_3$-factorisation algebra $\Phi^{\mathrm{FA}}_3(X)$ is the stage-two
specialisation $\Sp^{\mathrm{ch}}_{\Sigma_{d-1}, C}$ as a
Koszul-functor; its Koszul dual on $C$ is the universal curve-side
deformation complex, which for $X = K3$ is the Hilbert-scheme chiral
algebra $\mathrm{Hilb}^n(K3)$-chiral-algebra shadow on $C$.

\item \emph{Vertex~$V_4 = \bfH_{\Delta_5}$.} Koszul-dual of
$\bfH_{\Delta_5}$ is the Feigin--Frenkel dual
$\bfH_{\Delta_5}^! = V_{\mathrm{Muk}(K3)}^\vee \otimes \Lambda^\bullet(\mathfrak g^{\mathrm{im}}_{\Delta_5})$,
the Koszul-dual lattice vertex algebra tensored with the exterior
algebra on the imaginary simple roots. Central charge of the dual:
$K^! = -K + 2\dim(\mathfrak h) = -8 + 2\cdot 3 = -2$; the
complementarity $K + K^! = 8 - 2 = 6$ is the Vol III $\sfB$-row
six-form-spread witness (off the principal $K + K^! = 8$ ceiling by
the rank-three Cartan shift).

\item \emph{Vertex~$V_5 = \{\hbar^2 K = -1\}$.} Koszul-dual of the
universal trace identity is the Hochschild cochain complex
$\HH^\bullet(A_X, A_X)$ computing the derived centre
$Z^{\mathrm{der}}_{\mathrm{ch}}(A_X)$; its Koszul dual on $C$ is the
curve-side evaluation of the Serre trace pairing
$\tr : Z^{\mathrm{der}}_{\mathrm{ch}}(A_X)|_C \to \Cbb[-3]$, landing
in $\Cbb[-3]$ for CY-3 and pinned at $\tr(1) = -1$ by Serre duality.
\end{itemize}

The resulting Verdier-dual decagon is a frontier conjecture: under
\textup{(B1)--(B4)} and the additional construction of these Koszul
duals, it should refine Conjecture~\ref{conj:atiyah-connes-pentagon}.

% =====================================================================
% End of VOA supplement
% =====================================================================


% ==================================================
% Supplement: Local-global + arithmetic refinements for the
% universal anomaly comparison. Five refinement cycles inscribed:
% (1) Atiyah scope separation; (2) Koszul locus chart vs moduli;
% (3) Two loci intersection via Kuga--Satake period map;
% (4) Nine Heegner d_K as conditional rank-8 arithmetic problem;
% (5) Motivic Galois scope of Bridgeland finiteness.
% ==================================================

\section{Local-global scope and arithmetic refinements}
\label{sec:local-global-arithmetic}

Five refinements sharpen the scope of the conditional bridge. Each names a chart-local/global dichotomy or an arithmetic-theoretic refinement that the main statements assume implicitly.

\subsection{Atiyah scope separation: four global statements that admit no chart-local avatar}
\label{subsec:atiyah-scope-separation}

The four statements
\begin{itemize}
  \item the Hodge location of Corollary~\ref{cor:hodge}, namely $[m_3, B^{(2)}]_X \in H^2(X,\Omega^1_X)=H^{1,2}(X)$;
  \item the Yukawa non-vanishing $\int_{Q_5} \at_{Q_5}^{\cup 3} \neq 0$ of Proposition~\ref{prop:quintic-nonvanish};
  \item the K\"unneth additivity $\at_{K3 \times E} = \at_{K3} \boxplus \at_E$ of Proposition~\ref{prop:k3e-vanish};
  \item the vanishing $\alpha_X=0 \in H^2(X, \Omega^1_X)$ defining the Atiyah--Connes zero-locus candidate of Definition~\ref{def:atiyah-connes-zero-locus}
\end{itemize}
each live strictly at \emph{global} scope on $X$. On any single affine chart $U_\alpha \subset X$ trivialising $T_X$ --- witnessed by a chart-local Chern connection $\nabla_\alpha$ --- the restriction $\at_\cE|_{U_\alpha} \in \Ext^1_{U_\alpha}(\cE|_{U_\alpha}, \cE|_{U_\alpha} \otimes \Omega^1_{U_\alpha})$ vanishes identically: $\nabla_\alpha$ is a primitive, and every chart-local formula written in local coordinates with a named $\nabla_\alpha$ is tautologous at the level of $\at$. The four global statements are not shadowed by chart-local computations; they are born from patching combinatorics.

Conversely, the chain-level constructions of §\ref{sec:cocycle} and Theorem~\ref{thm:clemens-extension} are intrinsically chart-local. The Massey primitive $\alpha$ with $d\alpha = f \circ g$ exists chart-locally; the Kotake--Kato product heat-kernel is chart-local; the Costello--Li parametrix pinning is chart-local on each Kuranishi chart. The content of each vanishing theorem is the obstruction to patching these chart-local primitives into one global primitive. Local vanishing is automatic. Global vanishing is the theorem.

\subsection{Koszul locus: chart-local on $X$ versus global in $\cM_{\mathrm{cx}}(X)$}
\label{subsec:koszul-scope}

The locus $\cU^{\mathrm{adm}}(X) \subset \cM_{\mathrm{cx}}(X)$ is a subscheme of the Kuranishi moduli space of complex structures on the underlying smooth manifold, not a subset of $X$. Two equivalent global definitions serve. Globally on $\cM_{\mathrm{cx}}(X)$, the relative cocycle
\[
  [m_3, B^{(2)}]_{X/\cM_{\mathrm{cx}}(X)} \;\in\; R^2 (\pi_\cM)_* \Omega^1_{X_{\mathrm{univ}}/\cM_{\mathrm{cx}}(X)}
\]
on the universal family $\pi_\cM : X_{\mathrm{univ}} \to \cM_{\mathrm{cx}}(X)$ defines a coherent sheaf. The vanishing locus of the corresponding section is the obstruction-theoretic candidate for $\cU^{\mathrm{adm}}(X)$; any openness statement requires a separate smoothness or constant-rank hypothesis. Equivalently, $\cU^{\mathrm{adm}}(X)$ parametrises $[J] \in \cM_{\mathrm{cx}}(X)$ for which the Positselski contracting homotopy $h$ with $[d, h] = \id - p$ exists at every order of the completed modular cyclic obstruction filtration. Theorem~\ref{thm:main} makes the equivalence of these two definitions conditional on the bridge map \textup{(B1)}.

A Kuranishi chart $V_\beta \subset \cM_{\mathrm{cx}}(X)$ is \emph{not} a chart $U_\alpha \subset X$: the former is an open subset of moduli space parametrising a Kodaira--Spencer slice at some $[J_\beta]$; the latter is an open subset of the threefold. Two propagations must not be conflated: chart-overlap Čech patching on $U_\alpha \cap U_{\alpha'} \subset X$ of §\ref{subsec:atiyah-local-global} (gluing local connections into $\at_X$), and chart-overlap gluing on $V_\beta \cap V_{\beta'} \subset \cM_{\mathrm{cx}}$ (gluing Kuranishi obstruction data into $\cU^{\mathrm{adm}}$). Openness of $\cU^{\mathrm{adm}}$ is the Kodaira--Spencer compatibility of chart-uniform Koszulness witnesses across moduli-chart overlaps.

The three witnesses of §\ref{sec:clemens} live on three distinct moduli components: $\cM_{\mathrm{cx}}(Q_5)$ is $101$-dimensional and is a falsifier for the unconditional pentagon (Corollary~\ref{cor:quintic-fails}); $\cM_{\mathrm{cx}}(K3 \times E)$ is $21$-dimensional, but its admissible locus is conditional on the scalar computation $\lambda_{\mathrm{Cyc}}=0$ and on any Borcea--Voisin recomputation (Corollary~\ref{cor:k3e-locus}); the Clemens stratum lies in a non-K\"ahler component on which Dolbeault cohomology is ill-defined, and $\cU^{\mathrm{adm}}$ is given by the Fr\"olicher-degenerate reformulation of Theorem~\ref{thm:clemens-extension}.

\subsection{Two Koszul loci intersection: Atiyah side versus Humbert side via Kuga--Satake}
\label{subsec:two-loci-kuga-satake}

Theorem~\ref{thm:two-loci} intersects two a-priori-unrelated Koszul loci. The \emph{Atiyah-cocycle locus} $\cU^{\mathrm{adm}}_{\at}(X) \subset \cM_{\mathrm{cx}}(X)$ is the moduli-global vanishing locus of $[m_3, B^{(2)}]_{X/\cM_{\mathrm{cx}}(X)}$ of §\ref{subsec:koszul-scope}, an Atiyah-theoretic condition on $X$ itself. The \emph{Humbert-period locus} $\cP^{-1}(U^{\mathrm{adm}}) \subset \cM_{\mathrm{cx}}(X)$ is the preimage of the Kuga--Satake period image of the Bruinier admissible locus, a period-theoretic condition on $\overline{\cM_{\mathrm{A_2}}} = \Gamma \backslash D$ (the arithmetic quotient of the period domain, a Shimura variety).

The two a-priori-unrelated loci intersect in the stage-one canonical locus: this is the content of Theorem~\ref{thm:two-loci}. The period map $\cP$ reads Hodge-theoretic data; the Atiyah cocycle reads local Kuranishi-obstruction data. Their intersection equals $\cU^{\mathrm{canonical}}_{\Phi_3}$ by stage-one canonicality plus Griffiths transversality.

The Kuga--Satake period map (Kuga--Satake 1967, Morrison 1984 Theorem~7.9) is the embedding $\mathrm{KS} : D \to \cA_g^{\mathrm{KS}}$ lifting the K3 period domain $D = \mathrm{O}^+(2, 20)/(\mathrm{SO}(2) \times \mathrm{O}(20))$ to the Siegel modular variety $\cA_g^{\mathrm{KS}}$ of principally polarised abelian varieties of dimension $g = 2^{19}$. The lift is Hodge-compatible: the $(2,0)$-part of a K3 becomes the holomorphic $1$-forms of its Kuga--Satake abelian variety. Humbert divisors on $\cA_g^{\mathrm{KS}}$ classify abelian varieties whose endomorphism ring contains a real quadratic order; their preimages under $\mathrm{KS}$ and $\cP$ form the loci on which Bruinier's positivity cut imposes a definite signature at the Mukai-enhanced central charge $K = 8$. The passage through Kuga--Satake is necessary: the same positivity condition on the K3 side involves the transcendental lattice, which is directly readable on the abelian variety side via endomorphism structure. The factorisation $\cP = \mathrm{KS} \circ \cP_{K3}$ for K3-fibred $X$ identifies $\cU^{\mathrm{canonical}}_{\Phi_3}(X)$ as the Atiyah condition \emph{intersected with} the period image of the Bruinier locus: two independent conditions, the second Kuga--Satake-arithmetic.

\subsection{Nine Heegner discriminants: conditional rank-$8$ arithmetic problem}
\label{subsec:heegner-chenevier}

Proposition~\ref{prop:heegner} retains
$d_K \in \{3, 4, 7, 8, 11, 15, 19, 20, 24\}$ as a conditional
rank-$8$ arithmetic candidate. As currently justified, this set is the
initial segment $d\leq24$ of fundamental negative discriminants with
class number at most $2$, not a proved Bruinier--Chenevier selection.
The classical class-number-one fundamental discriminants remain
\[
  \{3,4,7,8,11,19,43,67,163\}.
\]

Chenevier's determinants and Bruinier's Heegner-divisor Chern-class
formula are relevant technologies, but they do not by themselves force
$K=8$, $\hbar^2=-1/8$, or the nine-element list. To make the refinement
theorem-grade one must define an explicit rank-$8$ admissibility
predicate $\mathcal C_8(d)$, compute the Bruinier coefficient/sign in
that predicate, and prove that its solution set is exactly the stated
nine discriminants.

The nine discriminants $\{3, 4, 7, 8, 11, 15, 19, 20, 24\}$ should not
be justified by the unit-group condition: every imaginary quadratic
order contains $\{\pm1\}$, while only $d_K=3,4$ have exceptional extra
units. Nor should Chenevier determinants be cited as selecting the list
without a residual pseudodeformation datum and a comparison map.

Thus the arithmetic refinement is conditional: either the
rank-$8$ predicate and Bruinier coefficient calculation must be
supplied, or the statement should revert to the classical Heegner list
with a separate positivity problem.

\subsection{Motivic Galois scope of Bridgeland finiteness}
\label{subsec:motivic-galois-scope}

Corollary~\ref{cor:bridgeland}'s conditional Bridgeland finiteness sits inside a larger arithmetic scope governed by the motivic Galois group. The \emph{motivic Galois group} $G_{\mathrm{mot}}(X)$ of a Calabi--Yau threefold $X$ is the Tannakian dual of the category of mixed motives over $X$ (Deligne--Milne 1982, Andr\'e 1996). Conjecturally, $G_{\mathrm{mot}}(X)$ is a reductive linear algebraic group acting on $\bigoplus_n H^n_{\mathrm{mot}}(X) \cong \bigoplus_n H^n_{\mathrm{dR}}(X)$ by restrictions of the Mumford--Tate group of the complex Hodge structure.

The Mumford--Tate--Galois--Sato--Lannes principle, in this manuscript's
working nomenclature, asserts that for a Calabi--Yau threefold
admitting a K3 fibration, the motivic Galois group of the Kuga--Satake
lift factors through the orthogonal group of the Mukai lattice:
$G_{\mathrm{mot}}(\mathrm{KS}(X)) \subseteq \mathrm{O}(L_{K3,\Qbb})$.
This is a programme-level conjectural input, not a
Chenevier--Lannes theorem.

The Bridgeland finiteness statement of Corollary~\ref{cor:bridgeland} is conditional on Bridgeland's conjecture that $\Stab(D^b(X))$ is a finite-dimensional complex manifold with locally-closed period maps. The arithmetic refinement: if MGSL holds for $X$, then the Koszul-admissible locus $\cU^{\mathrm{adm}}(X) \cap \cM_{\mathrm{cx}}(X)$ decomposes into $G_{\mathrm{mot}}(X)$-orbits under the motivic Galois action on the period domain; each orbit is a Shimura subvariety, and Bridgeland finiteness on each orbit follows from the arithmetic-volume finiteness of Shimura subvarieties (Zhang 2014). The motivic-Galois conditional statement is therefore strictly stronger than the Bridgeland conditional: MGSL $\Rightarrow$ Bridgeland on $\cU^{\mathrm{adm}}(X)$ but not conversely.

Two scope-separations are load-bearing. First, the motivic Galois scope applies only to $X$ in the K3-fibered regime where Kuga--Satake lifts exist; Clemens' non-K\"ahler stratum has no motivic Galois group in the conjectural sense (the mixed motives of a non-K\"ahler threefold lack the Hodge structure required for the Tannakian construction). Second, the Bridgeland scope applies more broadly but is weaker: it gives finiteness on compact subsets of $\Stab(D^b(X))$ but not global decomposition into Shimura subvarieties. Theorem~\ref{thm:main} holds in both scopes; its refinements via Heegner discriminants (Proposition~\ref{prop:heegner}) and two-loci intersection (Theorem~\ref{thm:two-loci}) use Kuga--Satake and hence live in the motivic-Galois scope.

\begin{remark}[Bibliographic anchors]
\label{rem:arithmetic-biblio}
The arithmetic refinements invoke: Chenevier 2014 (\emph{Astérisque} 324, determinants of Galois representations), Deligne--Milne 1982 (Tannakian categories), Andr\'e 1996 (motivic Galois group), Chenevier--Lannes 2019 (mass formula for Niemeier lattices with $M_{23}$ twist), Morrison 1984 Theorem 7.9 (Kuga--Satake construction), Morrison 1984 Theorem 8.1 ($K3 \times E$ motivic Galois), Zhang 2014 (arithmetic volume of Shimura subvarieties). None enter the proof of Theorem~\ref{thm:main}; they refine the scope of the corollaries (Theorem~\ref{thm:two-loci}, Proposition~\ref{prop:heegner}, Corollary~\ref{cor:bridgeland}, Corollary~\ref{cor:m23}) into arithmetic-precise statements.
\end{remark}


% ==================================================
\section*{Acknowledgements}
% ==================================================
The author thanks his co-advisors at the Perimeter Institute for extensive comments on the programme within which this obstruction theorem emerged, and the Vol~I / Vol~II / Vol~III manuscripts within which the comparison targets were developed.

% ==================================================
\begin{thebibliography}{99}
% ==================================================

\bibitem{Ayala-Francis-2015} D. Ayala, J. Francis, \emph{Factorization homology of topological manifolds}, J.~Topol. 8 (2015), 1045--1084.

\bibitem{AFT-2017} D. Ayala, J. Francis, H. L. Tanaka, \emph{Factorization homology of stratified spaces}, Selecta Math. 23 (2017), 293--362. Theorem 4.3.1.

\bibitem{Costello-2013} K. Costello, \emph{Supersymmetric gauge theory and the Yangian}, arXiv:1303.2632, 2013.

\bibitem{BBDJS-2015} C. Brav, V. Bussi, D. Dupont, D. Joyce, B. Szendr\H{o}i, \emph{Symmetries and stabilization for sheaves of vanishing cycles}, J.~Singul. 11 (2015), 85--151.

\bibitem{Beilinson-1978} A. Beilinson, \emph{Coherent sheaves on $\Pbb^n$ and problems of linear algebra}, Funct.~Anal.~Appl. 12 (1978), 214--216.

\bibitem{Besse-1987} A. Besse, \emph{Einstein manifolds}, Springer, 1987.

\bibitem{Bondal-Kapranov-1989} A. Bondal, M. Kapranov, \emph{Representable functors, Serre functors, and mutations}, Math.~USSR-Izv. 35 (1990), 519--541.

\bibitem{Bondal-Orlov-2001} A. Bondal, D. Orlov, \emph{Reconstruction of a variety from the derived category and groups of autoequivalences}, Compositio~Math. 125 (2001), 327--344. Theorem 3.1.

\bibitem{Bondal-Van-den-Bergh-2003} A. Bondal, M. Van den Bergh, \emph{Generators and representability of functors in commutative and noncommutative geometry}, Moscow Math.~J. 3 (2003), 1--36. Theorem 3.1.1.

\bibitem{Cohn-2016} L. Cohn, \emph{Differential graded categories are $k$-linear stable infinity-categories}, arXiv:1308.2587, 2016. \S3.

\bibitem{Borcea-1997} C. Borcea, \emph{$K3$ surfaces with involution and mirror pairs of Calabi--Yau manifolds}, in: \emph{Mirror symmetry, II}, AMS/IP Stud.~Adv.~Math., 1997.

\bibitem{Borcherds-1998} R. Borcherds, \emph{Automorphic forms with singularities on Grassmannians}, Invent.~Math. 132 (1998), 491--562.

\bibitem{Bruinier-2002} J. Bruinier, \emph{Borcherds products on $\mathrm{O}(2, l)$ and Chern classes of Heegner divisors}, Lecture Notes in Math. 1780, Springer, 2002. Proposition 5.1.

\bibitem{Caldararu-2005} A. C\u{a}ld\u{a}raru, \emph{The Mukai pairing, II: the Hochschild--Kostant--Rosenberg isomorphism}, Adv.~Math. 194 (2005), 34--66.

\bibitem{Caldararu-Willerton-2010} A. C\u{a}ld\u{a}raru, S. Willerton, \emph{The Mukai pairing, I: a categorical approach}, New York J.~Math. 16 (2010), 61--98. Theorem 1.6.

\bibitem{Calaque-Van-den-Bergh-2010} D. Calaque, M. Van den Bergh, \emph{Hochschild cohomology and Atiyah classes}, Adv.~Math. 224 (2010), 1839--1889.

\bibitem{CDOGP-1991} P. Candelas, X. de la Ossa, P. Green, L. Parkes, \emph{A pair of Calabi--Yau manifolds as an exactly soluble superconformal field theory}, Nucl.~Phys.~B 359 (1991), 21--74.

\bibitem{CG-FA-2} K. Costello, O. Gwilliam, \emph{Factorization algebras in quantum field theory, Volume 2}, Cambridge Univ.\ Press, 2021.

\bibitem{CGP-2018} K. Costello, O. Gwilliam, \emph{Factorization algebras in quantum field theory, Volume 1}, Cambridge Univ.\ Press, 2016.

\bibitem{Chenevier-Lannes-2019} G. Chenevier, J. Lannes, \emph{Automorphic forms and even unimodular lattices: Kneser neighbours of Niemeier lattices}, Ergebnisse der Mathematik 69, Springer, 2019.

\bibitem{Niemeier-1973} H.-V. Niemeier, \emph{Definite quadratische Formen der Dimension 24 und Diskriminante 1}, J.~Number Theory 5 (1973), 142--178.

\bibitem{Clemens-1983} H. Clemens, \emph{Double solids}, Adv.~Math. 47 (1983), 107--230.

\bibitem{Connes-1985} A. Connes, \emph{Non-commutative differential geometry}, Publ.~Math.~IHES 62 (1985), 257--360.

\bibitem{Costello-2011} K. Costello, \emph{Renormalization and effective field theory}, AMS, 2011.

\bibitem{Costello-Li-2016} K. Costello, S. Li, \emph{Twisted supergravity and its quantization}, arXiv:1606.00365, 2016. Proposition 5.2.

\bibitem{DVV-1997} R. Dijkgraaf, E. Verlinde, H. Verlinde, \emph{Counting dyons in $\cN = 4$ string theory}, Nucl.~Phys.~B 484 (1997), 543--561.

\bibitem{DGMS-1975} P. Deligne, P. Griffiths, J. Morgan, D. Sullivan, \emph{Real homotopy theory of K\"ahler manifolds}, Invent.~Math. 29 (1975), 245--274.

\bibitem{Eichler-Zagier-1985} M. Eichler, D. Zagier, \emph{The theory of Jacobi forms}, Progr.~Math. 55, Birkh\"auser, 1985.

\bibitem{Friedman-1986} R. Friedman, \emph{Simultaneous resolution of threefold double points}, Math.~Ann. 274 (1986), 671--689. Theorem 5.10.

\bibitem{Getzler-1994} E. Getzler, \emph{Batalin--Vilkovisky algebras and two-dimensional topological field theories}, Commun.~Math.~Phys. 159 (1994), 265--285.

\bibitem{Getzler-Jones-1994} E. Getzler, J.D.S. Jones, \emph{Operads, homotopy algebra and iterated integrals for double loop spaces}, preprint, hep-th/9403055, 1994. Theorem 2.1.

\bibitem{Getzler-Kapranov-1995} E. Getzler, M. Kapranov, \emph{Cyclic operads and cyclic homology}, in: \emph{Geometry, topology, and physics}, Conf.~Proc.~Lecture Notes Geom.~Topol.~4, 1995, 167--201.

\bibitem{Gilkey-1995} P. Gilkey, \emph{Invariance theory, the heat equation, and the Atiyah--Singer index theorem}, CRC Press, 2nd ed., 1995.

\bibitem{Griffiths-1969} P. Griffiths, \emph{On the periods of certain rational integrals I, II}, Ann.~Math. 90 (1969), 460--541.

\bibitem{Gritsenko-1999} V. Gritsenko, \emph{Modular forms and moduli spaces of abelian and K3 surfaces}, St.~Petersburg Math.~J. 6 (1995), 1179--1208.

\bibitem{Gritsenko-Nikulin-1998} V. Gritsenko, V. Nikulin, \emph{Automorphic forms and Lorentzian Kac--Moody algebras}, Int.~J.~Math. 9 (1998), 153--199.

\bibitem{Illusie-1971} L. Illusie, \emph{Complexe cotangent et d\'eformations I, II}, Lect.~Notes in Math. 239, 283, Springer, 1971--72. Chapter IV \S2.3.

\bibitem{Kapranov-1999} M. Kapranov, \emph{Rozansky--Witten invariants via Atiyah classes}, Compositio~Math. 115 (1999), 71--113. Proposition 4.4.

\bibitem{Keller-2003} B. Keller, \emph{Derived invariance of higher structures on the Hochschild complex}, preprint, 2003.

\bibitem{Kodaira-Morrow-1971} K. Kodaira, J. Morrow, \emph{Complex manifolds}, Holt, Rinehart, Winston, 1971.

\bibitem{Kontsevich-1999} M. Kontsevich, \emph{Operads and motives in deformation quantization}, Lett.~Math.~Phys. 48 (1999), 35--72.

\bibitem{Kontsevich-2003} M. Kontsevich, \emph{Deformation quantization of Poisson manifolds}, Lett.~Math.~Phys. 66 (2003), 157--216.

\bibitem{Kontsevich-Soibelman-2000} M. Kontsevich, Y. Soibelman, \emph{Deformations of algebras over operads and Deligne's conjecture}, in: \emph{Confér. Moshé Flato 1999}, Kluwer, 2000, 255--307. Theorem 2.4.

\bibitem{Kontsevich-Vlassopoulos-2022} M. Kontsevich, Y. Vlassopoulos, \emph{Weak Calabi--Yau structures and moduli of $A_\infty$-algebras}, 2022.

\bibitem{Kotake-Kato-1984} T. Kotake, S. Kato, \emph{Fundamental solutions of parabolic equations}, Osaka J.~Math. 21 (1984), 55--78.

\bibitem{Kuga-Satake-1967} M. Kuga, I. Satake, \emph{Abelian varieties attached to polarized $K3$ surfaces}, Math.~Ann. 169 (1967), 239--242.

\bibitem{Kuznetsov-2004} A. Kuznetsov, \emph{Hyperplane sections and derived categories}, Izv.~Math. 70 (2006), 447--547.

\bibitem{Loday-1998} J.-L. Loday, \emph{Cyclic homology}, Grundlehren Math.~Wiss. 301, Springer, 2nd ed., 1998. \S6.1.

\bibitem{Lurie-HA} J. Lurie, \emph{Higher Algebra}, available online, 2017. Definition 2.1.1.10, Theorem 4.1.8.4, \S4.8.1, Theorem 5.5.3.8.

\bibitem{Lurie-HTT} J. Lurie, \emph{Higher Topos Theory}, Ann.~of Math.~Stud. 170, Princeton UP, 2009. \S3.2.

\bibitem{Lunts-Orlov-2010} V. Lunts, D. Orlov, \emph{Uniqueness of enhancement for triangulated categories}, J.~AMS 23 (2010), 853--908. Theorem 2.8.

\bibitem{Markarian-2009} N. Markarian, \emph{The Atiyah class, Hochschild cohomology and the Riemann--Roch theorem}, J.~Lond.~Math.~Soc. 79 (2009), 129--143. Theorem 1.1.

\bibitem{MSS-2002} M. Markl, S. Shnider, J. Stasheff, \emph{Operads in algebra, topology and physics}, Math.~Surv.~Monogr. 96, AMS, 2002. \S5.4.

\bibitem{Michelsohn-1982} M.-L. Michelsohn, \emph{On the existence of special metrics in complex geometry}, Acta~Math. 149 (1982), 261--295. Theorem 6.8.

\bibitem{Morrison-1984} D. Morrison, \emph{The Kuga--Satake variety of an abelian surface}, J.~Algebra 92 (1985), 454--476.

\bibitem{Mukai-1987} S. Mukai, \emph{On the moduli space of bundles on $K3$ surfaces, I}, in: \emph{Vector bundles on algebraic varieties}, TIFR, 1987.

\bibitem{Nikulin-1979} V. Nikulin, \emph{Integer symmetric bilinear forms and some of their geometric applications}, Math.~USSR-Izv. 14 (1980), 103--167.

\bibitem{Positselski-2011} L. Positselski, \emph{Two kinds of derived categories, Koszul duality, and comodule-contramodule correspondence}, Mem.~AMS 212 (2011), no.~996. §6.5.

\bibitem{Quillen-1969} D. Quillen, \emph{Rational homotopy theory}, Ann.~Math. 90 (1969), 205--295. §9.

\bibitem{Ramadoss-2008} A. Ramadoss, \emph{The relative Riemann--Roch theorem from Hochschild homology}, New York J.~Math. 14 (2008), 643--717. Theorem 3.3.1.

\bibitem{Scheithauer-2008} N. Scheithauer, \emph{The Weil representation of $\mathrm{SL}_2(\Zbb)$ and some applications}, Int.~Math.~Res.~Not. 2008, rnn072.

\bibitem{Sen-2008} A. Sen, \emph{Black hole entropy function, attractors and precision counting of microstates}, Gen.~Relativ.~Gravit. 40 (2008), 2249--2431.

\bibitem{Strominger-1986} A. Strominger, \emph{Superstrings with torsion}, Nucl.~Phys.~B 274 (1986), 253--284.

\bibitem{Tamarkin-2007} D. Tamarkin, \emph{What do dg-categories form?}, Compositio Math. 143 (2007), 1335--1358.

\bibitem{Toen-2014} B. To\"en, \emph{Derived algebraic geometry}, EMS Surv.~Math.~Sci. 1 (2014), 153--240. Theorem 1.4.

\bibitem{Toen-Vezzosi-2015} B. To\"en, G. Vezzosi, \emph{Caract\`eres de Chern, traces \'equivariantes et g\'eom\'etrie alg\'ebrique d\'eriv\'ee}, Selecta Math. 21 (2015), 449--554. Theorem 2.6.1.

\bibitem{Voisin-1993} C. Voisin, \emph{Miroirs et involutions sur les surfaces $K3$}, in: \emph{Journ\'ees de G\'eom\'etrie Alg\'ebrique d'Orsay}, Ast\'erisque 218 (1993), 273--323.


\bibitem{Adler-Bardeen-1969} S. Adler, W.~A. Bardeen, \emph{Absence of higher-order corrections in the anomalous axial-vector divergence equation}, Phys.~Rev. 182 (1969), 1517--1536.

\bibitem{Banerjee-Gupta-Sen-2011} S. Banerjee, R.~K. Gupta, A. Sen, \emph{Logarithmic corrections to extremal black hole entropy from quantum entropy function}, JHEP 03 (2011), 147.

\bibitem{Bardeen-Zumino-1984} W.~A. Bardeen, B. Zumino, \emph{Consistent and covariant anomalies in gauge and gravitational theories}, Nucl.~Phys.~B 244 (1984), 421--453.

\bibitem{Bhattacharyya-Sen-2008} J. Bhattacharyya, A. Sen, \emph{Quarter BPS dyons in $\mathcal{N}=4$ supersymmetric type-II string theories}, JHEP 04 (2008), 012.

\bibitem{Fu-Yau-2008} J.-X. Fu, S.-T. Yau, \emph{The theory of superstring with flux on non-K\"ahler manifolds and the complex Monge--Amp\`ere equation}, J.~Diff.~Geom. 78 (2008), 369--428.

\bibitem{Green-Schwarz-1984} M.~B. Green, J.~H. Schwarz, \emph{Anomaly cancellation in supersymmetric D=10 gauge theory and superstring theory}, Phys.~Lett.~B 149 (1984), 117--122.

\bibitem{Gualtieri-2011} M. Gualtieri, \emph{Generalized complex geometry}, Ann.~Math. 174 (2011), 75--123. Theorem 3.8.

\bibitem{Hitchin-2010} N. Hitchin, \emph{Lectures on generalized geometry}, in: \emph{Surveys in differential geometry}, vol.~XVI, Int.~Press, 2010, 79--124.

\bibitem{Hull-Townsend-1995} C.~M. Hull, P.~K. Townsend, \emph{Unity of superstring dualities}, Nucl.~Phys.~B 438 (1995), 109--137.

\bibitem{Malikov-Schechtman-1999} F. Malikov, V. Schechtman, \emph{Chiral de Rham complex. II}, Amer.~Math.~Soc.~Transl. 194 (1999), 149--188.

\bibitem{MSV-1999} F. Malikov, V. Schechtman, A. Vaintrob, \emph{Chiral de Rham complex}, Commun.~Math.~Phys. 204 (1999), 439--473.

\bibitem{Maldacena-Moore-Strominger-1999} J. Maldacena, G. Moore, A. Strominger, \emph{Counting BPS black holes in toroidal type II string theory}, arXiv:hep-th/9903163, 1999. \S3.

\bibitem{Maloney-Witten-2007} A. Maloney, E. Witten, \emph{Quantum gravity partition functions in three dimensions}, JHEP 02 (2010), 029.

\bibitem{Narain-1986} K.~S. Narain, \emph{New heterotic string theories in uncompactified dimensions $<10$}, Phys.~Lett.~B 169 (1986), 41--46.

\bibitem{Polchinski-1998} J. Polchinski, \emph{String theory. Vol. I}, Cambridge Univ.~Press, 1998. \S4.6.

\bibitem{Shioda-Inose-1977} T. Shioda, H. Inose, \emph{On singular $K3$ surfaces}, in: \emph{Complex analysis and algebraic geometry}, Cambridge Univ.~Press, 1977, 119--136.

\bibitem{SYZ-1996} A. Strominger, S.-T. Yau, E. Zaslow, \emph{Mirror symmetry is $T$-duality}, Nucl.~Phys.~B 479 (1996), 243--259.

\bibitem{tHooft-1980} G. 't Hooft, \emph{Naturalness, chiral symmetry, and spontaneous chiral symmetry breaking}, in: \emph{Recent developments in gauge theories}, NATO Adv.~Study Inst., 1980, 135--157.

\bibitem{Weinberg-Vol-II} S. Weinberg, \emph{The quantum theory of fields, Vol.~II: Modern applications}, Cambridge Univ.~Press, 1996. \S22.5.

\bibitem{Wess-Zumino-1971} J. Wess, B. Zumino, \emph{Consequences of anomalous Ward identities}, Phys.~Lett.~B 37 (1971), 95--97.

\bibitem{Witten-1987} E. Witten, \emph{Elliptic genera and quantum field theory}, Commun.~Math.~Phys. 109 (1987), 525--536.

\bibitem{Witten-2007} E. Witten, \emph{Two-dimensional models with $(0, 2)$ supersymmetry: perturbative aspects}, Adv.~Theor.~Math.~Phys. 11 (2007), 1--63. \S3.

\bibitem{Zumino-1985} B. Zumino, \emph{Chiral anomalies and differential geometry}, in: \emph{Relativity, groups and topology II}, Les Houches 1983, North-Holland, 1985, 1291--1322.

\bibitem{Beilinson-1984} A. Beilinson, \emph{Higher regulators and values of $L$-functions}, J.~Soviet Math. 30 (1985), 2036--2070 (translated from Sovremennye Problemy Matematiki 24 (1984), 181--238).

\bibitem{Beilinson-1985} A. Beilinson, \emph{Notes on absolute Hodge cohomology}, in: \emph{Applications of algebraic $K$-theory to algebraic geometry and number theory}, Contemp.~Math. 55, AMS, 1986, 35--68.

\bibitem{Bloch-1974} S. Bloch, \emph{$K_2$ and algebraic cycles}, Ann.~Math. 99 (1974), 349--379.

\bibitem{Bloch-1986} S. Bloch, \emph{Algebraic cycles and higher $K$-theory}, Adv.~Math. 61 (1986), 267--304.

\bibitem{Esnault-Viehweg-1988} H. Esnault, E. Viehweg, \emph{Deligne--Be\u{\i}linson cohomology}, in: \emph{Be\u{\i}linson's conjectures on special values of $L$-functions}, Perspect.~Math. 4, Academic Press, 1988, 43--91.

\bibitem{Green-1989} M. Green, \emph{Griffiths' infinitesimal invariant and the Abel--Jacobi map}, J.~Differential Geom. 29 (1989), 545--555. Theorem 2.1.

\bibitem{Deligne-1971} P. Deligne, \emph{Th\'eorie de Hodge II}, Publ.~Math.~IHES 40 (1971), 5--57.

\bibitem{Voisin-2002} C. Voisin, \emph{Hodge theory and complex algebraic geometry I, II}, Cambridge Studies in Advanced Mathematics 76, 77, Cambridge Univ.\ Press, 2002/2003. Theorem 19.5.

\bibitem{Deligne-Goncharov-2005} P. Deligne, A. Goncharov, \emph{Groupes fondamentaux motiviques de Tate mixte}, Ann.~Sci.~\'Ecole Norm.~Sup. 38 (2005), 1--56. Proposition 5.27.

\bibitem{Kontsevich-1997} M. Kontsevich, \emph{Formality conjecture}, in: \emph{Deformation theory and symplectic geometry}, Math.~Phys.~Stud. 20, Kluwer, 1997, 139--156. Theorem 4.6.

\bibitem{Tamarkin-1998} D. Tamarkin, \emph{Another proof of M. Kontsevich's formality theorem}, arXiv:math/9803025, 1998. \S4.

\bibitem{Cattaneo-Felder-2001} A. Cattaneo, G. Felder, \emph{On the globalization of Kontsevich's star product and the perturbative Poisson sigma model}, Prog.~Theor.~Phys.~Suppl. 144 (2001), 38--53. Proposition 4.3.

\bibitem{Van-den-Bergh-2007} M. Van den Bergh, \emph{On global deformation quantization in the algebraic case}, J.~Algebra 315 (2007), 326--395.

\bibitem{Shoikhet-2003} B. Shoikhet, \emph{A proof of the Tsygan formality conjecture for chains}, Adv.~Math. 179 (2003), 7--37. Theorem 2.1.

\bibitem{Tsygan-1999} B. Tsygan, \emph{Formality conjectures for chains}, in: \emph{Differential topology, infinite-dimensional Lie algebras, and applications}, AMS Transl.~Ser.~2 194, AMS, 1999, 261--274.

\bibitem{Markl-Shnider-Stasheff-2002} M. Markl, S. Shnider, J. Stasheff, \emph{Operads in algebra, topology and physics}, Math.~Surveys Monogr. 96, AMS, 2002. Lemma 2.9.

\bibitem{Bezrukavnikov-1996} R. Bezrukavnikov, \emph{Koszul DG-algebras arising from configuration spaces}, Geom.~Funct.~Anal. 4 (1994), 119--135. Theorem 1.

\bibitem{Kriz-1994} I. Kriz, \emph{On the rational homotopy type of configuration spaces}, Ann.~Math. 139 (1994), 227--237. Theorem III.3.

\bibitem{Hain-1987} R. Hain, \emph{The de Rham homotopy theory of complex algebraic varieties I}, $K$-Theory 1 (1987), 271--324. \S6.

\bibitem{Shklyarov-2013} D. Shklyarov, \emph{Hirzebruch--Riemann--Roch-type formula for DG algebras}, Proc.~London Math.~Soc. 106 (2013), 1--32. Theorem 1.3.

\bibitem{Brav-Dyckerhoff-2019} C. Brav, T. Dyckerhoff, \emph{Relative Calabi--Yau structures}, Compositio Math. 155 (2019), 372--412. \S3, Corollary 3.18.

\bibitem{McClure-Smith-2002} J. McClure, J. Smith, \emph{A solution of Deligne's Hochschild cohomology conjecture}, in: \emph{Recent progress in homotopy theory}, Contemp.~Math. 293, AMS, 2002, 153--193.

\bibitem{Kaledin-2018} D. Kaledin, \emph{Spectral sequences for cyclic homology}, in: \emph{Algebra, geometry, and physics in the 21st century} (Kontsevich Festschrift), Progr.~Math. 324, Birkh\"auser, 2017, 99--129.

\bibitem{Gaitsgory-Rozenblyum-2017} D. Gaitsgory, N. Rozenblyum, \emph{A study in derived algebraic geometry, Vol. I--II}, Math. Surv. Monogr. 221, AMS, 2017.

\bibitem{Lurie-DAG-X} J. Lurie, \emph{DAG X: Formal moduli problems}, preprint, 2011. Theorem 5.2.3.11.

\bibitem{Hartshorne-1977} R. Hartshorne, \emph{Algebraic geometry}, Grad.~Texts~Math. 52, Springer, 1977. Theorem III.2.7.

\bibitem{Kapranov-Soibelman-2009} M. Kapranov, Y. Soibelman, \emph{Stability conditions on $A_\infty$-categories}, arXiv:math/0606233, 2009. \S4.7.

\bibitem{Ayala-Francis-Tanaka-2017} D. Ayala, J. Francis, H.L. Tanaka, \emph{Factorization homology of stratified spaces}, Selecta Math. 23 (2017), 293--362. Theorem 4.3.1.

\bibitem{Atiyah-1957} M.F. Atiyah, \emph{Complex analytic connections in fibre bundles}, Trans.~AMS 85 (1957), 181--207. \S4.

\end{thebibliography}

\end{document}
